% ===========================================================================
%  LiteraLevel — Documento final del Módulo 5 (Integrador)
%  Diplomado en Ciencia de Datos, FES Acatlán (UNAM), Generación 34
%  Autor: Miguel Reyes Carrillo · Profesor: Mtro. Fernando Barranco Rodríguez
%
%  Compila con pdflatex, lualatex, xelatex o tectonic (+ biber para APA).
%  Todas las figuras están dibujadas con TikZ/pgfplots: no hay imágenes
%  externas ni capturas de pantalla en este documento.
%
%    tectonic documento.tex          (necesita biber en el PATH)
%    pdflatex documento && biber documento && pdflatex documento (x2)
% ===========================================================================

\documentclass[11pt,a4paper]{article}

% --- Motor y codificación -------------------------------------------------
\usepackage{iftex}
\usepackage{graphicx}
\ifPDFTeX
  \usepackage[T1]{fontenc}
  \usepackage[utf8]{inputenc}
  \usepackage{lmodern}
\else
  \usepackage{fontspec}
\fi

\usepackage[spanish,es-noquoting,es-nodecimaldot]{babel}
\usepackage{csquotes}

% --- Página ---------------------------------------------------------------
\usepackage[a4paper,margin=2.5cm,headheight=14pt]{geometry}
\usepackage{microtype}
\setlength{\emergencystretch}{1em}
\usepackage{setspace}
\usepackage{fancyhdr}
\usepackage{titlesec}
\usepackage{enumitem}

% --- Matemáticas ----------------------------------------------------------
\usepackage{amsmath,amssymb}

% --- Cuadros --------------------------------------------------------------
\usepackage{booktabs}
\usepackage{array}
\usepackage{tabularx}
\usepackage{longtable}
\usepackage{multirow}
\usepackage[font=small,labelfont=bf,justification=justified,singlelinecheck=false]{caption}

% --- Figuras: TODO se dibuja aquí -----------------------------------------
\usepackage{xcolor}
\usepackage{pgfplots}
\usepackage{tikz}
\usetikzlibrary{arrows.meta,positioning,shapes.geometric,shapes.misc,calc,fit,backgrounds,decorations.pathreplacing}
\pgfplotsset{compat=1.18}
\usepgfplotslibrary{groupplots}
\usepgfplotslibrary{fillbetween}
\pgfplotsset{
  /pgf/number format/.cd, use comma, 1000 sep={\,},
}

% --- Referencias APA ------------------------------------------------------
\usepackage[style=apa,backend=biber,natbib=true]{biblatex}
\DeclareLanguageMapping{spanish}{spanish-apa}
\addbibresource{referencias.bib}

% --- Enlaces --------------------------------------------------------------
\usepackage[colorlinks=true,linkcolor=llazul,citecolor=llazul,urlcolor=llazul]{hyperref}

% ===========================================================================
%  Paleta y estilos propios
% ===========================================================================
\definecolor{llazul}{HTML}{1F4E79}
\definecolor{llcian}{HTML}{2E86AB}
\definecolor{llverde}{HTML}{2A9D8F}
\definecolor{llambar}{HTML}{E9A13B}
\definecolor{llrojo}{HTML}{C1495B}
\definecolor{llgris}{HTML}{5A5A5A}
\definecolor{llgrisclaro}{HTML}{E8ECF0}
\definecolor{llfondo}{HTML}{F7F9FB}

\newcolumntype{L}[1]{>{\raggedright\arraybackslash}p{#1}}
\newcolumntype{C}[1]{>{\centering\arraybackslash}p{#1}}
\newcolumntype{R}[1]{>{\raggedleft\arraybackslash}p{#1}}
\newcolumntype{Y}{>{\raggedright\arraybackslash}X}

\newcommand{\cod}[1]{\texttt{\small #1}}
% babel-spanish redefine \% con un ajuste de espaciado que choca dentro de
% tabularx/longtable ("Incompatible glue units"). Se recupera el signo llano.
\AtBeginDocument{\def\%{\char`\%\relax}}
\newcommand{\pct}{\,\%}

\titleformat{\section}{\normalfont\Large\bfseries\color{llazul}}{\thesection}{0.7em}{}
\titleformat{\subsection}{\normalfont\large\bfseries\color{llazul!85!black}}{\thesubsection}{0.6em}{}
\titleformat{\subsubsection}{\normalfont\normalsize\bfseries\color{llgris}}{\thesubsubsection}{0.6em}{}

\pagestyle{fancy}
\fancyhf{}
\fancyhead[L]{\footnotesize\color{llgris}LiteraLevel — Módulo 5 (Integrador)}
\fancyhead[R]{\footnotesize\color{llgris}\thepage}
\renewcommand{\headrulewidth}{0.4pt}
\renewcommand{\footrulewidth}{0pt}

% Estilo común de los ejes de pgfplots -------------------------------------
\pgfplotsset{
  ejell/.style={
    width=\linewidth,
    axis line style={llgris!60},
    tick style={llgris!60},
    label style={font=\small},
    tick label style={font=\footnotesize},
    title style={font=\small\bfseries, color=llazul},
    legend style={font=\footnotesize, draw=llgris!40, fill=white, fill opacity=0.9,
                  text opacity=1, rounded corners=2pt},
    grid style={llgris!20, line width=0.3pt},
    every axis plot/.append style={line width=0.9pt},
  },
  barrall/.style={
    ejell,
    xmajorgrids=true,
    nodes near coords style={font=\scriptsize, color=llgris},
  },
}

\tikzset{
  cajall/.style={
    draw=llazul!55, fill=llfondo, rounded corners=3pt, line width=0.7pt,
    align=center, font=\footnotesize, inner sep=4pt, minimum height=10mm,
    text width=23mm
  },
  cajamodelo/.style={cajall, fill=llcian!14, draw=llcian!75, font=\footnotesize\bfseries},
  cajasalida/.style={cajall, fill=llverde!14, draw=llverde!75},
  cajarechazo/.style={cajall, fill=llrojo!10, draw=llrojo!65, text width=22mm},
  flechall/.style={-{Stealth[length=2.6mm]}, draw=llgris, line width=0.8pt},
  etiqll/.style={font=\scriptsize\itshape, color=llgris, inner sep=2pt},
  notall/.style={font=\scriptsize\itshape, color=llgris, inner sep=1pt,
                 align=center, text width=30mm},
}

% ===========================================================================
\begin{document}
% ===========================================================================

% ---------------------------------------------------------------- Portada --
\begin{titlepage}
\centering

% --- Escudos institucionales -------------------------------------------
% Misma maquetacion que el proyecto del modulo 4, y con sus mismos archivos.
% Son las UNICAS imagenes externas del documento: un escudo es una marca
% institucional, no una figura de datos. Las doce figuras siguen en TikZ.
\begin{minipage}{0.25\textwidth}
    \raggedright
    \includegraphics[height=2.7cm]{logos/logo_unam.png}
\end{minipage}
\hfill
\begin{minipage}{0.45\textwidth}
    \centering
    {\Large\bfseries Universidad Nacional Autónoma de México\par}

    \vspace{0.2cm}

    {\large FACULTAD DE ESTUDIOS SUPERIORES ACATLÁN\par}
\end{minipage}
\hfill
\begin{minipage}{0.25\textwidth}
    \raggedleft
    \includegraphics[height=2.7cm]{logos/logo_fes_acatlan.png}
\end{minipage}

\vspace{2.0cm}

{\Huge\bfseries\color{llazul} LiteraLevel\par}

\vspace{0.6cm}

{\LARGE Medición automática de la dificultad de lectura\\[0.25cm]
en inglés y su entrega a un lector real\par}

\vspace{1.8cm}

{\Large\bfseries Proyecto integrador --- Módulo 5\par}

\vspace{1.2cm}

{\large
\textbf{Alumno:}\\
Miguel Reyes Carrillo
\par}

\vspace{0.7cm}

{\large
\textbf{Profesor:}\\
Mtro. Fernando Barranco Rodríguez
\par}

\vspace{0.7cm}

{\large
Diplomado en Ciencia de Datos --- Generación 34
\par}

\vfill

% El sello de cifras al pie: apoya, no encabeza.
\begin{tikzpicture}[font=\small]
  \node[draw=llazul!35, fill=llfondo, rounded corners=6pt, line width=0.8pt,
        inner sep=9pt, text width=12.4cm, align=center] (caja) {
    \begin{minipage}{12.1cm}
      \centering
      \begin{tabular}{@{}C{3.6cm}C{4.0cm}C{3.3cm}@{}}
        {\large\bfseries\color{llcian} 0,908} & {\large\bfseries\color{llverde} A1 $<\cdots<$ C1} & {\large\bfseries\color{llambar} 97,2\pct} \\[2pt]
        {\scriptsize Spearman, validación interna} & {\scriptsize orden MCER creciente en corpus externo} & {\scriptsize dentro de $\pm 1$ nivel} \\
      \end{tabular}
    \end{minipage}};
\end{tikzpicture}

\vspace{0.6cm}

{\normalsize Ciudad de México, septiembre de 2026\par}

\end{titlepage}

% ---------------------------------------------------------------- Resumen --
\section*{Resumen}
\addcontentsline{toc}{section}{Resumen}

Aprender inglés leyendo funciona, pero solo si el libro está al nivel de quien lo
abre. Un lector que elige un libro demasiado difícil no aprende: lo abandona. Y hoy
no hay forma sencilla de saber, antes de empezar, si \emph{ese} libro concreto está a
su alcance.

LiteraLevel es una aplicación que responde esa pregunta. Le da al lector un número
entre 0 y 2 que dice cuánta dificultad tiene un texto en inglés, y lo hace leyendo el
texto de verdad, no contando páginas. Con ese número le busca libros que le queden a
mano, se los abre dentro de la propia aplicación, le enseña la dificultad capítulo
a capítulo —los libros no son igual de difíciles de principio a fin—, le explica
cualquier palabra en la frase en la que aparece y le corrige la pronunciación cuando
lee en voz alta.

Debajo hay un modelo de lenguaje pequeño (\cod{distilroberta-base}) afinado para
predecir dificultad como un número continuo. En la validación interna —144 documentos
reservados del corpus de anclaje, los mismos con los que se eligió el punto de control—
alcanza una correlación de Spearman de 0,908 con el nivel de referencia y, sobre un corpus externo
que no se usó para ajustar los pesos —aunque sí se exploró durante el desarrollo, y eso
se detalla en la limitación 6—, ordena los niveles del Marco Común Europeo de forma
estrictamente
creciente: A1 0,5040 ($n{=}3$) · A2 0,5561 ($n{=}8$) · B1 0,9276 ($n{=}12$) ·
B2 1,2414 ($n{=}12$) · C1 1,4952 ($n{=}12$). Ese orden —no el acierto exacto— es el
resultado que sostiene el trabajo, \textbf{con una salvedad que conviene leer junto a las
cifras}: el escalón A1--A2 es de 0,052 sobre desviaciones de 0,111 y 0,223, y descansa en
tres textos. Está dentro del ruido. Los tres escalones siguientes, no.

Conviene decirlo desde la primera página: \textbf{la escala 0--2 es propia y
LiteraLevel no certifica ningún nivel del Marco Común Europeo}. Las letras A1 a C2
fueron el ancla con la que se entrenó el modelo, no un título que la aplicación
expida. Sirven para elegir qué leer, no para acreditar nada.

La entrega usa \textbf{Angular como interfaz y una API FastAPI sobre el motor Python};
Gradio 6.26 \parencite{abid2019gradio} se conserva como respaldo. Las pruebas de usabilidad
de la interfaz anterior con Esmeralda y Laura motivaron esta separación, pero no equivalen
a una validación de usuario de todos los cambios posteriores. El catálogo medido
contiene \textbf{53.179 libros} en inglés de dominio
público de Project Gutenberg, leídos enteros por el modelo, capítulo a capítulo. La última
tanda, en una GPU gratuita, tardó 55 minutos: el cuello de botella nunca fue el modelo,
sino leer y guardar los archivos.

\vspace{0.4cm}
\noindent\textbf{Palabras clave:} legibilidad automática, MCER, adquisición de
segundas lenguas, lectura extensiva, transformers, Project Gutenberg, Angular.

\vspace{0.8cm}
\noindent\fbox{\begin{minipage}{0.97\textwidth}
\small\textbf{Nota sobre la estructura de la entrega.} La \emph{Estrategia}, que la
rúbrica pide como entregable propio y como sección exigida del documento, se desarrolla
completa en el \textbf{apartado \ref{sec:estrategia}} y no en un anexo: ocupa el lugar
que le corresponde en el orden de secciones y sostiene por sí sola el resto del texto.
Las cifras que allí se citan y que se miden más adelante van referidas con su cuadro o
figura. \textbf{Todas las figuras de este documento están dibujadas en \LaTeX{}}
(TikZ/pgfplots) a partir de los datos medidos que acompañan a la entrega; no hay
ninguna captura de pantalla ni imagen externa.
\end{minipage}}

\clearpage
\tableofcontents
\clearpage
\listoffigures
\clearpage

% ===========================================================================
\section{Introducción y motivación}
% ===========================================================================

Unos 1.500 millones de personas estudian inglés como segunda lengua. De todos los
métodos con respaldo empírico, la \textbf{lectura extensiva} —leer mucho, por gusto, en
textos ligeramente por encima del nivel actual— es uno de los que mejor se sostiene: el
metaanálisis de \textcite{nakanishi2015meta} sobre décadas de estudios encuentra
efectos consistentes en comprensión y vocabulario, y \textcite{day1998extensive} la
formalizan como programa de aula. \textcite{krashen1985input} llamó a esa franja
«\emph{input} comprensible»: material que se entiende casi todo, con lo justo de
novedad para aprender.

El problema es operativo, no teórico. Para que la lectura extensiva funcione, el texto
tiene que caer en esa franja, y \textcite{hu2000unknown} la cuantifican: se necesita
reconocer alrededor del 98\pct{} de las palabras corridas de un texto para comprenderlo
sin ayuda —umbral coherente con el trabajo posterior de \textcite{nation2013learning}
sobre tamaño de vocabulario y cobertura léxica—. Un punto porcentual por debajo y la
lectura se convierte en trabajo de diccionario. Es decir: \textbf{el mismo libro es
excelente para un lector y una pérdida de tiempo para otro que está dos niveles por
debajo}, y nada en la portada lo dice.

Hoy el lector hispanohablante que quiere leer en inglés se enfrenta a decenas de miles
de libros gratuitos en dominio público y a ninguna señal fiable sobre cuál le
corresponde. Las etiquetas que circulan (A2, B1, C1) casi nunca las puso alguien que
haya leído la obra completa, y además tratan el libro como si tuviera una sola
dificultad, cuando en realidad cambia de capítulo a capítulo. La
figura~\ref{fig:curva_wp} lo muestra sobre un libro entero medido: la dificultad de
\emph{War and Peace} recorre casi setenta centésimas de la escala pública entre su
capítulo más llevadero y el más exigente.

\vspace{0.2cm}
\noindent\textbf{Pregunta que guía el trabajo:} ¿es posible estimar automáticamente,
con un modelo de lenguaje entrenado sobre corpus públicos, una medida de dificultad de
un texto en inglés que (a) ordene correctamente los niveles del MCER
\parencite{coe2020cefr} en material externo que no intervino en el ajuste de los pesos, y (b) sea lo
bastante rápida y comprensible para ponerse en manos de un lector adolescente sin
entrenamiento previo?

Este documento responde a las dos mitades. La primera es un problema de modelado y se
resuelve en los apartados \ref{sec:datos} a \ref{sec:modelacion}. La segunda es un
problema de producto y de usabilidad, que es lo que el módulo declara evaluar, y se
desarrolla en los apartados \ref{sec:estrategia} y \ref{sec:resolucion}.

% ===========================================================================
\section{Planteamiento del problema}
% ===========================================================================

\subsection{Formulación}
\label{sec:formulacion}

Sea $T$ un texto en inglés en prosa continua. Se busca una función
$f(T) \rightarrow [0, 2]$ tal que, para dos textos $T_i$ y $T_j$ cuyos niveles
oficiales de referencia cumplen $\text{nivel}(T_i) < \text{nivel}(T_j)$, se cumpla
$f(T_i) < f(T_j)$ con la mayor frecuencia posible.

Obsérvese que \textbf{el objetivo es ordinal, no de clasificación exacta}. La frontera
entre B1 y B2 es una convención pedagógica, no un hecho físico: confundir B1 con B2 es
un error muy distinto de confundir B1 con C2. Por eso las métricas principales del
trabajo son la correlación de Spearman y el \textbf{error ordinal medio} (distancia
media en niveles entre lo predicho y lo real), y no la exactitud de clase.

Un texto largo no se puntúa de una vez. Se parte en bloques de 250 palabras (mínimo
150) y se puntúa cada bloque; la salida del documento es la media sin ponderar de sus
bloques. Formalmente, si $B(T) = \{b_1, \dots, b_n\}$ son los bloques de $T$ y $g$ es
el modelo por bloque,

\begin{equation}
f(T) \;=\; 0{,}4 \cdot \frac{1}{n}\sum_{k=1}^{n} \operatorname{clip}\!\left(g(b_k),\; 0,\; 5\right),
\label{eq:f}
\end{equation}

\noindent donde el factor 0,4 convierte la escala interna del modelo (0--5, un valor
por cada nivel A1\dots C2) en la escala pública 0--2 que ve el usuario.

\textbf{El recorte va dentro del sumatorio, y el orden importa.} Cada bloque se lleva a
0--5 \emph{antes} de promediar, no después. No es un detalle de notación: la cabeza de
regresión no tiene ninguna activación que la obligue a quedarse en el intervalo, y el
\textbf{13,96\pct} de las ventanas de prosa literaria puntúa por encima de 5 —medido
sobre 530 ventanas de \emph{Alicia}, \emph{Orgullo y prejuicio} y \emph{Crimen y
castigo}; ninguna baja de 0—. Recortar después de promediar dejaría que un bloque que
puntúa 5,3 arrastrase hacia arriba a todo el capítulo. Las dos versiones de la fórmula
llegan a separarse \textbf{0,021} en un libro real, cinco veces la diferencia de
cualquier otra decisión de agregación del sistema. El recorte va dentro del sumatorio,
que es lo que hace el código. La
figura~\ref{fig:pipeline} recorre esa cadena completa, del texto crudo a la banda que
se muestra en pantalla, y la figura~\ref{fig:escala} sitúa las anclas MCER sobre la
recta 0--2 resultante.

% --------------------------------------------------------- FIGURA: pipeline
\begin{figure}[htbp]
\centering
\begin{tikzpicture}[node distance=7mm and 6mm]

  \node[cajall] (texto) {Texto crudo\\[1pt]{\scriptsize (libro, capítulo o pegado por el usuario)}};
  \node[cajall, right=of texto] (filtro) {Filtro de\\prosa continua\\[1pt]{\scriptsize sobre texto \emph{sin} normalizar}};
  \node[cajall, right=of filtro] (norm) {Normalización\\de espacios\\[1pt]{\scriptsize \cod{\textbackslash s+} $\rightarrow$ espacio}};
  \node[cajall, right=of norm] (vent) {Ventanas de\\250 palabras\\[1pt]{\scriptsize mínimo 150}};

  \node[cajarechazo, below=9mm of filtro] (rech) {No es prosa:\\no se puntúa\\[1pt]{\scriptsize menú, horario, tabla}};

  \node[cajamodelo, below=32mm of vent] (tok) {Tokenización\\y truncado\\a 256 \emph{tokens}};
  \node[cajamodelo, left=of tok, text width=30mm] (modelo) {\texttt{\scriptsize distilroberta-base}\\6 capas · 82 M\\cabeza de regresión};
  \node[cajall, left=of modelo] (media) {Media sin\\ponderar\\[1pt]{\scriptsize escala interna 0--5}};
  \node[cajasalida, left=of media] (escala) {$\times\,0{,}4$\\Escala pública\\\textbf{0--2}};

  \node[cajasalida, below=9mm of media] (banda) {Banda\\{\scriptsize básico / intermedio / avanzado}};
  \node[cajasalida, below=9mm of escala] (letra) {Letra MCER\\{\scriptsize solo orientativa}};

  \draw[flechall] (texto) -- (filtro);
  \draw[flechall] (filtro) -- node[etiqll,above]{sí} (norm);
  \draw[flechall] (norm) -- (vent);
  \draw[flechall] (filtro) -- node[etiqll,right]{no} (rech);
  \draw[flechall] (vent) -- (tok);
  \draw[flechall] (tok) -- (modelo);
  \draw[flechall] (modelo) -- (media);
  \draw[flechall] (media) -- (escala);
  \draw[flechall] (escala) -- (letra);
  \draw[flechall] (media) -- (banda);

  \node[notall, above=1mm of tok] {En prosa inglesa real,\\ el 100\pct{} de las
    ventanas llega a 256};
  \node[notall, below=1mm of modelo] {8,2 a 2.022 ventanas/s según el equipo};

\end{tikzpicture}
\caption[Cadena de medición de LiteraLevel]{Cadena de medición, de la ecuación~\eqref{eq:f}
al número que ve el lector. El filtro de prosa se aplica sobre el texto crudo, antes de
normalizar espacios: hacerlo al revés lo desactivaba en producción (cuadro~\ref{cuad:textos},
T1). En prosa inglesa real, las ventanas de un texto de más de 250 palabras ocupan
siempre 256 \emph{tokens} tras truncar, así que no llevan relleno y reagruparlas no movió ningún decimal en la
comprobación de 300 ventanas del apartado~\ref{sec:escalado}. Cuando se mezclan capítulos
cortos en un mismo lote sí hay relleno, y la diferencia medida es de
$4{,}8\times 10^{-7}$: por debajo de los tres decimales que se guardan.}
\label{fig:pipeline}
\end{figure}

% ----------------------------------------------------------- FIGURA: escala
\begin{figure}[htbp]
\centering
\begin{tikzpicture}[x=6.6cm, y=1cm]

  % Bandas de la interfaz (fondo)
  \fill[llverde!12]  (0,-0.32) rectangle (0.6,0.32);
  \fill[llambar!14]  (0.6,-0.32) rectangle (1.4,0.32);
  \fill[llrojo!12]   (1.4,-0.32) rectangle (2.0,0.32);

  % Recta
  \draw[line width=1.2pt, llgris] (0,0) -- (2.0,0);

  % Cortes MCER (datos_figuras/escala_cefr.dat)
  \foreach \x/\lab in {0.20/{0,20}, 0.60/{0,60}, 1.00/{1,00},
                      1.40/{1,40}, 1.80/{1,80}}{
    \draw[llgris!70, line width=0.6pt] (\x,-0.20) -- (\x,0.20);
    \node[font=\scriptsize, color=llgris, below=1pt] at (\x,-0.20) {\lab};
  }
  \node[font=\scriptsize, color=llgris, below=1pt] at (0,-0.20) {0,00};
  \node[font=\scriptsize, color=llgris, below=1pt] at (2.0,-0.20) {2,00};

  % Anclas MCER (centro de cada tramo = nivel entero x 0,4)
  \foreach \x/\lab in {0.0/A1, 0.4/A2, 0.8/B1, 1.2/B2, 1.6/C1, 2.0/C2}{
    \fill[llazul] (\x,0) circle (2.2pt);
    \node[font=\small\bfseries, color=llazul, above=4pt] at (\x,0) {\lab};
  }

  % Bandas de producto
  \node[font=\footnotesize\bfseries, color=llverde!60!black] at (0.3,-0.85) {básico};
  \node[font=\footnotesize\bfseries, color=llambar!60!black] at (1.0,-0.85) {intermedio};
  \node[font=\footnotesize\bfseries, color=llrojo!70!black]  at (1.7,-0.85) {avanzado};

  \draw[decorate, decoration={brace, amplitude=4pt, mirror}, llgris!70] (0,-0.55) -- (0.6,-0.55);
  \draw[decorate, decoration={brace, amplitude=4pt, mirror}, llgris!70] (0.6,-0.55) -- (1.4,-0.55);
  \draw[decorate, decoration={brace, amplitude=4pt, mirror}, llgris!70] (1.4,-0.55) -- (2.0,-0.55);

  \node[font=\footnotesize, color=llgris, above=16pt] at (1.0,0) {escala pública $f(T)$};

\end{tikzpicture}
\caption[La escala 0--2 y sus anclas MCER]{La escala pública 0--2 y sus anclas. Los
puntos azules son los seis niveles MCER tal como los coloca la ecuación~\eqref{eq:f}
(nivel entero $\times\,0{,}4$); las marcas grises son los cortes que decide qué letra se
muestra, y las llaves inferiores, las tres bandas con las que la interfaz comunica la
incertidumbre en lugar de fingir dos decimales de precisión. \emph{La letra es
orientación, no certificación.} Cortes tomados de \cod{datos\_figuras/\allowbreak escala\_cefr.dat}.}
\label{fig:escala}
\end{figure}

\subsection{Lo que el problema no es}

Tres delimitaciones que evitan malentendidos y que la aplicación repite en pantalla:

\begin{enumerate}[leftmargin=*, itemsep=3pt]
  \item \textbf{No es certificación.} El modelo se ancló a etiquetas MCER durante el
        entrenamiento; eso no lo convierte en un examinador acreditado. Devuelve una
        posición en una recta continua, y la letra que muestra es simplemente el ancla
        más cercana.
  \item \textbf{No es un medidor universal de texto.} El modelo solo vio prosa continua.
        Una carta de menú, un horario de vuelos o una tabla no son prosa, y puntuarlos
        daría un número sin significado. La aplicación los rechaza con un filtro
        estructural explícito (apartado~\ref{sec:proc_texto}).
  \item \textbf{No mide el interés ni la calidad literaria.} Mide dificultad de lectura.
        Un libro fácil puede ser magnífico y uno difícil, aburrido.
\end{enumerate}

\subsection{Criterio de éxito}

El proyecto se considera resuelto si se cumplen tres condiciones simultáneas, recogidas
en el cuadro~\ref{cuad:exito}.

\begin{table}[htbp]
\centering
\small
\begin{tabularx}{\textwidth}{@{}c L{6.2cm} Y@{}}
\toprule
\textbf{\#} & \textbf{Condición} & \textbf{Cómo se verifica} \\
\midrule
C1 & El orden por nivel MCER se conserva en corpus externo no visto & Medias por nivel estrictamente crecientes \\
C2 & La señal supera a la trivial & Correlación parcial descontando la longitud del texto \\
C3 & Un usuario sin formación técnica completa el recorrido & Prueba con lectora ajena al proyecto \\
\bottomrule
\end{tabularx}
\caption{Criterios de éxito del proyecto. C1 y C2 se cumplen y se documentan en los
apartados~\ref{sec:validacion} y~\ref{sec:baseline}. C3 se cumplió con la interfaz anterior, de
Gradio (\S\ref{sec:esmeralda}); sobre el \emph{front} de Angular que se entrega está pendiente
(accionable A8).}
\label{cuad:exito}
\end{table}

% ===========================================================================
\section{Estrategia}
\label{sec:estrategia}
% ===========================================================================

\subsection{El marco: Peldaño, una compañía ficticia}
\label{sec:compania}

Este trabajo se plantea como el primer producto de \textbf{Peldaño}, una compañía
ficticia de inteligencia artificial aplicada a la educación. La ficción es un requisito
del módulo y conviene decir desde el principio \textbf{qué es ficticio y qué no}: la
empresa lo es; el modelo, la aplicación desplegada, el corpus, las métricas y cada cifra
de este documento, no. Todo lo que se afirma aquí está medido, y qué parte se puede
además \emph{reproducir} y cuál no está dicho sin adornos en el
apartado~\ref{sec:validacion} y en las limitaciones: con lo público —los pesos del v4 en
Hugging Face, CLEAR, el catálogo de Project Gutenberg y el manifiesto de hashes de British
Council— se rehacen la inferencia y la cadena entera del cuaderno; los guiones de
verificación (\cod{validar\_local.py}, \cod{tools/\allowbreak bajar\_corpus\_bc.py},
\cod{entrenamiento\_v4/}) no viajan con los entregables y se comparten a petición; y el
corpus de British Council no se redistribuye, se reconstruye con su manifiesto.

El nombre no es decorativo. En adquisición de segundas lenguas, la hipótesis del
\emph{input} comprensible sostiene que se aprende con material que está justo un
escalón por encima de lo que uno ya domina \parencite{krashen1985input}. Ese
escalón es el peldaño, y encontrarlo es exactamente el problema que el lector no puede
resolver solo.

\subsubsection*{Nicho}

Peldaño no compite en el mercado general de aprendizaje de idiomas, que está saturado y
lo dominan aplicaciones de repetición espaciada. Se sitúa en un hueco más estrecho y
peor atendido: \textbf{la lectura extensiva en inglés para hispanohablantes que ya
tienen base}, entre A2 y B2. Es el tramo donde el estudiante ya no cabe en el material
graduado, todavía no aguanta una novela sin ayuda, y abandona. El apartado
\ref{sec:competencia} documenta contra quién se compite y dónde se pierde.

\subsubsection*{Propuesta de valor, en una frase}

Convertir una biblioteca gratuita e inabarcable —Project Gutenberg— en una estantería
personal ordenada por lo que cada lector puede leer hoy.

\subsubsection*{Cómo se sostendría}

La pregunta de negocio no es cómo cobrar por el modelo, que es pequeño y replicable,
sino quién paga la parte cara. Medido en este proyecto: \textbf{medir un libro, leerlo,
buscar y consultar el diccionario no cuestan nada} —corren en local sobre datos
abiertos—, y solo tres acciones salen a un servicio externo de pago: explicar una
palabra en su frase, resumir un libro y corregir la pronunciación. Esa separación, que
en la aplicación es una puerta de acceso con cupo por perfil, es también la frontera
natural entre un nivel gratuito y uno de pago, y no hubo que diseñarla: salió de medir
qué costaba cada cosa.

El cliente institucional —el docente de inglés del
apartado~\ref{sec:usuario-docente}— es el que justificaría una tarifa, porque su
necesidad no es leer sino \emph{asignar}: repartir lecturas de dificultad medida, y
medida capítulo a capítulo, sin tener que leerse cuarenta libros antes. Conviene decir
qué es esto y qué no: hoy el producto \textbf{no hace seguimiento de alumnos}, ni puede
hacerlo tal como está construido —los datos de cada lector se guardan separados y no se
comparten—, y la Declaración de alcance del final de este documento descarta
expresamente usar esta medida para decisiones con consecuencias para el evaluado. El
producto que hay sirve al docente para \emph{elegir textos}; un panel de seguimiento
sería otro producto, y está en Trabajo futuro.

\subsubsection*{Lo que Peldaño no promete}

No promete enseñar inglés, ni certificar un nivel, ni sustituir a un profesor. Mide la
dificultad de un texto en una escala propia de 0 a 2 y la pone al lado del nivel
estimado del lector. El apartado~\ref{sec:limitaciones} enumera los límites de esa
medida, y la aplicación los dice en pantalla en vez de esconderlos.

\subsection{La apuesta, en un párrafo}

La lectura extensiva es uno de los métodos con más respaldo para aprender una segunda
lengua, y sin embargo casi nadie la practica bien, porque elegir el libro adecuado es
un problema que el lector no puede resolver solo. LiteraLevel convierte una medida
continua de dificultad —producida por un modelo que lee el texto entero, no la
portada— en tres decisiones cotidianas: \textbf{qué leer ahora}, \textbf{por dónde se
va a complicar} y \textbf{qué significa esta palabra en esta frase}. El activo
defendible no es el modelo, que es pequeño y replicable: es la cadena completa desde el
corpus hasta el pulgar de una persona de quince años en un teléfono.

\subsection{Descripción de usuarios}

Tres perfiles, con lo que cada uno necesita y con lo que hoy hace en su lugar. No son
personas inventadas para rellenar: el primero está observado directamente
(apartado~\ref{sec:esmeralda}) y los otros dos se derivan de las peticiones reales
recogidas durante el desarrollo.

\subsubsection{Usuaria primaria: la estudiante de preparatoria}

\textbf{Perfil.} 14 a 18 años. Estudia inglés en la escuela, entiende bastante más de
lo que cree, no sabe qué es el MCER y no le interesa. Entra desde el teléfono, casi
siempre en vertical, y en pantallas de unos 390~px de ancho. Su tolerancia a la
fricción es de segundos. Sus necesidades y la respuesta del producto están en el
cuadro~\ref{cuad:primaria}.

\begin{table}[htbp]
\centering
\small
\begin{tabularx}{\textwidth}{@{}L{5.0cm} Y@{}}
\toprule
\textbf{Necesidad} & \textbf{Cómo se atiende} \\
\midrule
Saber por dónde empezar sin sentirse examinada &
Seis pasos, uno por pantalla: tres preguntas sobre hábitos observables —con qué
frecuencia lee, si ha terminado un libro entero, qué puede contar de una página que acaba
de leer—, intercaladas con los tres fragmentos de la fila siguiente, y
\textbf{cero jerga gramatical} \\
\addlinespace
Juzgarse ante algo real y no ante una descripción &
Tres fragmentos literarios en escalera (\emph{Alice}, Sherlock Holmes,
\emph{Frankenstein}), pintados con la tipografía del lector para que se lean como un
libro y no como un formulario. Pesan 0,65 en el cálculo del nivel: delante de un texto
real la gente se juzga mejor que describiéndose \\
\addlinespace
Que le propongan algo, no una biblioteca &
La pestaña «A mi nivel» de \emph{Encontrar}, que propone lecturas sin escribir nada \\
\addlinespace
Poder usarlo con el pulgar &
Objetivos táctiles de 44~px, controles repartidos en filas que caben en 390~px, mandos
de navegación también debajo del texto \\
\addlinespace
No perderse &
Al terminar el cuestionario, un enlace que salta directamente al buscador \textbf{con
la recomendación ya hecha} \\
\bottomrule
\end{tabularx}
\caption{Necesidades de la usuaria primaria y respuesta del producto. Qué hace hoy en
su lugar: elige por portada o por título conocido, abandona a las tres páginas y
concluye que el problema es ella.}
\label{cuad:primaria}
\end{table}

\subsubsection{Usuario secundario: el adulto autodidacta}

\textbf{Perfil.} 25 a 55 años. Ya lee en inglés con esfuerzo. Su problema no es
empezar, es \textbf{subir de escalón sin abandonar}. Tiene libros a medias y no
recuerda por dónde iba en ninguno.

\begin{itemize}[leftmargin=*, itemsep=3pt]
  \item \textbf{Continuidad.} «Mis libros» guarda cada libro empezado con su barra de
        avance y el capítulo por el que iba, y lo reabre exactamente ahí. Es la
        petición literal que originó el módulo: \emph{«un apartado con mis libros, los
        que estoy leyendo, que se queden guardados, en qué partes voy»}.
  \item \textbf{Saber a qué se enfrenta antes de comprometerse.} La ficha del libro dice
        cuánto tarda en leerlo \textbf{a 180 o 110 palabras por minuto} según le quede por
        debajo o por encima de su nivel, cuánto es diálogo y qué palabras se
        repiten. Las palabras y el tiempo salen del fichero completo; la longitud de frase,
        el diálogo y el vocabulario, de una muestra con tope de 60.000 palabras.
  \item \textbf{Vocabulario que persista.} «Mis palabras» guarda cada consulta con su
        frase y su libro, que es la unidad que sirve para recordar. Deliberadamente
        \textbf{no hay puntos ni rachas}: la evidencia sobre recompensas externas
        \parencite{deci1999meta} indica que premiar una actividad que ya gusta la
        deteriora, y leer ya es gratificante.
  \item \textbf{Pronunciación.} Lee bien y suena mal, y nadie se lo corrige. La
        aplicación marca palabra por palabra qué sonó bien y qué no.
\end{itemize}

\subsubsection{Usuario institucional: el docente de inglés}
\label{sec:usuario-docente}

\textbf{Perfil.} Da clase a grupos heterogéneos y necesita repartir lecturas graduadas.
No puede leerse cuarenta libros para clasificarlos.

\begin{itemize}[leftmargin=*, itemsep=3pt]
  \item Un catálogo \textbf{medido}, no etiquetado a ojo, con la dificultad por capítulo
        para poder asignar fragmentos y no obras completas.
  \item Una declaración honesta de lo que la herramienta no es, para poder usarla sin
        incurrir en una promesa que no puede cumplir: \textbf{LiteraLevel no certifica
        ningún nivel del MCER}.
  \item Reproducibilidad: poder comprobar de dónde salen los números.
\end{itemize}

Este perfil es el que justifica dos decisiones que a un usuario final le dan igual:
publicar el guion de validación y mantener el aviso de alcance visible.

\subsection{El ciclo de usuario real: Esmeralda}
\label{sec:esmeralda}

La única prueba que vale es la que hace alguien que no sabe dónde hay que pulsar.
\textbf{Esmeralda, 15 años, estudiante de preparatoria}, prima del autor, fue la primera
persona ajena al proyecto que usó la aplicación. Lo que dijo, sin edulcorar:

\begin{quote}\itshape
«Está bien difícil de usar en celular.» — ¿Qué parte? — «Casi todo.»\par\medskip
«No me gustan los colores\dots{} la página se siente antigua, todo muy recto, muy
marcado.»\par\medskip
«¿Qué es eso de 10?»\par\medskip
«O yo soy muy tonta para esto.»
\end{quote}

La cuarta frase es la importante, y es la que ordena todo lo que vino después.
\textbf{Cuando un usuario concluye que el tonto es él, el fallo es del diseño.}
\textcite{krug2014dont} formula la misma idea como regla de producto: si el usuario
tiene que pensar para entender un control, el control está mal. A partir de ahí se
diagnosticó qué producía cada queja y se corrigió (cuadro~\ref{cuad:esmeralda}).

\begin{small}
\begin{longtable}{@{}L{3.1cm} L{3.5cm} L{4.6cm} L{2.7cm}@{}}
\toprule
\textbf{Lo que dijo} & \textbf{Diagnóstico} & \textbf{Corrección} & \textbf{Verificación} \\
\midrule
\endfirsthead
\toprule
\textbf{Lo que dijo} & \textbf{Diagnóstico} & \textbf{Corrección} & \textbf{Verificación} \\
\midrule
\endhead
\bottomrule
\endfoot
\bottomrule
\caption{Del comentario de una usuaria a un cambio verificable.}
\label{cuad:esmeralda}
\endlastfoot
«¿Qué es eso de 10?» & La etiqueta «Resultados» junto a un número no dice qué cuenta & «¿Cuántos te propongo?», con explicación & Etiqueta reescrita \\
\addlinespace
Creyó que «Recomiéndame por nivel» tendría en cuenta sus gustos & La etiqueta prometía más de lo que hacía & «Sorpréndeme con algo de mi nivel» & Etiqueta reescrita \\
\addlinespace
— & «Ajustar a mi nivel» es abstracto & «Acercar los resultados a mi nivel · no descarta libros» & Etiqueta reescrita \textbf{dos veces}: la primera versión, «Solo libros que pueda leer», prometía un filtro que la búsqueda no hace \\
\addlinespace
— & «Medir la dificultad de este capítulo» es lenguaje de laboratorio & «¿Cómo de difícil es este capítulo?» & Etiqueta reescrita \\
\addlinespace
— & «Vistazo rápido a la dificultad» no dice cuánto tarda & «Vistazo rápido (unos segundos)» & Etiqueta reescrita \\
\addlinespace
— & «Analizar el libro entero» no advierte del coste & «Leer el libro entero (tarda más)» & Etiqueta reescrita \\
\addlinespace
«No me gustan los colores» & El color es gusto, y el gusto no se discute: se ofrece & \textbf{Cinco paletas} con pestaña propia y cambio instantáneo: \emph{hielo} (la original), \emph{amanecer} (clara, crema y naranja), \emph{menta}, \emph{papel} (para leer horas) y \emph{medianoche} & Una prueba automática comprueba las cinco contra el criterio de contraste \textbf{WCAG AA (4,5:1)} \parencite{w3c2023wcag}. Las cinco pasan en las \textbf{ochenta y cinco parejas} de color que se comprueban —diecisiete por paleta, cada una con su umbral: 4,5:1 para el texto normal y 3:1 para el título grande—. La más ajustada del texto normal mide \textbf{4,53} y la del título, \textbf{3,11} \\
\addlinespace
«Todo muy recto, muy marcado» & Esto no es gusto, es diseño, y se corrige para todas las paletas & Esquinas de 14~px, bordes finos, sombras suaves y las opciones de lista como fichas tocables & — \\
\addlinespace
«Difícil de usar en celular\dots{} casi todo» & En un móvil, el círculo de un botón de radio es un blanco de 12~px & La ficha entera pasa a ser el blanco: \textbf{44~px} & — \\
\end{longtable}
\end{small}

\subsubsection{Lo que ya se había corregido antes, por la misma vía}

El ciclo de Esmeralda no fue el primero. Probar la aplicación en un teléfono real —y no
en el navegador del escritorio— había descubierto antes fallos que \textbf{ninguna de
las 303 pruebas automáticas podía encontrar, porque todos necesitaban un dedo y una
pantalla pequeña}:

\begin{itemize}[leftmargin=*, itemsep=3pt]
  \item El doble toque para consultar una palabra \textbf{no funcionaba en el móvil}: el
        código escuchaba solo el evento del ratón, que iOS no emite al seleccionar con
        el dedo. La instrucción en pantalla prometía algo que el dispositivo no hacía.
        Hoy la palabra se consulta manteniéndola presionada —con ratón, doble clic— y el
        lector lee la selección, que en iOS llega por \cod{selectionchange}.
  \item Los títulos de la tabla de resultados eran enlaces a \cod{gutenberg.org}, así
        que tocarlos \textbf{sacaba al lector fuera de la aplicación}.
  \item El resultado del nivel se pintaba en la parte superior de la página, fuera de la
        pantalla: la usuaria no se enteraba de que había terminado.
  \item Al pedir «traer una frase del capítulo» sin frase disponible, el campo se vaciaba
        y \textbf{borraba lo que el usuario había escrito}, sin decir nada.
  \item «Mis libros» mostraba 24 entradas llamadas \cod{vacio.txt}, \cod{falso.txt} y
        similares: la batería de pruebas anotaba sus ficheros de ataque en la biblioteca
        \textbf{real} de quien las ejecutara. Los datos personales del usuario eran el
        banco de pruebas.
\end{itemize}

\subsubsection{El segundo ciclo: Laura}
\label{sec:laura}

El 12 de septiembre de 2026 la probó \textbf{Laura}, hermana del autor y desarrolladora
web, sobre la aplicación \textbf{desplegada}: primero en la laptop y después en su
teléfono Android. Mandó tres capturas y dos notas de voz, que se transcribieron y se
conservan íntegras en el repositorio (\cod{auditoria/\allowbreak feedback\_laura\_2026-09-12.md}).
Su crítica es de usuaria y de diseñadora a la vez:

\begin{quote}\itshape
«Hago el movimiento de deslizar con el dedo y dos, tres segundos y la página se empieza
a mover.»\par\medskip
«No dejas un espacio entre el cuadro y lo que escribes\dots{} se come el palito de la
P.»\par\medskip
«Doy a siguiente capítulo y me deja ahí mismo abajo.»\par\medskip
«Lo que pones en los tres circulitos son configuraciones, no son funcionalidades.»\par\medskip
«No entiendo el sonido W que es diferente a la letra B.»
\end{quote}

\begin{small}
\begin{longtable}{@{}L{3.1cm} L{3.8cm} L{4.2cm} L{2.8cm}@{}}
\toprule
\textbf{Lo que dijo} & \textbf{Diagnóstico} & \textbf{Corrección} & \textbf{Verificación} \\
\midrule
\endfirsthead
\toprule
\textbf{Lo que dijo} & \textbf{Diagnóstico} & \textbf{Corrección} & \textbf{Verificación} \\
\midrule
\endhead
\bottomrule
\endfoot
\bottomrule
\caption{El segundo ciclo, de la crítica a la verificación.}
\label{cuad:laura}
\endlastfoot
Retraso de 2--3~s al deslizar en el teléfono & \cod{backdrop-filter} en más de cien bloques y fondo con \cod{background-\allowbreak attachment:\allowbreak fixed}: repintado completo en cada fotograma & Los dos efectos, solo en dispositivos con ratón & Con emulación táctil quedan desactivados; \textbf{falta confirmarlo en un teléfono real} \\
\addlinespace
El texto toca el borde redondeado & Los bloques de Markdown de la interfaz de entonces (Gradio) no tienen relleno, pero aquí todos llevan fondo de panel & Relleno interior a los bloques con contenido & Medido y capturado a 390 y 1440~px \\
\addlinespace
En modo oscuro no distingue las secciones & Panel al 72\,\% sobre casi negro y borde al 16\,\% en \emph{hielo}, la paleta por defecto & Panel y borde más claros en las dos paletas oscuras & Las 85 parejas de contraste WCAG siguen pasando \\
\addlinespace
Quiere una barra de navegación con las funciones y los ajustes aparte & \emph{Colores} y \emph{Estado} eran pestañas como leer o pronunciar. \textbf{El profesor había dicho lo mismo}: «Colores toma demasiada presencia» & Botón \textbf{Ajustes} en la cabecera; barra de pestañas fija arriba al bajar y centrada en escritorio & Probado en el navegador; a 390~px «Pronunciar» sigue en el menú «\dots» \\
\addlinespace
«Siguiente capítulo» la deja al final & Los botones no movían la página & Al cambiar de capítulo sube al principio del texto & Teléfono: el lector queda justo bajo la barra con el capítulo nuevo \\
\addlinespace
Borrar a mano la frase para practicar otra & No había forma de empezar de nuevo & Botón \textbf{«Otra frase»} bajo la corrección & Probado en el navegador \\
\addlinespace
No entiende «/w/ diferente a la b» y pide «los símbolos» & La explicación comparaba con una letra, no enseñaba el AFI y ni siquiera sabía qué se había oído & Explicación anclada a palabras del español y a lo que se entendió; símbolo AFI; transcripción de la palabra del diccionario local; 41 audios pregrabados, uno por fonema, que no gastan créditos ($\approx$\,0,19~USD una sola vez) & Probada de punta a punta con «fix» por «fox»: ahora culpa a la vocal y no a la /f/. \textbf{Los audios aún no los ha escuchado nadie} \\
\addlinespace
Captura: «don’t» marcado como «don» y «t» & El apóstrofo tipográfico de los libros partía la contracción & Se normaliza antes de comparar & Dos pruebas nuevas \\
\addlinespace
Captura: «Tolstoy, Leo, graf» & El título nobiliario de Gutenberg se tomaba por parte del nombre & El autor corto se queda con apellido y nombre & Prueba nueva \\
\end{longtable}
\end{small}

\textbf{Dos de las correcciones fallaron la primera vez, y las 306 pruebas automáticas no
lo vieron.} La primera versión del arreglo del lector sacaba la cabecera del libro letra a
letra. La del botón de siguiente capítulo no movía nada, porque el texto nuevo llega
después y los botones que pulsó Laura eran los de debajo del texto. Las dos se
descubrieron midiendo la página en un navegador controlado a ancho de teléfono: es la
misma lección del apartado anterior, repetida.

\textbf{Lo que vino después de ese ciclo cambió la interfaz entera.} Ese mismo 12 de
septiembre, en la revisión del avance, el profesor Barranco pidió que la aplicación
entregada no fuera la de Gradio sino un \emph{front} propio, separado del motor. La
interfaz pasó a Angular sobre una API (apartado~\ref{sec:front}). Las correcciones del
cuadro~\ref{cuad:laura} se hicieron primero en Gradio, que sigue en marcha como motor, y
las que tocan la pantalla se llevaron al \emph{front}: los ajustes fuera de la navegación,
el lector que sube al cambiar de capítulo, «Otra frase» y la ficha de cada sonido.

\subsubsection{El tercer ciclo: Ismael}
\label{sec:ismael}

Los días 14 y 15 de septiembre de 2026 la probó \textbf{Ismael}, hermano del autor, que
abrió en un teléfono Android el enlace que le mandaron por WhatsApp, sin que nadie le
explicara nada y ---cosa que se supo después, y que cambia la lectura de lo que
contó--- \textbf{sin iniciar sesión}. Mandó dos notas de voz:

\begin{quote}\itshape
«Cuesta un poquito entenderle\dots{} le falta hacerla más intuitiva, o a lo mejor unos
cuadros de diálogo que te diga qué hace cada cosa.»\par\medskip
«Entré y ya inicié sesión, pero le pisas al micrófono y no te dice que hables, no se
activa la grabación\dots{} a lo mejor no tan saturada la página, los textos están largos y
tienes que hacerle scroll, y te pierdes, no sabes qué hacer.»
\end{quote}

\begin{small}
\begin{longtable}{@{}L{3.1cm} L{3.8cm} L{4.2cm} L{2.8cm}@{}}
\toprule
\textbf{Lo que dijo} & \textbf{Diagnóstico} & \textbf{Corrección} & \textbf{Verificación} \\
\midrule
\endfirsthead
\toprule
\textbf{Lo que dijo} & \textbf{Diagnóstico} & \textbf{Corrección} & \textbf{Verificación} \\
\midrule
\endhead
\bottomrule
\endfoot
\bottomrule
\caption{El tercer ciclo, de la nota de voz al código publicado.}
\label{cuad:ismael}
\endlastfoot
«Le falta hacerla más intuitiva\dots{} unos cuadros de diálogo que te diga qué hace cada
cosa» & Recorrido a 390~px sin cuenta: «Entrar» abría \emph{Ajustes} por arriba y «Tu
perfil» quedaba fuera de la pantalla; el botón de la cabecera era un icono sin texto; la
tarjeta «Mídeme» prometía una estimación desde la primera pregunta que ya no ocurre &
Primera tanda: «Entrar» abre «Tu perfil» con el cursor en «Usuario», el botón dice
\textbf{Ajustes} y «Mídeme» pierde la promesa. Tercera tanda: bienvenida de tres pasos la
primera vez en Inicio, con «Saltar» y «Ver la bienvenida otra vez» en \emph{Ajustes} &
5 pruebas de \cod{entrar-perfil} y 6 de \cod{bienvenida} en el \emph{front}; recorrido en
el navegador a 390~px \\
\addlinespace
«Le pisas al micrófono y no te dice que hables, no se activa la grabación» &
\textbf{La causa de su caso está en revisión.} Al investigar la queja apareció un fallo
real y medido: en el \emph{front} anterior, el aviso de error del micrófono y el
«Grabando\dots{}» se pintaban \textbf{debajo de la barra fija} ---de 750 a 849, con la
barra empezando en 687---, así que no se veían. Pero Ismael estaba \textbf{sin sesión}, y
sin cuenta no hay corrección de pronunciación; no está establecido que lo que él vio fuera
ese fallo, ni consta que su navegador fuera un WebView & Una línea de estado
\textbf{dentro} de la barra, con un mensaje por caso —«Grabando: habla ahora», permiso
pendiente, permiso negado, sin micrófono, micrófono ocupado, sin frase— (\cod{601788b}).
El aviso para el WebView de Android se queda como precaución, no como explicación de este
caso & 11 pruebas de \cod{pronunciar.\allowbreak microfono} y el navegador a 390~px.
\textbf{Falta medir qué ve alguien sin sesión al pulsar el micrófono} \\
\addlinespace
«No tan saturada la página, los textos están largos\dots{} te pierdes» & Medido a 390~px:
el lector enseñaba dos avisos a la vez y dejaba ver unas dos líneas del libro, y
\emph{Pronunciar} ponía 46 palabras antes de la caja de la frase & Segunda tanda: un solo
aviso en el lector —«Mantén presionada una palabra para ver qué significa; desliza para
cambiar de página.»— y, en \emph{Pronunciar}, una línea de ayuda en lugar del subtítulo;
las dos se cierran con una cruz y vuelven con «?» & Comprobado en el \emph{front}
publicado a 390~px táctil el 15 de septiembre \\
\end{longtable}
\end{small}

\textbf{La lección de este ciclo es la contraria de la que parecía.} Al principio se dio
por hecho que «no se activa la grabación» era la barra fija tapando el aviso, y así se
escribió; después se supo que Ismael no había iniciado sesión, y esa causa dejó de
sostenerse. Quedan en pie dos cosas distintas: \textbf{lo que se arregló}, que estaba mal lo
viera él o no, y \textbf{lo que su queja sigue señalando}, que es que la aplicación no dice
con suficiente claridad qué se puede hacer sin cuenta. Esa medición está en marcha y la
causa se escribirá cuando esté, no antes. Y de los sitios donde se midió saturación queda
sin tocar el índice de la ficha, a casi dos pantallas del principio.

\subsubsection{La lección, en una línea}

\textbf{El valor de una prueba de usuario no está en el promedio de satisfacción, sino
en la frase que duele.} Con $n = 3$ no se puede estimar nada, pero se pueden encontrar
los fallos que hacen que una persona se culpe a sí misma, y ya hay una coincidencia: a los
tres les costó el teléfono, y a las dos lectoras no les gustó la paleta por defecto
(cuadros~\ref{cuad:esmeralda}, \ref{cuad:laura} y~\ref{cuad:ismael}). \textcite{nielsen2000five}
sitúa en cinco usuarios el punto donde se detecta la mayor parte de los problemas de
usabilidad; el accionable A8 lleva la prueba hasta ahí.

\subsection{Usabilidad del modelo}

Un modelo es usable cuando el usuario puede confiar en él sin entenderlo. Eso se
construye con cinco decisiones, todas medidas, y con la interfaz que las pone delante del
lector (apartado~\ref{sec:front}).

\subsubsection{La espera se traslada, no se pide}

Medir un libro largo son minutos. Pedirle esos minutos al lector en el momento en que
quiere leer es perderlo. La estrategia es \textbf{mover el cómputo antes de la
consulta}: el catálogo se precalcula fuera de línea y, cuando el lector busca, el
trabajo ya está hecho. Los números que hacen viable esa decisión están en el
cuadro~\ref{cuad:velocidad} y en la figura~\ref{fig:velocidad}: 8,19 ventanas/s en CPU,
112,7 en la GPU local y \textbf{2.022 en una A100}, con lo que los 61.642 libros en
inglés de Project Gutenberg son \textbf{unas tres horas} de cómputo puro. De extremo
a extremo son \textbf{unas diecisiete}, y la diferencia no la pone la tarjeta sino el
disco; el §8.3 lo desglosa.

Cuando el cómputo sí tiene que ocurrir en vivo, se ofrecen \textbf{dos velocidades y se
dice cuál se usó}: un vistazo (30 capítulos repartidos, 2 ventanas de cada uno) o el
libro entero sin muestrear. Y el atajo está cuantificado, no prometido: 8 ventanas dan
\textbf{1,784} y el recorrido completo del mismo capítulo, \textbf{1,743}.

\subsubsection{La incertidumbre se comunica en bandas}

El modelo devuelve un número continuo, pero su escala está comprimida
($\text{predicho} \approx 0{,}72 \times \text{real} + 0{,}38$; véase la
figura~\ref{fig:compresion}). Presentar «1,27 sobre 2» como si fueran dos decimales de
precisión sería mentir con exactitud. La interfaz muestra \textbf{bandas} —básico
$< 0{,}6 \leq$ intermedio $< 1{,}4 \leq$ avanzado, tal como las dibuja la
figura~\ref{fig:escala}— y la letra MCER únicamente como orientación, siempre junto al
número y con la tabla de correspondencia a la vista.

Un fallo real de esta decisión, encontrado y corregido, ilustra el nivel de cuidado que
exige: el redondeo por defecto de Python resuelve los empates hacia el par, de modo que
un 1,00 exacto caía en un nivel y la tabla publicada decía otro. Un usuario que
comparase su cifra con la tabla habría encontrado dos respuestas para el mismo número.
Ahora hay una sola regla, y una prueba fija los cinco límites.

\subsubsection{Nada falla en silencio, y nada se rompe del todo}

La regla transversal del sistema es la \textbf{degradación elegante}: sin clave de
servicio externo, la aplicación sigue dando diccionario, fonética, buscador y
puntuación, porque nada de eso sale a la red; el audio de Wiktionary es un enlace a
\cod{upload.wikimedia.org} que descarga el navegador del lector, así que necesita
conexión, pero no clave ni servicio de pago. Solo se apagan tres cosas —la explicación
en contexto, la corrección de pronunciación y el resumen de «¿de qué trata?» de la
ficha—, y se dice con palabras que el lector
entienda —«se ha agotado el presupuesto mensual de explicaciones; el diccionario, la
fonética y el audio siguen funcionando»— en lugar del error técnico en crudo.

La regla de diseño es que los fallos vuelvan como texto legible en vez de excepciones
sin tratar. Las pruebas cubren escenarios concretos; no demuestran que ningún manejador
pueda fallar. El corte de verificación actual se distingue de las baterías históricas
en el apartado~\ref{sec:front}.

\subsubsection{Cada número dice de dónde viene}

En la ficha del libro conviven dos clases de información y \textbf{la separación se ve
en pantalla}: lo \emph{contado} sobre el fichero (palabras, tiempo de lectura,
proporción de diálogo, vocabulario) sale solo al abrirlo, es instantáneo y no necesita
red; el \emph{resumen} escrito por un modelo generativo va detrás de un botón, se
guarda en disco y viene marcado como lo que es. \textbf{Ninguna cifra de la ficha la
dice un modelo generativo.} En el diccionario, la regla equivalente: el modelo de
lenguaje explica el contexto y la gramática sin sustituir una acepción utilizable de
Wiktionary. Si no existe, desde el 13 de septiembre puede proponer una definición
marcada como generada por IA, solo con acceso autorizado. La definición de diccionario y el texto generado no son lo mismo.
Y si la elección entre acepciones es dudosa ---menos de 0,05 de afinidad
entre las dos más cercanas--- y no hay explicación de OpenAI, desde \cod{bf3a9f2}
la ficha no afirma una: dice «Puede ser» y enseña la acepción principal de
cada entrada del diccionario, si hay al menos dos.

\subsubsection{La latencia se reparte para que la pantalla no se quede quieta}

En una consulta léxica nueva, casi todo el tiempo es \textbf{espera al servicio
externo}: la explicación contextual se midió en \textbf{14,9 segundos} con el esfuerzo
de razonamiento \cod{low} que usa la aplicación —eran 67,7 con el esfuerzo medio de
antes—, mientras todo el trabajo local —lema, diccionario, desambiguación, fonética—
son \textbf{40 milisegundos} de trabajo en el servidor; en el teléfono, con la red de
por medio, la hoja tarda unos 0,7~s en caliente y unos 5~s la primera consulta de la
sesión (medido el 15 de septiembre de 2026 a 390~px). Por eso la ficha se pinta en dos
etapas: lo local aparece primero y la explicación se añade cuando llega. El trabajo tarda lo mismo; deja de
parecer que la aplicación se colgó. La misma idea gobierna la curva de dificultad, que
se repinta cada pocos capítulos mientras mide (el repintado cuesta 25~ms, medido) y los
recursos pesados, que se cargan al primer uso y no al arrancar.

\subsubsection{La interfaz entregada: un \emph{front} de Angular sobre una API}
\label{sec:front}

Hasta el 12 de septiembre de 2026 la aplicación era una sola pieza: ocho pestañas de
Gradio que hacían de interfaz y de motor a la vez. En la revisión del avance de ese día,
el profesor Barranco pidió que la entrega fuera un \emph{front} propio, separado del
motor. Se hizo al día siguiente, y es lo que se entrega:

\begin{itemize}
  \item \textbf{El \emph{front}} es una aplicación de Angular 22 (componentes
        independientes, sin \cod{zone.js}) publicada como \textbf{Space estático} de
        Hugging Face (\cod{MikeMaki1763/\allowbreak literalevel-front}), gratuito y siempre
        despierto: es HTML, CSS y JavaScript ya compilados. Tiene cuatro sitios en vez de
        ocho pestañas —\emph{Leer}, \emph{Encontrar}, \emph{Pronunciar} y \emph{Tu
        nivel}—, más \emph{Mis palabras}; los ajustes (paleta, tamaño de letra, idioma de
        las explicaciones) y el estado del servicio salen de la navegación, que es lo que
        Laura y el profesor habían pedido por separado.
  \item \textbf{La API} es un servicio FastAPI (\cod{app/api.py}, 40 rutas bajo \cod{/api} a fecha de \cod{8ee42f0}, el \cod{app/} que publica el Space en \cod{7c09d90}) \textbf{montado en el mismo Space de Gradio} que ya existía, con Gradio
        en la raíz como motor y respaldo. \textbf{Reutiliza el motor y ayudantes de Gradio},
        pero sus adaptadores contienen lógica propia: topes, cobro, aislamiento y
        degradación requieren pruebas también en la API. Compartir funciones no basta
        para demostrar equivalencia. El inicio de sesión frena los intentos fallidos en una hora: 15 por IP
        y, desde \cod{514786b}, 20 por usuario, porque la IP sale de una cabecera que
        puede escribir el cliente. Solo cuentan los que fallan, y son valores por
        defecto que se cambian por variable de entorno. Medir tampoco es ilimitado: cada medida ocupa la CPU del único proceso del Space, así que desde el 14 de septiembre hay topes por IP y hora ---40 medidas de nivel o de capítulo y 6 curvas nuevas--- y como mucho 4 curvas en cola a la vez; lo ya medido no gasta ficha. En el Space, a fecha de \cod{2bb1276}, la séptima curva seguida dio 429 y un libro ya medido siguió respondiendo. La identidad del lector
        viaja en una cabecera (\cod{X-Lector}) y sus datos siguen en el bucket montado en
        \cod{/data}. El primer montaje falló ---el proxy de renderizado en servidor de
        Gradio no dejaba pasar \cod{/api}--- y el segundo, sin ese proxy, es el que está en
        marcha.
  \item \textbf{Regla de entrega:} una función existe cuando tiene ruta en la API
        \textbf{y} pantalla en el \emph{front}. Lo que solo vive en Gradio no cuenta como
        entregado.
  \item \textbf{Verificación del 13 de septiembre, corte actualizado el 14:} 32 de 35
        pruebas API pasaron en el entorno aislado; tres quedaron sin validar por
        dependencias NLP ausentes. Pasaron además 7 pruebas de concurrencia de voz,
        10 del despliegue y 11 del pipeline de portadas. No se suman como una prueba
        completa de usuario ni sustituyen la batería histórica del motor. El backend
        publicado \cod{e56cf08} se cotejó por hash y se comprobó salud JSON, lectura de
        Alicia en Angular, diccionario sin perfil y rechazo de voz anónima. No se probó
        OpenAI real ni micrófono físico en ese cierre. Evidencia:
        \cod{auditoria/\allowbreak despliegue\_codex\_\allowbreak 2026-09-13\_noche.md}.
        Recuento sin claves del 15 de septiembre sobre el árbol del despliegue
        \cod{0319290} (\cod{1643c4b}, cada suite corrida sola): 66 de 66 pruebas API,
        351 pruebas extremas del motor sin ningún fallo y 48 de 48 de corpus
        (\cod{DESPLIEGUE.md}). Las 34 de 34 de lectura natural son del corte anterior,
        \cod{d58b0be}.
\end{itemize}

\medskip\noindent\textbf{Lo que el \emph{front} añadió a la usabilidad en un día.} Tres
cosas, y las tres salieron de mirar cómo se usa un teléfono:

\begin{itemize}
  \item \textbf{Uso con una sola mano.} Un estudio propio
        (\cod{investigacion/\allowbreak una\_mano.md}) fijó diez reglas a partir de la
        literatura sobre el alcance del pulgar y el tamaño de los objetivos; las primeras ya
        están en el \emph{front}: en el lector, una barra única abajo con capítulo y página
        anterior y siguiente y el índice, que se esconde con la navegación; nada primario
        en el cuarto superior de la pantalla; en \emph{Pronunciar}, una barra de acción
        fija abajo con el micrófono en el centro, que no se mueve entre grabar, parar y
        corregir.
  \item \textbf{Corrección de pronunciación con fonema determinista.} Con la misma
        grabación, el modelo generativo señalaba un sonido distinto en cada llamada. Ahora
        el fonema que falla sale de \textbf{alinear la transcripción de lo que había que
        decir con la de lo que se entendió} ---las dos de diccionario: Wiktionary y, donde
        no la trae, CMUdict \parencite{cmudict}, 126.052 palabras bajo licencia BSD--- y el consejo de boca,
        lengua y aire sale de una ficha fija por fonema. Del modelo queda el resumen.
        La selección del fonema es determinista dadas las mismas palabras y
        transcripciones; esto no hace determinista al reconocimiento de voz.
  \item \textbf{Animación del corte de la boca.} Para cada uno de los 41 sonidos del
        inglés americano (y los 44 del británico), un corte sagital animado ---labios,
        dientes, paladar, velo, lengua y mandíbula--- que va del reposo a la posición del
        sonido y vuelve; y para la palabra entera, la secuencia de sus sonidos a partir de
        su transcripción. Son SVG animados con CSS, sin JavaScript, generados una vez en el
        servidor y servidos por la API; el \emph{front} los enseña en la ficha de cada
        sonido y en la corrección ---primero el sonido que falla y después la palabra---, y
        con «reducir movimiento» del sistema los deja quietos.
\end{itemize}

\noindent\textbf{Boca 3D y portadas, a fecha del \emph{front} \cod{ec957d0} (repositorio \cod{4e6276e}).} La boca 3D ya está publicada: el \emph{front} sirve \cod{boca3d/cara.glb} y \cod{lengua.glb}, idénticos byte a byte a los del repositorio, con 72 mallas entre los dos y respaldo 2D si no cargan. La piel de la media cabeza sale de la malla base de MakeHuman 1.3.0, publicada bajo CC0 1.0 (no obliga a atribuir) \parencite{makehuman_licencia}; paladar, velo, encías, dientes y lengua son modelos propios, y sus 53 poses se validan contra ellos. Es un corte esquemático, no anatomía validada clínicamente. Las portadas publicadas son \textbf{tipográficas}: una de 12 familias según los temas del catálogo, con una paleta de cinco colores sacada de la huella del identificador y el título a un contraste de al menos 4,5:1. El lector toma el fondo y el acento de esa paleta, con el texto a 7:1 o más y el acento a 4,5:1 sobre el fondo (sobre el color suave del acento, la pista del lector baja a 4,17:1 en algunos libros en tema claro, como SWV369), y el ajuste «Color del libro en el lector» lo apaga; la prueba del \emph{front} que fija esos contrastes recorre 12 familias, 200 libros y dos temas (\cod{b2003e5}). Las portadas FLUX no se generaron: el pipeline
exige saldo, tarifa actual, aprobación de muestras y permiso separado para publicar.
La selección local de 1.200 libros contiene 746 motivos temáticos y 454 genéricos.
Su calidad visual sigue sin comprobarse. El enlace del lector es
\url{https://mikemaki1763-literalevel-front.static.hf.space/}.

El 14 de septiembre se volvieron a contar con pandas \textbf{53.179 IDs únicos y 585.111 capítulos} en el catálogo local. A fecha de \cod{740eb84}, el catálogo externo de textos libres tiene 15.367 filas y los vectores de búsqueda suman 68.546 (53.179 + 15.367). El Space publicado cuenta los mismos 68.546 libros en \cod{/api/estado} a fecha de \cod{384a206}, el reinicio 21. Es una comprobación del artefacto, no una
nueva inferencia. Las 40 rutas API se contaron en el código a fecha de \cod{8ee42f0}. En el cierre remoto
anterior, \emph{house} eligió un sentido no adecuado para una frase sobre la ubicación
de un edificio: disponibilidad del diccionario no demuestra desambiguación correcta
en todos los casos. Estos límites también forman parte de la entrega.

\textbf{Lectura natural, medida con diez frases.} En \emph{Pronunciar}, además del sonido que falla, la aplicación mide el ritmo de la lectura y enseña cómo se dice la frase de corrido: énfasis, enlaces entre palabras y reducciones (\emph{want to} \textrightarrow{} \emph{wanna}). Lo propone un modelo generativo y lo filtra un validador: las reducciones solo salen de una lista cerrada y los enlaces se comprueban con la pronunciación de diccionario. Con diez frases reales del catálogo evaluadas a mano (\cod{auditoria/\allowbreak bateria\_lectura\_natural\_\allowbreak 2026-09-14.md}), el modelo propuso 11 reducciones y 29 enlaces; llegaron al lector 2 reducciones y 4 enlaces, \textbf{todos correctos}, y las diez versiones naturales quedaron fieles a la frase (8 reconstruidas por el validador). La exhaustividad es baja: se quedan fuera enlaces claros como \emph{an argument}. En esa batería todo lo que llegó al lector era correcto, y eso no es una garantía: desde el 15 de septiembre el validador exige además un verbo detrás de \emph{gonna}, \emph{wanna} o \emph{gotta} (antes aceptaba «I am gonna school» si el modelo lo proponía) y no enlaza a través de un punto. La frase de la demostración, «Don't you want to come with us?», da en el Space a fecha de \cod{7c09d90} las reducciones \emph{doncha} y \emph{wanna}, marcadas como generadas por IA, y ningún enlace. Esa comprobación es anterior a la pantalla de hoy, donde la lectura natural va plegada bajo «Para sonar más natural»; como la propone un modelo generativo, la respuesta puede cambiar, y el guion de la demostración la vuelve a comprobar con cuenta antes de enseñarla.

\subsection{Accionables}

Diez acciones concretas (cuadro~\ref{cuad:accionables}), ordenadas por puntos de rúbrica
protegidos u obtenidos por hora de trabajo. Cada una lleva coste real y criterio de
terminado. No son «se podría»: son «se hace».

\begin{small}
\begin{longtable}{@{}c L{3.2cm} L{3.6cm} L{2.0cm} L{2.9cm}@{}}
\toprule
\textbf{\#} & \textbf{Acción} & \textbf{Por qué} & \textbf{Coste} & \textbf{Terminado cuando} \\
\midrule
\endfirsthead
\toprule
\textbf{\#} & \textbf{Acción} & \textbf{Por qué} & \textbf{Coste} & \textbf{Terminado cuando} \\
\midrule
\endhead
\bottomrule
\endfoot
\bottomrule
\caption{Accionables, con coste y criterio de terminado.}
\label{cuad:accionables}
\endlastfoot
\textbf{A1} & \textbf{Publicar el modelo y el diccionario} como repositorios abiertos & Desbloquea que el cuaderno corra en un Colab ajeno. Es el mayor rendimiento por hora de todo el proyecto & $\sim$30 min; gratuito. Requiere un token que solo el autor puede crear & \textbf{HECHO a medias}: el modelo está en \cod{MikeMaki1763/\allowbreak literalevel-v4}, público y con la descarga sin credenciales verificada; el diccionario (243 MB) sigue sin publicar \\
\addlinespace
\textbf{A2} & \textbf{Desplegar en enlace permanente} & El enlace anterior era un túnel temporal que moría al apagar el equipo, y era el mayor riesgo de la entrega & \textbf{7 USD/mes}: el Space de Gradio que aloja el motor y la API exige cuenta Hugging Face Pro (9 USD de tarifa, 7 con descuento); el hardware \cod{cpu-basic} no añade coste, y el Space estático del \emph{front} es gratuito & \textbf{HECHO}: \cod{MikeMaki1763/\allowbreak literalevel-\allowbreak front} (el \emph{front}) y \cod{MikeMaki1763/\allowbreak literalevel} (API y motor), públicos y en marcha; el disco persistente es A4 \\
\addlinespace
\textbf{A3 · HECHO} & \textbf{Terminar la remedición del catálogo} con el modelo v4, libro entero & Sustituye las etiquetas heredadas —predicciones de un modelo v1 sobre las primeras $\sim$190 palabras— por medidas reales & T4 gratuita; la descarga por el volcado de Common Pile (57.226 textos en 7,7 minutos) y los resultados en Hugging Face, no en Drive, que pedía un permiso en cada sesión (\S\ref{sec:escalado}) & \textbf{Cumplido el 12 de septiembre de 2026}: 53.179 libros y 585.111 capítulos medidos, y la aplicación los usa \\
\addlinespace
\textbf{A4 · HECHO} & \textbf{Montar almacenamiento persistente} para los datos del lector & Un servicio que duerme y despierta con la estantería vacía deja en ridículo «Mis libros» & 18 USD por TB y mes en repositorio privado; los datos del lector ocupan 10,9 MB medidos, así que el coste es despreciable & \textbf{Cumplido y verificado}: bucket montado en \cod{/data}; se fijó un perfil, se reinició el Space y el perfil seguía ahí \\
\addlinespace
\textbf{A5} & \textbf{Conseguir un micrófono} (auriculares con micro) & El equipo actual no tiene. Sin él no se puede demostrar la corrección de pronunciación, que es la función más vistosa y la que mejor encaja con «usabilidad para el usuario final» & Coste mínimo de hardware & Se graba una frase desde el equipo del autor y la corrección palabra por palabra aparece en pantalla \\
\addlinespace
\textbf{A6 · HECHO} & \textbf{Sustituir el catálogo heredado} de 1.197 registros & Incluía 1.049 registros sin texto descargable, el 87,6\pct{} de aquella versión, no de la actual & Incluido en A3 & El catálogo medido actual contiene 53.179 IDs de Gutenberg; a fecha de \cod{db366b8} se excluyen al recomendar y al buscar 600 que no son lectura (índices, catálogos, directorios) y, desde \cod{3524357}, 163 textos libres que no son lectura en inglés (relleno o en otra lengua), y por nivel solo se ofrecen obras literarias. Una medida del precálculo se corrige al cargar (\cod{correcciones\_\allowbreak catalogo.csv}): PG2242, \emph{A Midsummer Night's Dream} del First Folio, guardó 0,697 y releída por el mismo camino da 1,713 (C1) \\
\addlinespace
\textbf{A7} & \textbf{Calibrar la escala} con una transformación afín de dos parámetros ajustada sobre entrenamiento, sin tocar los pesos & La escala está comprimida y un A2 puntúa cerca de B1. La técnica ya se aplicó con éxito a una variante anterior & Una tarde & Las medias por nivel siguen estrictamente crecientes \textbf{y} se acercan a sus anclas \\
\addlinespace
\textbf{A8} & \textbf{Repetir la prueba de usuario con cinco lectores} de perfiles distintos sobre Angular, con guion de tareas y observación en silencio & Los dos ciclos previos encontraron fallos graves; no validan todos los cambios del nuevo front \parencite{nielsen2000five} & Cinco sesiones de 20 min & Existe una lista priorizada de hallazgos y las tres primeras correcciones están hechas \\
\addlinespace
\textbf{A9} & \textbf{Instrumentar la métrica que de verdad importa}: proporción de libros empezados que se terminan, con recomendación y sin ella & Todo lo demás son métricas de proceso. Esta es la única que valida la propuesta de valor & Una jornada de instrumentación, respetando el aislamiento por sesión & Hay una cifra base con la que comparar dentro de un mes \\
\addlinespace
\textbf{A10} & \textbf{Ampliar el corpus de anclaje MCER} más allá de los 654 documentos actuales & Es el eslabón más delgado de la cadena de evidencia y el que limita la confianza en la letra que se muestra & Trabajo de corpus, sin coste de cómputo & El corpus de anclaje supera los 1.000 documentos con niveles de fuente institucional \\
\end{longtable}
\end{small}

\noindent\textbf{Orden de ejecución actualizado:} ensayar la voz con micrófono real
(A5), recorrer Angular con usuarios (A8), grabar el video y comprobar los archivos
de entrega. A2, A3, A4 y la sustitución del catálogo de A6 ya tienen evidencia de
cierre; no se vuelven a presentar como trabajo por empezar. Las mejoras visuales
pendientes no deben desplazar el ensayo de la demo.

\subsection{Métricas de éxito}

\begin{table}[htbp]
\centering
\small
\begin{tabularx}{\textwidth}{@{}L{1.9cm} L{6.0cm} Y@{}}
\toprule
\textbf{Nivel} & \textbf{Métrica} & \textbf{Estado} \\
\midrule
Modelo & Orden estrictamente creciente por nivel MCER en corpus externo & \textbf{Cumplido}: 0,5040 · 0,5561 · 0,9276 · 1,2414 · 1,4952 \\
Modelo & Señal por encima de la longitud del texto & \textbf{Cumplido}: correlación parcial 0,668 \\
Producto & Un usuario sin formación técnica completa el recorrido & \textbf{Cumplido con reservas, sobre la interfaz anterior}: lo completó en Gradio, y de esa sesión salieron las nueve correcciones del Cuadro~1 de la estrategia ---seis de ellas etiquetas---; sobre el \emph{front} de Angular, pendiente (A8) \\
Producto & Accesibilidad de color & \textbf{Cumplido}: las ochenta y cinco parejas de color de las cinco paletas superan WCAG AA; las más ajustadas, 4,53 (texto) y 3,11 (título) \\
Producto & Robustez & \textbf{Parcial}: recuento sin claves del 15 de septiembre sobre el árbol del despliegue \cod{0319290} ---API 66 de 66, 351 pruebas extremas sin ningún fallo y corpus 48 de 48---, más lectura natural 34 de 34 y voz concurrente 10 de 10 del corte anterior \cod{d58b0be}; no equivalen a una prueba completa de usuario \\
Operación & Enlace permanente disponible & \textbf{Cumplido}: el \emph{front} en \cod{MikeMaki1763/\allowbreak literalevel-front} y la API con el motor en \cod{MikeMaki1763/\allowbreak literalevel}, los dos en Hugging Face Spaces, públicos y en marcha, con los datos del lector en un bucket montado en \cod{/data} \\
Operación & Catálogo medido, no heredado & \textbf{Cumplido} el 12 de septiembre de 2026: 53.179 libros y 585.111 capítulos medidos por el v4 (A3) \\
Impacto & Proporción de libros terminados & \textbf{Sin instrumentar} (A9) \\
\bottomrule
\end{tabularx}
\caption{Métricas por nivel, con su estado real.}
\label{cuad:metricas}
\end{table}

\subsection{Riesgos}

\begin{table}[htbp]
\centering
\small
\begin{tabularx}{\textwidth}{@{}L{5.2cm} L{2.4cm} Y@{}}
\toprule
\textbf{Riesgo} & \textbf{Impacto} & \textbf{Mitigación} \\
\midrule
El enlace de la demostración muere durante la evaluación & Alto & A2: despliegue permanente, probado desde un dispositivo ajeno antes de la fecha \\
El servicio externo de lenguaje agota su presupuesto en mitad de una demostración & Medio & Ya mitigado: la aplicación degrada con un mensaje legible y conserva diccionario, fonética, audio, buscador y puntuación \\
Que la letra MCER se lea como una certificación & Alto, reputacional & Declaración de alcance en el documento y aviso permanente en pantalla \\
Que la remedición del catálogo no termine a tiempo & \textbf{Resuelto} & Terminada el 12 de septiembre de 2026: sin Drive, con los resultados en Hugging Face tras cada lote, 0 errores y ningún libro medido dos veces \\
Sobrecarga del espejo de descarga de Project Gutenberg & Medio, ético & Seis hilos y no treinta, cada uno con su pausa entre peticiones; y el vocabulario se obtuvo del SPGC, no rascando el sitio \\
Contenido inapropiado para una lectora menor de edad & Medio & Decisión razonada de \textbf{no filtrar por contenido} y avisar en la ficha; se argumenta en el apartado~\ref{sec:producto} \\
\bottomrule
\end{tabularx}
\caption{Riesgos y mitigaciones.}
\label{cuad:riesgos}
\end{table}

% ===========================================================================
\section{Descripción de los datos}
\label{sec:datos}
% ===========================================================================

Todos los datos de este trabajo provienen de \textbf{corpus académicos publicados,
APIs y descarga directa de repositorios de dominio público}. No se usó Kaggle como
fuente de datos en ningún punto del proyecto, conforme a la restricción del módulo.
Cuando un corpus tiene licencia restrictiva se dice, y su texto crudo nunca sale del
entorno de cómputo: solo estadísticas agregadas.

\subsection{Inventario de fuentes}

\begin{table}[htbp]
\centering
\footnotesize
\begin{tabularx}{\textwidth}{@{}L{2.4cm} L{3.2cm} L{2.2cm} L{2.8cm} Y@{}}
\toprule
\textbf{Fuente} & \textbf{Qué aporta} & \textbf{Tamaño} & \textbf{Origen} & \textbf{Licencia / acceso} \\
\midrule
\textbf{CLEAR Corpus} & Etiqueta continua de legibilidad (\cod{BT\_easiness}) para el modelado supervisado & 4.724 extractos & Repositorio público del proyecto CLEAR \parencite{crossley2023clear} & Uso investigador, descarga directa \\
\addlinespace
\textbf{One\-Stop\-English} & Corpus de \textbf{afinado} de la variante v2 (151 artículos), validación ordinal en 3 niveles sobre los 38 reservados y \textbf{validación externa del v4} (560 unidades) & 567 textos (189 por nivel) & GitHub \parencite{vajjala2018onestop} & CC BY-SA 4.0 \\
\addlinespace
\textbf{British Council Learn\-English} & Validación externa contra niveles MCER oficiales & 59 textos A1--C1 & Descarga de los PDF públicos, con la URL y el SHA-256 de cada uno en un manifiesto & Con derechos: no se distribuye ni se imprime texto, pero se reconstruye y se verifica byte a byte \\
\addlinespace
\textbf{SPGC} (Standar\-dized Project Guten\-berg Corpus) & Lista global de frecuencia de palabras (solo conteos, sin texto) & 45.345 libros en inglés de 55.905 & Zenodo, registro 2422561 \parencite{gerlach2018spgc,gerlach2020standardized} & Abierta \\
\addlinespace
\textbf{Project Gutenberg} & Catálogo de obras a medir y a leer dentro de la app & 79.337 entradas; 61.642 libros de texto en inglés, instantánea del 7 de septiembre de 2026 & Catálogo oficial \parencite{gutenberg2026} + espejo \cod{gutenberg.\allowbreak pglaf.org} & Dominio público \\
\addlinespace
\textbf{Kaikki / Wiktio\-nary} & Diccionario, IPA británica y americana, audio & 1.487.639 filas en SQLite (243 MB); IPA en 46.794 y audio en 7.003 & Volcado de kaikki.org \parencite{ylonen2022wiktextract} & Libre \\
\addlinespace
\textbf{ELG CEFR English} & Corpus de anclaje MCER: afinado y validación interna del v4 & 712 documentos; 654 aceptados & European Language Grid, descargado de Hugging Face \parencite{breuker2023cefr} & No comercial: CC BY-NC-SA 4.0 según el guion de entrenamiento; CC BY-NC 4.0 según su ficha en UniversalCEFR \\
\addlinespace
\textbf{Story\-Weaver} (Pratham Books) & Textos libres del catálogo, medidos con el v4 & 14.506 textos & Plataforma abierta \parencite{prathambooks_storyweaver} & CC BY 4.0 \\
\addlinespace
\textbf{Let's Read} (The Asia Foundation) & Textos libres del catálogo, medidos con el v4 & 328 textos & Biblioteca digital \parencite{asiafoundation_letsread} & CC BY 4.0 (318) y CC BY 3.0 (10) \\
\addlinespace
\textbf{African Story\-book} (Saide) & Textos libres del catálogo, medidos con el v4 & 280 textos & Biblioteca digital \parencite{saide_africanstorybook} & CC BY \\
\addlinespace
\textbf{VOA Learning English} & Textos libres del catálogo, medidos con el v4 & 253 textos & Sitio de Voice of America \parencite{voa_learningenglish} & Dominio público \\
\addlinespace
\textbf{Common Pile} (\cod{project\_\allowbreak gutenberg}) & Texto de los libros de Gutenberg para medir el catálogo & 57.226 textos descargados & Hugging Face \parencite{kandpal2025commonpile} & Dominio público \\
\addlinespace
\textbf{CMUdict} & Fonética americana donde Wiktionary no la trae (corrección de voz) & 126.052 palabras & GitHub de Carnegie Mellon University \parencite{cmudict} & BSD de dos cláusulas \\
\addlinespace
\textbf{Catálogo heredado} (Story\-Weaver / Global Digital Library) & Registros del catálogo anterior, retirado el 12 de septiembre de 2026; no son los 14.506 textos de arriba & 1.049 registros & API pública de contenidos & Abierta, pero sin descarga de texto \\
\bottomrule
\end{tabularx}
\caption{Fuentes de datos, con volumen y procedencia. Los recuentos de textos libres son de \texttt{740eb84}, el catálogo publicado en el reinicio 21.}
\label{cuad:fuentes}
\end{table}

\subsection{Relevancia de cada fuente}

\textbf{CLEAR} es el corpus central del modelado. Su variable objetivo,
\cod{BT\_easiness}, no es una fórmula de legibilidad: es una escala construida a partir
de comparaciones pareadas juzgadas por lectores humanos, lo que evita que el modelo
aprenda a reproducir un contador de sílabas. Además trae precalculadas once métricas
clásicas (Flesch, Flesch-Kincaid, ARI, SMOG, Dale-Chall, CAREC, CAREC\_M, CML2RI,
conteos de palabras, oraciones y párrafos) y dos variables de estratificación
—\cod{Categ} (literario/informativo) y \cod{Lexile Band}— que permiten comprobar si el
modelo funciona igual de bien en narrativa que en divulgación. Esa comprobación por
dominio se hizo y aparece en el apartado~\ref{sec:dominio}.

\textbf{OneStopEnglish} aporta algo que CLEAR no puede: el mismo artículo reescrito en
tres niveles. Como el contenido es idéntico y solo cambia la complejidad lingüística,
permite medir el orden dentro de cada artículo, una prueba mucho más exigente y limpia
que comparar textos distintos entre sí.

\textbf{British Council} es la única fuente con niveles MCER asignados por una
institución de enseñanza de idiomas, que es exactamente la escala que el producto
muestra al usuario. Es el juez final.

\textbf{SPGC} existe por una razón de higiene: descargar decenas de miles de libros uno
a uno desde \cod{gutenberg.org} viola su política de acceso automatizado. El SPGC
publica los conteos de palabras ya calculados, con lo que se obtiene la distribución
léxica del catálogo completo sin tocar el servidor de Gutenberg.

\textbf{Kaikki/Wiktionary} sustituyó a una API comercial de diccionario. La decisión fue
de producto y de derechos: un SQLite local de 1,48 millones de entradas no tiene límite
de peticiones, no cuesta dinero por consulta y no puede desaparecer si un proveedor
cambia sus condiciones. Se descartó explícitamente Oxford Dictionaries.

\subsection{Calidad de los datos: lo que está bien y lo que no}
\label{sec:calidad}

Un apartado de calidad que solo diga «los datos están limpios» no informa. Estos son los
problemas reales encontrados y qué se hizo con cada uno.

\paragraph{(a) El corpus de validación MCER no es homogéneo.} De los 59 textos de
British Council, una parte no es prosa: en A1 hay un menú de restaurante, un horario de
estudio, un tablero de salidas de aeropuerto, tarjetas de visita y mensajes de texto.
Se definió un filtro de prosa continua por dos vías independientes —por palabras clave
del título y por estructura ($\geq 7$ palabras por línea y densidad de puntuación final
de frase en el intervalo $(0;\,0{,}11]$)—, que coinciden en 53 de 59 casos, el 89,8\pct.

\textbf{Ese filtro es el del análisis, no el que aplica la aplicación}, y conviene no
confundirlos porque dan resultados distintos. Lo que corre en producción es
\cod{scoring.continuous\_prose}, con otras reglas: al menos 150 palabras, al menos cuatro
frases, entre 5 y 55 palabras de media por frase, y rechazo si más del 17,5\% de las
palabras viven en líneas de maquetado —viñetas o líneas de cuatro palabras o menos— o más
del 35\pct{} de líneas cortas. Coinciden en el propósito —tirar menús, listas y tablas—
pero no en los umbrales.

\textbf{Y aquí hay que contar la secuencia real, porque el filtro no se fijó antes de
mirar ningún resultado.} El orden fue el contrario: primero se puntuaron los 59 textos, se vio
que el modelo daba resultados absurdos con el menú y el tablero de salidas, y \textbf{de
ese diagnóstico salió la idea del filtro}. Que las dos vías se definieran sin volver a
mirar las puntuaciones, y que coincidan entre sí en el 89,8\pct{} de los casos, es lo que
impide que sea un filtro hecho a medida de las notas. Pero no es un criterio \emph{a
priori}, y llamarlo así sería regalarse una garantía que no se tiene.

\paragraph{(b) Las etiquetas del catálogo heredado no eran fiables, y por eso se sustituyó.}
El catálogo con el que la aplicación funcionó hasta el 12 de septiembre de 2026 tenía
\textbf{1.197 libros} —no
12.752, como sugiere un \cod{wc -l} ingenuo: el campo \cod{hook} contiene saltos de
línea entrecomillados y hay que contar con un lector de CSV real—. Su columna
\cod{nivel\_cefr} \textbf{no es anotación humana ni metadato del editor}: son
predicciones de un modelo anterior (v1) sobre las primeras $\sim$190 palabras de cada
obra, con cortes MCER obtenidos por cuantiles. El resultado es un reparto que no resiste
el sentido común (cuadro~\ref{cuad:catalogo} y figura~\ref{fig:catalogo}).

\begin{table}[htbp]
\centering
\small
\begin{tabular}{@{}l R{2.2cm} R{2.2cm}@{}}
\toprule
\textbf{Nivel} & \textbf{Libros} & \textbf{\%} \\
\midrule
A1 & 723 & 60,4 \\
A2 & 185 & 15,5 \\
C2 & 115 & 9,6 \\
B1 & 69 & 5,8 \\
C1 & 53 & 4,4 \\
B2 & 52 & 4,3 \\
\midrule
\textbf{Total} & \textbf{1.197} & \textbf{100} \\
\bottomrule
\end{tabular}
\caption{Reparto de niveles del catálogo heredado. Son predicciones de un modelo v1
sobre las primeras $\sim$190 palabras, no anotación fiable.}
\label{cuad:catalogo}
\end{table}

% ------------------------------------------------------- FIGURA: catálogo
\begin{figure}[htbp]
\centering
\begin{tikzpicture}
\begin{axis}[
  barrall,
  height=5.4cm,
  xbar,
  bar width=8pt,
  xmin=0, xmax=830,
  xlabel={Libros del catálogo heredado},
  symbolic y coords={C1,B2,B1,C2,A2,A1},
  ytick=data,
  y dir=reverse,
  enlarge y limits=0.14,
  nodes near coords={\pgfmathprintnumber[fixed,precision=0,1000 sep={.}]{\pgfplotspointmeta}},
  every node near coord/.append style={font=\scriptsize, color=llgris, anchor=west},
  xmajorgrids=true,
]
\addplot[fill=llrojo!35, draw=llrojo!80] coordinates {
  (723,A1) (185,A2) (69,B1) (52,B2) (53,C1) (115,C2)};
\draw[llgris!70, dashed, line width=0.7pt]
  (axis cs:199.5,C1) -- (axis cs:199.5,A1);
\node[font=\scriptsize\itshape, color=llgris, anchor=south west, align=left]
  at (axis cs:210,B1) {reparto uniforme\\(199,5 libros por nivel)};
\end{axis}
\end{tikzpicture}
\caption[Reparto de niveles del catálogo heredado]{El sesgo del clasificador v1, visto de
una vez. Seis de cada diez libros de dominio público quedaron etiquetados A1 y el tramo
B1--B2 —donde está la mayoría de los lectores reales— reúne 121 de 1.197 obras. Una
distribución así no describe la literatura en dominio público: describe el clasificador
que la produjo. Datos del cuadro~\ref{cuad:catalogo}.}
\label{fig:catalogo}
\end{figure}

Hay cuentos infantiles etiquetados C1 y clásicos muy distintos entre sí todos en C2. Y
el tramo B1--B2, donde está la mayoría de los lectores reales, eran 121 libros de 1.197.
Esas etiquetas servían para ordenar y filtrar de forma aproximada, \textbf{no eran una
medida fiable libro a libro}. Es la razón de ser del precálculo descrito
en~\ref{sec:escalado}: remedir todo el catálogo, libro entero, con el modelo final.

\textbf{Y se remidió.} El 12 de septiembre de 2026 terminó: \textbf{53.179 libros de
Project Gutenberg y 585.111 capítulos}, medidos enteros por el v4, sin un solo libro medido
dos veces. Es lo que usa hoy la aplicación, y su reparto dice lo contrario que el heredado
(cuadro~\ref{cuad:catalogomedido}).

\begin{table}[htbp]
\centering
\small
\begin{tabular}{@{}l R{2.4cm} R{3.2cm}@{}}
\toprule
\textbf{Nivel} & \textbf{Libros medidos} & \textbf{Libros literarios} \\
\midrule
A1 & 1 & \\
A2 & 40 & \\
B1 & 557 & 288 entre A1 y B1 \\
B2 & 3.127 & 1.676 \\
C1 & 23.455 & \\
C2 & 25.999 & \\
\midrule
\textbf{Total} & \textbf{53.179} & \\
\bottomrule
\end{tabular}
\caption{Reparto del catálogo medido por el v4 (media 1,742 sobre 2), con la corrección de \cod{correcciones\_catalogo.csv} aplicada
(PG2242 pasa de B1 a C1), que es como lo cuenta la aplicación. «Literarios» son los
que la aplicación ofrece al recomendar por nivel: ficción, cuentos, lecturas escolares,
poesía y teatro. Entre los libros «fáciles» había índices, anuarios, mapas y tablas de
números, que puntúan bajo precisamente porque son listas.}
\label{cuad:catalogomedido}
\end{table}

La literatura del XIX es, sobre todo, C1--C2, y un lector B1 apenas tiene libros enteros a
su medida. Lo que sí tiene son \textbf{capítulos}: 3.644 en 1.499 libros quedan a su
alcance, y la aplicación se los ofrece. Dos cautelas: la letra es orientativa, porque la
escala del modelo sale comprimida y desplazada hacia arriba, y el reparto describe el
catálogo de Gutenberg, no la dificultad de la lectura en inglés en general.

\textbf{Y fuera de Gutenberg.} A fecha de \cod{740eb84}, publicado en el reinicio 21, el catálogo suma además \textbf{15.367 textos de licencia libre} medidos con el v4: 14.506 de StoryWeaver, 328 de Let's Read, 280 de African Storybook y 253 de VOA. Por nivel, 766 A1, 5.696 A2, 5.491 B1, 2.524 B2, 777 C1 y 113 C2. Desde ese reinicio, un bloque de tres o más cuentos se mide también sin cada uno de ellos: si quitar uno solo mueve la medida más de 0,5 (un margen fijado antes de medir la tanda), los demás se quedan con la medida sin él, y ese texto conserva la del bloque entero ---la regla no lo baja de nivel--- y queda marcado. Son 92 textos, y frente al reinicio 19 cambian de letra 225: 42 suben y 183 bajan. En dos de los bloques cuyos cuentos subieron de A2 a B2, quitar el texto marcado al bloque tal como estaba antes deja la medida en A2, así que esas subidas están por revisar (\cod{54d6591}). Lo que se quedó fuera se contó en \cod{377a650}, cuando eran 15.357: otros 894 textos libres se quedaron fuera; la salida de \cod{cerrar\_corpus.py} dice: «medidos: 15357 (texto 11039, bloque 4318) | sin medida, fuera: 894». Recontados sobre las medidas en disco: 776 tienen menos de 150 palabras, 878 cayeron en un bloque que el filtro de prosa continua rechazó (260 bloques: 215 con pinta de lista, menú o tabla, 32 con frases demasiado cortas o largas y 13 con menos de cuatro frases), 15 no llegaron a formar bloque y 1 se rechazó solo; 828 son de StoryWeaver. Es donde está la lectura de nivel bajo que Gutenberg no tiene: 11.953 textos entre A1 y B1, frente a 598 libros. Cuatro decisiones sobre ellos, publicadas a fecha de \cod{7c09d90}. \textbf{(1)} Un cuento corto medido dentro de un bloque lleva la dificultad del bloque, no la suya, y la API lo marca (\cod{medido\_en\_bloque}) para que la tarjeta y la ficha lo digan. \textbf{(2)} La atribución que exige su licencia se enseña al final del texto, pero no se puntúa: medida con el cuento, esa línea en español con su URL subía «Doggie is Happy» de 0,007 a 0,349 en «¿Cómo de difícil es?»; en el Space da 0,007. \textbf{(3)} El texto que se sirve sale sin caracteres U+FEFF, sin párrafos repetidos seguidos y sin páginas duplicadas dentro de un párrafo; el que se midió no cambia, y \cod{catalogo\_externo.csv} sale idéntico byte a byte al regenerarlo. Cambian 1.442 textos servidos, 185 de ellos por repeticiones, y los estribillos de una frase se conservan. \textbf{(4)} Como no tienen descargas de Gutenberg, al recomendar por nivel ocupan un tercio de las tarjetas (puestos 2, 5, 8\ldots), ordenados por cercanía al nivel del lector; con 30 tarjetas, A2 recibe 11 y B1 10, sin cuentos en común entre las dos franjas.

\paragraph{(c) De aquellos 1.197 libros, solo 148 se podían abrir.} Eran los de Project
Gutenberg. Los otros 1.049 provenían de StoryWeaver y no tenían fichero descargable, de
modo que aparecían en el desplegable de lectura como opciones muertas. Peor aún: su
identificador se había usado para construir enlaces a Gutenberg extrayendo sus dígitos,
con lo que \cod{SWgdlgdl201} apuntaba al libro 201 de Gutenberg, que es otra obra. Un
enlace al libro equivocado es peor que ningún enlace. \textbf{En el catálogo medido los
53.179 libros se pueden abrir}, y la aplicación los descarga en \cod{.txt}, el mismo texto
con el que se midieron, para que los capítulos del lector sean los de la curva.

\paragraph{(d) La licencia de Project Gutenberg se estaba puntuando como si fuera
novela.} El texto legal viene pegado al último capítulo: en \emph{Alice's Adventures in
Wonderland}, 2.900 de las 5.042 palabras del capítulo XII son licencia. Eso contaminaba
tanto la puntuación como las estadísticas de la ficha (\cod{gutenberg} y \cod{project}
aparecían entre las palabras más repetidas del libro). Se recorta por las marcas
START/END dentro del capítulo, sin descartar el capítulo, que contiene el final de la
obra.

\paragraph{(e) La portada no es un capítulo.} En un \cod{.txt} de Gutenberg, las 4.957
palabras previas al capítulo I (portada, índice, pies de ilustración) se estaban
puntuando como si fueran la novela. En el análisis de \emph{War and Peace}, el «capítulo
más fácil del libro» resultó ser la portadilla. Se descarta con la etiqueta que el
propio lector de ficheros asigna, «Antes del primer capítulo». Se probó descartar por el
filtro de prosa continua y \textbf{no sirve}: se llevaba por delante dos capítulos reales
de \emph{Alice}.

% ===========================================================================
\section{Análisis exploratorio}
\label{sec:eda}
% ===========================================================================

Siguiendo la organización canónica del módulo, el análisis se presenta en tres niveles:
univariado, bivariado y multivariado. Todas las figuras y cuadros citados se generan en
el cuaderno que acompaña a este documento.

\subsection{Análisis univariado}

\textbf{La variable objetivo.} \cod{BT\_easiness} es continua y aproximadamente
simétrica alrededor de cero. No hay clases desbalanceadas que corregir porque no hay
clases: el problema es de regresión. Cuando se necesitó una versión categórica para
comparar modelos, se construyó por cuantiles (\cod{qcut}), lo que garantiza seis grupos
equifrecuentes y hace que la exactitud sea comparable entre modelos.

\textbf{Las variables textuales derivadas.} Sobre cada extracto se calcularon, a partir
de sus \emph{tokens}, seis variables de ingeniería propia: longitud media de oración,
\emph{type-token ratio}, longitud media de palabra, proporción de palabras raras,
sílabas medias por palabra y porcentaje de polisílabos. Todas son positivas y con cola
derecha, como es habitual en variables de conteo sobre texto.

\textbf{La distribución del vocabulario inglés.} El agregado de los 45.345 libros del
SPGC arroja \textbf{2.838.258.803 palabras} y un vocabulario único de
\textbf{3.295.146 formas}. La distribución sigue la ley de Zipf
\parencite{zipf1949human} con precisión: unas pocas miles de palabras cubren la mayor
parte del texto corrido. Este hecho, que parece de trivia, tiene una consecuencia
práctica que se documenta en~\ref{sec:bivariado}.

\textbf{El catálogo.} El cuadro~\ref{cuad:catalogo} y la figura~\ref{fig:catalogo}
muestran el reparto de niveles del catálogo heredado. Su forma —60\pct{} en A1— es en sí
misma un hallazgo del análisis exploratorio: una distribución así no describe la
literatura en dominio público, describe el sesgo del clasificador que la produjo. El
catálogo medido por el v4 (cuadro~\ref{cuad:catalogomedido}) lo confirma por el lado
contrario: 53.179 libros, y solo 598 por debajo de B2.

\textbf{La dificultad dentro de un libro.} La unidad de análisis del producto no es el
libro sino el capítulo, y la figura~\ref{fig:curva_wp} muestra por qué: los 362
capítulos de \emph{War and Peace}, medidos sin muestrear, se reparten entre 1,333 y
1,990 con media 1,735. La figura~\ref{fig:curvas_comp} añade dos libros más y confirma
que la dispersión interna es la norma, no la excepción.

\subsection{Análisis bivariado}
\label{sec:bivariado}

\textbf{Cada variable frente a la dificultad.} La matriz de correlación de las variables
derivadas contra \cod{BT\_easiness} confirma el sentido esperado: los textos más fáciles
tienen oraciones más cortas, palabras más cortas, menos polisílabos y menor proporción
de palabras raras. La correlación de \cod{rare\_word\_ratio} con \cod{BT\_easiness} es
negativa, como debe ser.

\textbf{El experimento de saturación léxica.} La variable «palabra rara» depende de qué
lista de palabras familiares se use. Se probaron tres, de menor a mayor
(cuadro~\ref{cuad:saturacion}).

\begin{table}[htbp]
\centering
\small
\begin{tabular}{@{}l R{2.0cm} R{5.2cm}@{}}
\toprule
\textbf{Lista de palabras familiares} & \textbf{Libros} & \textbf{Correlación de \cod{rare\_word\_ratio} con \cod{BT\_easiness}} \\
\midrule
Corpus Gutenberg de NLTK & 18 & $-0{,}379$ \\
Muestra propia de Gutenberg & 397 & $-0{,}330$ \\
\textbf{SPGC completo, en inglés} & \textbf{45.345} & $\mathbf{-0{,}326}$ \\
\bottomrule
\end{tabular}
\caption{Efecto de escalar el corpus de referencia léxica. La correlación no mejora.}
\label{cuad:saturacion}
\end{table}

Este es un \textbf{resultado negativo y se reporta como tal}: multiplicar por 2.500 el
corpus de referencia no mejora la señal, e incluso la baja marginalmente. La exactitud
del clasificador de tres clases se mantuvo en 0,51 con las tres listas. La explicación
está en la ley de Zipf: la distribución de frecuencias del inglés converge muy rápido,
así que las 3.000 palabras más frecuentes de 18 libros y las de 45.345 libros son casi
el mismo conjunto. El valor del paso no está en la métrica sino en la cobertura: la
lista del SPGC generaliza mejor a texto nuevo arbitrario, que es lo que la aplicación
recibe.

\textbf{Dificultad frente a longitud, la variable incómoda.} Sobre los 47 textos de
British Council medibles en local, contar las palabras del texto —sin modelo alguno—
alcanza una correlación de Spearman de \textbf{0,807} con el nivel oficial. El modelo
alcanza 0,860. Es decir: en ese corpus, nivel y longitud están confundidos, porque los
materiales de examen gradúan también la extensión de la tarea. Publicar ese
\emph{baseline} antes de que lo pregunte el evaluador es más creíble que omitirlo; en el
apartado~\ref{sec:baseline} se cuantifica cuánta señal aporta el modelo por encima de él.

\textbf{Dificultad frente a dominio.} El nivel del texto está correlacionado con su tipo
documental. Al filtrar los 59 textos de British Council a prosa continua, el nivel A1
desaparece casi por completo y el rango evaluable \textbf{de esa tanda de análisis} se
encoge de cinco niveles a cuatro (A2--C1): el catálogo gratuito de British Council no
tiene C2, así que el corpus nunca tuvo seis. La reproducción local del
cuadro~\ref{cuad:validacion} sí conserva los cinco, con tres textos en A1, porque
filtra con otras reglas. Por eso el \emph{baseline} de contar palabras cae de 0,879 a 0,765 al filtrar:
parte de la señal del conjunto completo era simplemente que los textos A1 son cortos.
Las correlaciones antes y después del filtro \textbf{no son comparables entre sí}, y se
reportan por separado. Los conjuntos de textos que aparecen aquí y en el
apartado~\ref{sec:validacion} están todos, con lo que se puede y no se puede hacer con cada
uno, en el cuadro~\ref{cuad:conjuntos}.

\subsection{Análisis multivariado}

\textbf{Las variables clásicas se estorban entre sí.} Las métricas de legibilidad
tradicionales (Flesch, Flesch-Kincaid, ARI, SMOG, Dale-Chall) son transformaciones
distintas de los mismos dos ingredientes —longitud de oración y longitud de palabra— y
por tanto están fuertemente correlacionadas entre ellas. Meterlas todas en un modelo
lineal no añade información: añade multicolinealidad. Por eso el conjunto de trabajo se
redujo a \textbf{tres} variables poco redundantes entre sí y con interpretación
distinta: \cod{avg\_sentence\_length} (complejidad sintáctica),
\cod{rare\_word\_ratio} (carga léxica) y \cod{type\_token\_ratio} (diversidad de
vocabulario).

\textbf{Hasta dónde llegan las variables clásicas.} Con esas tres variables, un
clasificador de seis clases alcanza 0,29 de exactitud con regresión logística
\parencite{pedregosa2011scikit} y 0,28 con XGBoost \parencite{chen2016xgboost}. Que un
ensamble de árboles —capaz de capturar interacciones no lineales arbitrarias— no supere
a un modelo lineal es la prueba de que \textbf{el cuello de botella no era el modelo sino
la representación del texto}. Tres números por documento no bastan para separar seis
clases.

Conviene decir qué sostiene ese argumento y qué no. Lo que se comprueba es que
\textbf{subir la capacidad del algoritmo sobre estas variables no mueve el resultado}, y
eso se ve en cuatro modelos distintos. Lo que no se demuestra es que exista un límite
teórico: nadie ha probado todas las combinaciones posibles de conteos, y el cuadernillo
enseña un caso donde ampliar de tres a diecisiete variables sí mejora. La conclusión
honesta es empírica, no un techo.

\textbf{Representación densa.} Al sustituir las tres variables por un \emph{embedding} de
384 dimensiones de \cod{all-MiniLM-L6-v2} \parencite{reimers2019sentence}, la exactitud
salta a 0,37 y el error ordinal medio baja de 1,39 a 0,93 niveles. El modelo híbrido
(\emph{embedding} + las tres variables clásicas) da 0,37 y 0,91: prácticamente lo mismo.
Esto indica que \textbf{el \emph{embedding} ya contiene la información que aportaban las
tres variables}, y bastante más. La figura~\ref{fig:progresion} recoge la progresión
completa.

\textbf{Estructura por dominio.} La comparación por subconjuntos de CLEAR (literario
frente a informativo, y por banda Lexile) muestra que el rendimiento no es uniforme: el
modelo funciona mejor en texto informativo que en literario, y peor en la banda más
fácil (Lexile 500) que en la más difícil (Lexile 1300). Los números están en el
cuadro~\ref{cuad:dominio}.

% ===========================================================================
\section{Procesamiento de datos}
\label{sec:procesamiento}
% ===========================================================================

Cada transformación se documenta con su justificación. La regla que se siguió es la del
módulo 1 del diplomado: \textbf{no se transforma sin medir qué se rompió}.

\subsection{Sobre el texto de entrada}
\label{sec:proc_texto}

El cuadro~\ref{cuad:textos} recoge las ocho transformaciones que sufre el texto antes de
llegar al modelo; la figura~\ref{fig:pipeline} las sitúa en la cadena completa.

\begin{small}
\begin{longtable}{@{}c L{3.2cm} L{3.4cm} L{5.6cm}@{}}
\toprule
\textbf{\#} & \textbf{Transformación} & \textbf{Por qué} & \textbf{Evidencia / consecuencia} \\
\midrule
\endfirsthead
\toprule
\textbf{\#} & \textbf{Transformación} & \textbf{Por qué} & \textbf{Evidencia / consecuencia} \\
\midrule
\endhead
\bottomrule
\endfoot
\bottomrule
\caption{Transformaciones sobre el texto, con su justificación y su evidencia.}
\label{cuad:textos}
\endlastfoot
T1 & \textbf{Filtro de prosa continua} antes de puntuar: se rechaza si más del 17,5\% de las palabras viven en líneas de maquetado o más del 35\% de las líneas son cortas & El modelo solo vio prosa; puntuar un menú da un número sin significado & El filtro estaba muerto en producción: se normalizaba \emph{antes} de comprobar, y \cod{normalize()} colapsa los saltos de línea, así que el detector recibía siempre una sola línea. Una carta de menú de 432 palabras se puntuaba como prosa. Corregido: el filtro va sobre el texto crudo \\
\addlinespace
T2 & \textbf{Normalización de espacios} (\cod{\textbackslash s+} $\rightarrow$ espacio) después del filtro & Estabiliza el troceado en bloques & Sin efecto sobre la puntuación; solo sobre el corte \\
\addlinespace
T3 & \textbf{Troceado en bloques de 250 palabras, mínimo 150} & Es la unidad con la que se entrenó; un bloque más corto no tiene señal suficiente & Constante del entrenamiento: cambiarla invalidaría las métricas reportadas \\
\addlinespace
T4 & \textbf{Truncado a 256 \emph{tokens}} & Límite de contexto del modelo & Medido: de 200 bloques de 250 palabras, el \textbf{100\pct} llega exactamente a 256 \emph{tokens} tras truncar, porque en prosa inglesa real 250 palabras superan siempre los 256 BPE. Dentro de un texto troceado así no hay relleno; agrupando ventanas de capítulos distintos sí lo hay, y eso se trata en el ~\S\ref{sec:escalado} \\
\addlinespace
T5 & \textbf{Recorte de la licencia de Gutenberg} por las marcas START/END, dentro del capítulo & 2.900 de las 5.042 palabras del capítulo XII de \emph{Alice} son texto legal & Sin esto, el modelo puntuaba texto jurídico como narrativa y la ficha listaba «gutenberg» entre las palabras más frecuentes del libro \\
\addlinespace
T6 & \textbf{Descarte de portada, índice y portadilla} por la etiqueta «Antes del primer capítulo» & 4.957 palabras de paratexto se puntuaban como novela & Se probó descartar por el filtro de prosa continua y se rechazó: eliminaba dos capítulos reales de \emph{Alice} \\
\addlinespace
T7 & \textbf{Muestreo de hasta 8 bloques repartidos por el documento} para la puntuación rápida & Dar un número en segundos, no en minutos & Cuantificado: 8 bloques dan 1,784 y el recorrido completo de las 23 ventanas del mismo capítulo da 1,743. El atajo se desvía 0,041 sobre una escala de 0 a 2 \\
\addlinespace
T8 & \textbf{Recorte de entradas del usuario en el servidor} (40 caracteres por palabra, 400 por oración, 20 MB por fichero y, desde \cod{9eb3c91}, 100 MB descomprimidos por EPUB) & El límite del navegador es una sugerencia, no una defensa & Una cadena de 3.000 caracteres pasaba el filtro de «máximo cuatro palabras» y se pintaba como encabezado: 9.312 caracteres llenando la pantalla del móvil \\
\end{longtable}
\end{small}

\subsection{Sobre las variables del modelado clásico}

\begin{table}[htbp]
\centering
\small
\begin{tabularx}{\textwidth}{@{}c L{4.6cm} Y@{}}
\toprule
\textbf{\#} & \textbf{Transformación} & \textbf{Por qué} \\
\midrule
T9 & \textbf{Partición agrupada por libro} (\cod{GroupShuffleSplit} sobre \cod{Title}, 20\pct{} de prueba, semilla 42) & Un mismo libro aporta varios extractos a CLEAR. Con una partición aleatoria por fila, fragmentos del mismo libro caerían a ambos lados y el modelo memorizaría el estilo del autor en lugar de aprender dificultad. Es fuga de datos por diseño, y se evita por diseño \\
\addlinespace
T10 & \textbf{Estandarización} (\cod{StandardScaler}) ajustada solo en entrenamiento & Ajustar el escalador con los datos de prueba filtra su distribución al modelo \\
\addlinespace
T11 & \textbf{Discretización por cuantiles} (\cod{qcut}) para las versiones categóricas & Produce grupos equifrecuentes, con lo que la exactitud entre modelos es comparable y el azar tiene una probabilidad conocida ($1/6 = 16{,}7\pct$) \\
\addlinespace
T12 & \textbf{Lista de palabras familiares construida solo con el conjunto de entrenamiento} & Si se cuentan también las palabras de los extractos de prueba, esas palabras entran en la definición de una variable de entrada: es fuga de datos. Medida la diferencia, resulta despreciable —cambian 228 de 3.000 palabras y los resultados se mueven en la tercera decimal—, pero era fuga y se quitó \\
T13 & \textbf{Lista alternativa a escala}: las 3.000 más frecuentes del SPGC & Sustituye una lista de 18 libros por una de 45.345. Es un experimento de cobertura, no la lista que usan los modelos del cuadernillo, y \textbf{el sistema desplegado no consulta ninguna lista de palabras}: puntúa con el transformador \\
\bottomrule
\end{tabularx}
\caption{Transformaciones sobre las variables tabulares.}
\label{cuad:tabulares}
\end{table}

\subsection{Sobre el corpus de vocabulario (procesamiento distribuido)}

Los conteos del SPGC son 45.345 ficheros pequeños. Agregarlos secuencialmente es un mal
uso del tiempo y del sistema de ficheros, de modo que la agregación se hizo con
\textbf{PySpark} \parencite{zaharia2016spark}: los ficheros se distribuyen en 200
particiones, un \cod{flatMap} los convierte en pares (palabra, conteo) y un
\cod{reduceByKey} los suma en una única lista global de frecuencia. Es el patrón canónico
de conteo distribuido y es lo que hace tratable un corpus de 2.838 millones de palabras
en una sesión de Colab.

\subsection{Sobre el diccionario}

El volcado de Kaikki llega como JSONL de líneas independientes. Se filtra a entradas en
inglés con palabra no vacía y se ingesta en SQLite con un esquema plano (\cod{palabra},
\cod{pos}, \cod{sentidos} como JSON, \cod{ipa\_uk}, \cod{ipa\_us}, \cod{audio\_uk},
\cod{audio\_us}) y un índice por \cod{palabra}. La justificación es de latencia: la
primera consulta sobre 1.487.639 filas tarda \textbf{0,001~s} con el índice puesto y la
base abierta en modo solo lectura. Los 243 MB del fichero no cuestan nada en tiempo de
respuesta.

Dos datos de cobertura que condicionan el diseño de la aplicación. Sobre las 1.487.639
entradas, solo \textbf{7.003 (0,5\pct) traen audio de pronunciación} y \textbf{46.794
(3,1\pct) traen transcripción fonética}; con las dos variantes ---la británica y la
americana--- quedan \textbf{33.635 (2,3\pct)}. Pero eso cuenta entradas, y el diccionario
es en su mayoría palabra rara mientras que un libro es en su mayoría palabra corriente.
Medido por el lado del lector ---400 palabras tomadas al azar del texto de \emph{Crime and
Punishment} y resueltas por el mismo camino que usa la aplicación, lematización
incluida---, \textbf{el 58\pct{} llega con fonética y el 32\pct{} con audio grabado}.

Las dos mediciones apuntan a lo mismo desde sitios distintos: \textbf{dos de cada tres
consultas necesitan síntesis de voz}, y esa es la medición que justificó paralelizar las
dos llamadas de red de una consulta en lugar de encadenarlas. La ficha de la palabra lo
dice cuando toca: si no hay ni transcripción ni audio ni clave para generarlo, escribe
\emph{«de esta palabra no hay pronunciación»} en lugar de dejar el hueco en blanco.

% ===========================================================================
\section{Modelación}
\label{sec:modelacion}
% ===========================================================================

\subsection{Diseño experimental}

Todos los modelos comparten la \textbf{misma partición agrupada por libro} (3.774
extractos de entrenamiento y 950 de prueba) y la misma semilla, de modo que las
diferencias entre ellos son atribuibles al modelo y no al reparto de los datos. Se
comparan sobre los mismos 950 extractos reservados, con dos métricas: exactitud sobre
seis clases y \textbf{error ordinal medio}, que penaliza proporcionalmente a la distancia
del fallo.

Dos precisiones sobre ese conjunto, porque cambian cómo hay que leer las cifras. La
primera: esos 950 extractos son también los que sirvieron para \textbf{elegir la época}
del transformador, de modo que funcionan como validación de selección y no como un test
independiente; sus números pueden ser algo optimistas. La segunda: las seis clases salen
de partir \cod{BT\_easiness} en \textbf{sextiles de este corpus}, no de anotaciones del
Marco Común Europeo. Se les dan nombres A1--C2 para poder ordenarlos, pero no son niveles
MCER certificados ni equivalen a los del modelo desplegado, que se ancla con otro corpus.
Las dos cosas son ciertas.

\subsection{Progresión de modelos}
\label{sec:progresion}

\begin{table}[htbp]
\centering
\small
\setlength{\tabcolsep}{4pt}
\begin{tabular}{@{}l l R{2.1cm} R{2.1cm} l@{}}
\toprule
\textbf{Modelo} & \textbf{Representación del texto} & \textbf{Exactitud (6 clases)} & \textbf{Error ordinal medio} & \textbf{Tanda} \\
\midrule
Regresión logística & 3 variables clásicas & 0,29 & 1,39 & histórica \\
XGBoost & 3 variables clásicas & 0,28 & 1,42 & histórica \\
Regresión logística & 3 variables clásicas & \textbf{0,313} & \textbf{1,295} & recalculada \\
Regresión logística & \textbf{17 variables} & \textbf{0,352} & \textbf{1,062} & recalculada \\
XGBoost & \textbf{17 variables} & \textbf{0,337} & \textbf{1,091} & recalculada \\
Regresión logística & \emph{Embedding} MiniLM (384 d) & 0,37 & 0,93 & histórica \\
Regresión logística & Híbrido (\emph{embedding} + 3 variables) & 0,37 & 0,91 & histórica \\
\textbf{\cod{roberta-base} afinado} & \textbf{Texto crudo, atención} & \textbf{0,46} & \textbf{0,67} & histórica \\
\bottomrule
\end{tabular}
\caption{Progresión de modelos sobre los mismos 950 extractos reservados, con la tanda de
la que procede cada fila. La mejora acompaña siempre al enriquecimiento de la
representación, no a la complejidad del clasificador. Las filas recalculadas son las que
imprime el cuadernillo en la ejecución del 16 de septiembre de 2026; su tercera decimal se
mueve con la versión de las bibliotecas.}
\label{cuad:progresion}
\end{table}

% ---------------------------------------------------- FIGURA: progresión
\begin{figure}[htbp]
\centering
\begin{tikzpicture}
\begin{groupplot}[
  group style={group size=2 by 1, horizontal sep=0.6cm},
  barrall,
  width=0.47\linewidth,
  height=5.6cm,
  xbar,
  symbolic y coords={Transformer,Hibrido,MiniLM,XGBoost,LogReg},
  ytick={LogReg,XGBoost,MiniLM,Hibrido,Transformer},
  yticklabels={Reg. logística,XGBoost,\emph{Embedding} MiniLM,Híbrido,\cod{roberta-base}},
  enlarge y limits=0.16,
  every node near coord/.append style={font=\scriptsize, color=llgris, anchor=west},
]

\nextgroupplot[
  xmin=0, xmax=0.60,
  xlabel={Exactitud sobre 6 niveles \emph{(más es mejor)}},
  title={Exactitud},
  nodes near coords={\pgfmathprintnumber[fixed,precision=2]{\pgfplotspointmeta}},
]
\addplot[fill=llcian!35, draw=llcian!85] coordinates {
  (0.29,LogReg) (0.28,XGBoost) (0.37,MiniLM) (0.37,Hibrido) (0.46,Transformer)};
\draw[llrojo!80, dashed, line width=0.7pt]
  (axis cs:0.1667,LogReg) -- (axis cs:0.1667,Transformer);

\nextgroupplot[
  xmin=0, xmax=1.75,
  xlabel={Error ordinal medio \emph{(menos es mejor)}},
  title={Error ordinal},
  yticklabels={,,,,},
  ytick={LogReg,XGBoost,MiniLM,Hibrido,Transformer},
  nodes near coords={\pgfmathprintnumber[fixed,precision=2]{\pgfplotspointmeta}},
]
\addplot[fill=llambar!45, draw=llambar!90] coordinates {
  (1.39,LogReg) (1.42,XGBoost) (0.93,MiniLM) (0.91,Hibrido) (0.67,Transformer)};

\end{groupplot}
\end{tikzpicture}
\caption[Progresión de modelos: la representación manda]{La tesis metodológica del
trabajo, en dos ejes. Los dos primeros modelos comparten \emph{representación} (tres
variables clásicas) y solo cambian de clasificador: XGBoost no mejora a la regresión
logística, incluso empeora el error ordinal. Los saltos ocurren cuando cambia la
representación del texto: \emph{embedding} denso primero, texto crudo con atención
después. La línea roja discontinua del panel izquierdo marca el azar
($1/6 = 0{,}167$), que es lo que acertaría un clasificador que no leyese nada. Datos del
cuadro~\ref{cuad:progresion}.}
\label{fig:progresion}
\end{figure}

La lectura del cuadro~\ref{cuad:progresion} es la tesis metodológica del trabajo:
\textbf{cada vez que se le dio al modelo una representación más rica del texto, el
desempeño subió; y cuando solo se aumentó la potencia del clasificador sobre la misma
representación pobre, no subió}. XGBoost sobre tres variables es marginalmente peor que
una regresión logística sobre las mismas tres variables.

El modelo afinado directamente sobre el texto es un \textbf{\cod{roberta-base}} de 125
millones de parámetros con cabeza de regresión sobre \cod{BT\_easiness}, y alcanza además
\textbf{RMSE 0,560} y \textbf{correlación de Pearson 0,89} con la etiqueta continua real.

\textbf{Y aquí hay que evitar una confusión que el documento invitaba a hacer.} Ese
\cod{roberta-base} \textbf{no es el modelo que se despliega}. El que va en la aplicación es
un \cod{distil\allowbreak roberta-base} —la mitad de capas, 82 millones de parámetros— y se entrenó
sobre otro corpus, el de anclaje MCER, no sobre \cod{BT\_easiness}. Las cifras de este
apartado son del experimento sobre CLEAR; las del modelo desplegado están en
el~\S\ref{sec:validacion} y no se pueden encadenar con estas. La diferencia de calidad
entre las dos arquitecturas sobre CLEAR \textbf{no se midió}, así que la decisión de usar
la destilada se tomó por tamaño y velocidad, no porque se comprobara que costaba poco.

\textbf{Sobreajuste, visto y controlado.} El entrenamiento corrió seis épocas con tasa de
aprendizaje 2e-5, decaimiento de pesos 0,01, calentamiento y precisión mixta. El RMSE de
validación mejoró hasta la época 3 (0,560) y a partir de ahí empeoró de forma monótona
—época 4: 0,578; época 5: 0,637; época 6: 0,660— mientras la pérdida de entrenamiento
seguía bajando (figura~\ref{fig:epocas}). Es un caso de libro. Como se guardaba un punto
de control por época y se cargaba automáticamente el mejor al terminar, \textbf{el modelo
reportado es el de la época 3, no el de la última}.

% ------------------------------------------------------ FIGURA: épocas
\begin{figure}[htbp]
\centering
\begin{tikzpicture}
\begin{axis}[
  ejell,
  width=0.72\linewidth,
  height=5.6cm,
  xlabel={Época de entrenamiento},
  ylabel={RMSE de validación},
  xmin=2.6, xmax=6.4,
  ymin=0.53, ymax=0.70,
  xtick={3,4,5,6},
  ymajorgrids=true,
  legend pos=north west,
]
\addplot[color=llrojo, mark=*, mark size=2.2pt, mark options={fill=llrojo}]
  coordinates {(3,0.560) (4,0.578) (5,0.637) (6,0.660)};
\addlegendentry{RMSE de validación}

\draw[llverde, line width=0.8pt, dashed]
  (axis cs:2.6,0.560) -- (axis cs:6.4,0.560);
\node[font=\scriptsize, color=llverde!65!black, anchor=west]
  at (axis cs:4.4,0.548) {mejor punto de control: época 3};

\node[draw=llverde!70, fill=llverde!10, rounded corners=2pt, inner sep=3pt,
      font=\scriptsize, anchor=south west] at (axis cs:3.05,0.575)
  {modelo desplegado};
\draw[-{Stealth[length=2mm]}, llverde!80] (axis cs:3.12,0.574) -- (axis cs:3.02,0.562);
\end{axis}
\end{tikzpicture}
\caption[Sobreajuste visto y controlado]{El sobreajuste, medido. El RMSE de validación
toca su mínimo en la época 3 y empeora de forma monótona después, mientras la pérdida de
entrenamiento seguía bajando. Se guardó un punto de control por época y se cargó
automáticamente el mejor: el modelo desplegado es el de la época 3, no el de la sexta.}
\label{fig:epocas}
\end{figure}

\subsection{Dónde funciona y dónde no: análisis por dominio}
\label{sec:dominio}

Un modelo que se evalúa solo en agregado esconde sus propios puntos ciegos. Se midió la
correlación con la dificultad real por subconjunto, comparando dos variantes del modelo
entrenadas sobre corpus distintos (v1 sobre \cod{BT\_easiness}, v2 afinado sobre material
graduado para aprendices).

\begin{table}[htbp]
\centering
\small
\begin{tabular}{@{}l R{1.6cm} R{2.0cm} R{2.0cm}@{}}
\toprule
\textbf{Dominio} & \textbf{n} & \textbf{v1} & \textbf{v2} \\
\midrule
Literatura (CLEAR, \cod{Lit}) & 461 & \textbf{0,848} & 0,760 \\
Informativo (CLEAR, \cod{Info}) & 489 & \textbf{0,905} & 0,837 \\
Banda Lexile 500 (más fácil) & 62 & \textbf{0,651} & 0,459 \\
Banda Lexile 1300 (más difícil) & 272 & \textbf{0,877} & 0,773 \\
Periodístico (OneStopEnglish) & 567 & 0,458 & \emph{no comparable} \\
Exámenes MCER (British Council) & 59 & $-0{,}239$ & \textbf{0,175} \\
\bottomrule
\end{tabular}
\caption{Matriz modelo $\times$ dominio. Correlación orientada con la dificultad real;
más alto es mejor.}
\label{cuad:dominio}
\end{table}

% --------------------------------------------------- FIGURA: modelo x dominio
\begin{figure}[htbp]
\centering
\begin{tikzpicture}
\begin{axis}[
  barrall,
  width=0.86\linewidth,
  height=6.8cm,
  xbar,
  bar width=6.5pt,
  xmin=-0.35, xmax=1.05,
  xlabel={Correlación orientada con la dificultad real},
  symbolic y coords={examen,prensa,lex1300,lex500,info,lit},
  ytick={lit,info,lex500,lex1300,prensa,examen},
  yticklabels={Literatura (n=461),Informativo (n=489),Lexile 500 (n=62),Lexile 1300 (n=272),Periodístico (n=567),Exámenes MCER (n=59)},
  yticklabel style={font=\scriptsize},
  enlarge y limits=0.11,
  legend style={at={(0.5,-0.22)}, anchor=north, legend columns=2},
  every node near coord/.append style={font=\scriptsize, color=llgris},
  nodes near coords={\pgfmathprintnumber[fixed,precision=3]{\pgfplotspointmeta}},
  nodes near coords align={horizontal},
]
\addplot[fill=llcian!35, draw=llcian!85] coordinates {
  (0.848,lit) (0.905,info) (0.651,lex500)
  (0.877,lex1300) (0.458,prensa) (-0.239,examen)};
\addlegendentry{v1 (\cod{BT\_easiness})}
\addplot[fill=llverde!35, draw=llverde!85] coordinates {
  (0.760,lit) (0.837,info) (0.459,lex500)
  (0.773,lex1300) (0.175,examen)};
\addlegendentry{v2 (material graduado)}
\draw[llgris!70, line width=0.6pt] (axis cs:0,lit) -- (axis cs:0,examen);
\end{axis}
\end{tikzpicture}
\caption[Matriz modelo $\times$ dominio]{v1 domina en las cuatro filas de CLEAR, donde la
etiqueta son juicios de lectores nativos sobre prosa continua. \textbf{A v2 le falta la
barra de la fila periodística a propósito}: la cifra que ocupaba ese hueco estaba medida
sobre los 567 textos de OneStopEnglish, y v2 se afinó con 151 de sus 189 artículos, así
que el 80~\pct{} de ese conjunto era su propio entrenamiento. Lo detectó una auditoría
externa y se retira en lugar de maquillarse. Queda por tanto sin respaldo la lectura
simétrica —que v2 gane donde la etiqueta es una graduación para aprendices—, porque la
única fila que la sostenía estaba contaminada. Datos del cuadro~\ref{cuad:dominio}.}
\label{fig:dominio}
\end{figure}

\textbf{Por qué la casilla de OneStopEnglish está vacía para v2, y no es un olvido.} La
cifra que ocupaba ese hueco —un 0,840 que hacía ganar a
v2 con holgura— estaba medida sobre los 567 textos completos, y \textbf{v2 se afinó
precisamente con 151 de los 189 artículos de ese corpus}. El 80~\pct{} de esos textos era
su propio entrenamiento, y se comparaban contra un v1 que no los había visto nunca. No es
que la cifra fuera falsa: es que no medía lo que la tabla decía medir. Se retira en lugar
de maquillarla.

Eso obliga a rebajar la conclusión que colgaba de ella. Lo que sigue en pie es que
\textbf{v1 gana en dificultad de lectura tipo CLEAR} —juicios de lectores nativos sobre
prosa continua—, con ventaja en las cuatro filas medidas limpiamente. Lo que \textbf{ya no
se sostiene con esta evidencia} es la afirmación simétrica de que v2 gane donde la
etiqueta es una graduación para aprendices: la única fila que la respaldaba estaba
contaminada, y la de British Council, donde v2 sí aventaja a v1, tiene una correlación de
0,175, demasiado baja para sostener nada por sí sola.

La decisión de qué escala presentar al usuario precede a la de qué modelo desplegar; se
eligió la escala anclada al MCER, y con ella el modelo correspondiente. Esa decisión no
cambia, pero conviene saber que se tomó con menos evidencia de la que este cuadro
aparentaba dar.

\subsection{El modelo desplegado y su validación}
\label{sec:validacion}

El modelo en producción, \textbf{v4}, es un \cod{distilroberta-base}
\parencite{sanh2019distilbert,liu2019roberta} de seis capas y 82 millones de parámetros,
con cabeza de regresión, servido con \emph{Transformers} \parencite{wolf2020transformers}
y cuyos pesos ocupan 328 MB en disco. Puntúa bloques de 250 palabras truncados a 256
\emph{tokens} y promedia sin ponderar, con un tope de \textbf{ocho ventanas por
documento} repartidas a lo largo del texto: es lo que mantiene el análisis en \textbf{alrededor de un segundo}
en una CPU —0,96 a 1,04 s medidos sobre ocho ventanas en el equipo del autor, un
Ryzen 5 2600—, y por eso un vistazo mide una muestra del texto y no el texto
entero. El precálculo del catálogo, en cambio, recorre el libro completo sin
muestrear y publica una dificultad por capítulo; esa remedición \textbf{terminó el 12 de
septiembre de 2026} (accionable A3), y las letras del catálogo son ya las medidas por el
v4. Su escala
interna es 0--5 —un valor por nivel
MCER, con posiciones intermedias (A1+, A2+, B1+, B2+)— y la pública es esa misma escala
multiplicada por 0,4. \textbf{El modelo es final y no se reentrena.}

\textbf{Antes de las cifras, las reglas del juego.} Por este apartado y por el siguiente
circulan \textbf{cinco conjuntos de textos distintos}, y confundirlos es fácil porque
cuatro de los cinco salen del mismo material de British Council. No son subconjuntos
ordenados unos de otros y sus tamaños no se restan entre sí. El
cuadro~\ref{cuad:conjuntos} dice qué es cada uno, qué se mide con él y qué \emph{no} se
puede hacer con él; el resto del apartado los cita por su nombre.

\begin{footnotesize}
\begin{longtable}{@{}L{1.9cm} r L{3.0cm} L{3.6cm} L{4.3cm}@{}}
\caption{Los cinco conjuntos de textos que aparecen en la validación. Cuatro de los cinco
salen del mismo catálogo de British Council, con extracciones y filtros distintos: no son
subconjuntos unos de otros y sus tamaños no se restan.}
\label{cuad:conjuntos}\\
\toprule
\textbf{Conjunto} & \textbf{n} & \textbf{De dónde sale} & \textbf{Qué se mide con él} &
\textbf{Qué \emph{no} se puede hacer con él} \\
\midrule
\endfirsthead
\toprule
\textbf{Conjunto} & \textbf{n} & \textbf{De dónde sale} & \textbf{Qué se mide con él} &
\textbf{Qué \emph{no} se puede hacer con él} \\
\midrule
\endhead
\bottomrule
\endfoot
\bottomrule
\endlastfoot
\textbf{Validación interna} & 144 & Partición de validación del corpus de anclaje
MCER: 654 documentos aceptados, 510 de entrenamiento y 144 reservados & Exactitud,
acierto dentro de $\pm 1$ nivel y Spearman 0,908 del modelo desplegado
(cuadro~\ref{cuad:desplegado}) & Presentarlo como validación externa. Sale del mismo
corpus con el que se afinaron los pesos, y sobre él se eligió el punto de control, así que
su 0,908 no es comparable con las cifras de British Council \\
\addlinespace
\textbf{British Council, tanda de Colab} & 38 & Extracción de los PDF hecha en Colab
durante el entrenamiento, con su propio filtro & Spearman 0,829 y el 100\pct{} dentro de
$\pm 1$ (cuadro~\ref{cuad:desplegado}); las cuatro medias por nivel que ajustan la recta
de compresión (figura~\ref{fig:compresion}) & Restarlo de los 47. Es otra tanda, con otra
extracción y otro filtro: la diferencia no son «nueve textos concretos» que se puedan
nombrar \\
\addlinespace
\textbf{British Council, reproducción local} & 47 & Los 59 publicados, menos 2 sin cuerpo
localizable al extraer el PDF y 10 rechazados por el filtro de prosa
(cuadro~\ref{cuad:balance}) & \textbf{El resultado principal}: el orden creciente por
nivel (cuadro~\ref{cuad:validacion}) y las tres cifras del \emph{baseline} ---0,860,
0,807 y 0,668--- del cuadro~\ref{cuad:baseline} & Leerlo como «los textos que la
aplicación aceptaría». El validador local filtra más flojo que la app
(§\ref{sec:filtro}) \\
\addlinespace
\textbf{Los que pasan el filtro completo} & 36 & Esos mismos 47 refiltrados conservando
los saltos de línea, es decir con las cuatro reglas del filtro de prosa continua
activas & Cuánto depende el conjunto evaluado de una regla que el validador desactiva:
sobreviven 36 y A1 se queda con un solo texto (§\ref{sec:filtro}) & Sostener el orden
sobre cinco niveles. No se usa para ninguna cifra publicada en este documento \\
\addlinespace
\textbf{Publicados por British Council} & 59 & Todos los PDF descargados del catálogo
gratuito. Los nombres de los 59 están en el manifiesto del corpus externo & El
diagnóstico del §\ref{sec:calidad}: que el corpus no es homogéneo, y el 89,8\pct{} de
coincidencia entre las dos vías con que se definió el filtro & Usarlo como corpus de
validación. Incluye menús, pósters y horarios, y las correlaciones medidas sobre él no se
comparan con las de después de filtrar \\
\end{longtable}
\end{footnotesize}

El corpus de anclaje MCER con el que se afinó es \textbf{ELG CEFR English} \parencite{breuker2023cefr} (licencia no
comercial; cuadro~\ref{cuad:fuentes}), una sola fuente: 712 documentos, de los que se aceptan 654 como prosa —57 son cortos
y 1 tiene forma de frase anómala—, repartidos por documento y con semilla fija en 510 de
entrenamiento y 144 de validación. Esos 144 producen las cifras de la primera fila del
cuadro~\ref{cuad:desplegado}.

\textbf{De dónde salen esas cifras, y cómo se comprueban.} Del guion de entrenamiento del
v4 y de su tarjeta de métricas, que desde el 12 de septiembre de 2026 están en el
repositorio (\cod{entrenamiento\_v4/}). Ese día se \textbf{recalcularon} con el guion, los
pesos publicados y los datos públicos de ELG: la misma partición y las cuatro métricas
idénticas a cuatro decimales (exactitud 0,5347, dentro de un nivel 0,9722, MAE 0,4985,
Spearman 0,9084). Las 144 parejas de nivel y predicción están en
\cod{entrenamiento\_v4/\allowbreak validacion\_144\_recalculada.csv}. Y un matiz que no se
esconde: \textbf{no es un conjunto de prueba virgen}, porque el entrenamiento se quedó con el
punto de control que mejor puntuaba en esos mismos 144 documentos.

\begin{table}[htbp]
\centering
\small
\begin{tabular}{@{}l R{1.4cm} R{1.8cm} R{3.0cm} R{2.0cm}@{}}
\toprule
\textbf{Conjunto} & \textbf{n} & \textbf{Exacto} & \textbf{Dentro de $\pm 1$ nivel} & \textbf{Spearman} \\
\midrule
Validación interna (ELG reservado) & 144 & 53,5\pct & \textbf{97,2\pct} & \textbf{0,908} \\
British Council (desarrollo) & 38 & 50,0\pct & \textbf{100\pct} & 0,829 \\
\bottomrule
\end{tabular}
\caption{Métricas del modelo desplegado. MAE de 0,499 niveles de la escala interna
0--5 —0,200 en la escala pública 0--2— en el conjunto de
validación.}
\label{cuad:desplegado}
\end{table}

\textbf{Un 97,2\pct{} suena a mucho, y conviene poner al lado con qué se compara.} Con seis
niveles, acertar «dentro de $\pm 1$» es más fácil de lo que parece: contestar
\textbf{al azar} acierta el 44,4\pct{} de las veces, y contestar \textbf{siempre B1}, sin
mirar el texto, el 50\pct{} —ambas cifras suponiendo el nivel verdadero repartido por igual
entre los seis—. El modelo llega al 97,2\pct.

\textbf{Esas dos referencias son teóricas y hay que leerlas como tales.} Suponen el nivel
verdadero repartido por igual entre los seis, y \textbf{la distribución real de los 144
documentos de validación no está publicada}. Si ese conjunto estuviera concentrado en los
niveles centrales, «contestar siempre B1» acertaría bastante más del 50\pct. Lo honesto es
decir que un 97,2\pct{} está muy por encima de lo que da el azar bajo una hipótesis
razonable, no que la ventaja empírica sobre un predictor constante esté medida: no lo
está, porque faltan los recuentos por nivel de ese conjunto.

\textbf{El resultado que de verdad sostiene el trabajo no es ninguno de esos porcentajes,
sino el orden.} Esto es lo que se midió: los \textbf{47 textos de la reproducción local de
British Council}, cada uno con un nivel MCER oficial asignado por la institución que lo
publica y ninguno usado para ajustar los pesos, puntuados con \cod{validar\_local.py} y
agrupados después por su nivel. Y esto es lo que salió: \textbf{las medias por nivel quedan
estrictamente crecientes en los cinco niveles disponibles}, de A1 a C1 —0,5040 · 0,5561 ·
0,9276 · 1,2414 · 1,4952 en la escala 0--2—, y la correlación de Spearman con el nivel
oficial es de \textbf{0,860} (cuadro~\ref{cuad:validacion}, figura~\ref{fig:validacion} y
cuadro~\ref{cuad:baseline}). Ese orden —no el acierto exacto— es el resultado que el
trabajo afirma haber conseguido, y es el que se cita en las conclusiones.

\begin{table}[htbp]
\centering
\small
\begin{tabular}{@{}l R{2.4cm} R{3.4cm} R{2.6cm}@{}}
\toprule
\textbf{Nivel oficial} & \textbf{n (local)} & \textbf{Media en escala 0--2} & \textbf{Desviación} \\
\midrule
A1 & 3 & 0,5040 & 0,111 \\
A2 & 8 & 0,5561 & 0,223 \\
B1 & 12 & 0,9276 & 0,198 \\
B2 & 12 & 1,2414 & 0,236 \\
C1 & 12 & 1,4952 & 0,172 \\
\bottomrule
\end{tabular}
\caption{Reproducción local de la validación externa, sobre los 47 textos que quedan tras
la extracción de los PDF \textbf{y el filtro de prosa continua}. Fuente:
\cod{datos\_figuras/\allowbreak validacion.dat}.}
\label{cuad:validacion}
\end{table}

% ---------------------------------------------------- FIGURA: validación
\begin{figure}[htbp]
\centering
\begin{tikzpicture}
\begin{axis}[
  ejell,
  width=0.78\linewidth,
  height=6.2cm,
  ymin=0, ymax=1.95,
  xmin=-0.7, xmax=4.7,
  xtick={0,1,2,3,4},
  xticklabels={A1,A2,B1,B2,C1},
  xlabel={Nivel MCER oficial del texto \emph{(fuente institucional)}},
  ylabel={Puntuación media en la escala 0--2},
  ymajorgrids=true,
]

\addplot[
  ybar, bar width=22pt,
  fill=llcian!35, draw=llcian!90,
  error bars/.cd, y dir=both, y explicit,
  error bar style={llgris, line width=0.7pt},
  error mark options={llgris, mark size=2.5pt},
] table[x expr=\coordindex, y=media, y error=desv] {datos_figuras/validacion.dat};

% Anclas teóricas de cada nivel (nivel entero x 0,4)
\addplot[only marks, mark=-, mark size=11pt, llgris!85, line width=1.1pt]
  coordinates {(0,0.0) (1,0.4) (2,0.8) (3,1.2) (4,1.6)};

\node[font=\scriptsize, color=llgris, anchor=south] at (axis cs:0,0.72) {n\,=\,3};
\node[font=\scriptsize, color=llgris, anchor=south] at (axis cs:1,0.90) {n\,=\,8};
\node[font=\scriptsize, color=llgris, anchor=south] at (axis cs:2,1.19) {n\,=\,12};
\node[font=\scriptsize, color=llgris, anchor=south] at (axis cs:3,1.55) {n\,=\,12};
\node[font=\scriptsize, color=llgris, anchor=south] at (axis cs:4,1.72) {n\,=\,12};

\node[font=\scriptsize\itshape, color=llgris, anchor=north west, align=left]
  at (axis cs:-0.55,1.90) {los trazos grises son el ancla teórica\\del nivel (nivel entero $\times\,0{,}4$)};

\end{axis}
\end{tikzpicture}
\caption[Validación externa por nivel MCER]{El resultado que sostiene el trabajo: las
medias por nivel quedan \textbf{estrictamente crecientes} sobre textos externos con nivel
MCER asignado por una institución de enseñanza. Las barras de error son la desviación
dentro de cada nivel; con $n$ de 3 a 12 textos, lo que se afirma es el orden global y no
cada comparación por pares. Datos de \cod{datos\_figuras/\allowbreak validacion.dat}
(cuadro~\ref{cuad:validacion}).}
\label{fig:validacion}
\end{figure}

\subsubsection*{La letra pequeña de ese resultado}

\textbf{Aquí empieza la letra pequeña, y conviene leerla antes de citar la cifra.} Son
cuatro salvedades. Ninguna anula la afirmación anterior, pero cada una acota lo que puede
hacerse con ella.

\textbf{Una: el escalón A1--A2 está dentro del ruido.} Es de 0,052 sobre desviaciones de
0,111 y 0,223, y descansa en tres textos. Los tres escalones siguientes, no. Lo que se
afirma es el orden global de las cinco medias, no cada comparación por pares: con $n$ de 3
a 12 textos por nivel, el primer escalón no se sostendría por sí solo.

\textbf{Dos: es una reproducción local, no una reproducción independiente.} El
cuadro~\ref{cuad:validacion} se \textbf{recalcula} con \cod{validar\_local.py} en
cualquier máquina que tenga los mismos PDF y los mismos pesos, pero los PDF de British
Council tienen derechos y no acompañan al repositorio: quien quiera rehacerla los descarga
con el navegador y los verifica por hash con \cod{tools/\allowbreak bajar\_corpus\_bc.py}. Lo que se reproduce exactamente es \textbf{el orden};
las medias difieren ligeramente de las obtenidas en Colab (C1 coincide hasta la tercera
cifra decimal —1,4952 frente a 1,4950—, A2 se desvía 0,09).

\textbf{Tres: los 38 de Colab y los 47 de aquí son dos tandas distintas, y no se restan.}
La diferencia tiene dos causas —otra extracción de los PDF y otro filtro—, así que
$47-38$ no son nueve textos que se puedan nombrar. Contado sobre los archivos reales, el
balance de esta tanda cierra así:

\begin{table}[htbp]
\centering\small
\begin{tabular}{@{}l r@{}}
\toprule
 & \textbf{Textos} \\
\midrule
Publicados por British Council & 59 \\
menos: sin cuerpo localizable al extraer el PDF & $-2$ \\
menos: \textbf{rechazados por \cod{scoring.continuous\_prose}} & $\mathbf{-10}$ \\
\midrule
\textbf{Evaluados} & \textbf{47} \\
\bottomrule
\end{tabular}
\caption{De 59 a 47, paso a paso. La primera resta es de extracción; la segunda
es el filtro de prosa continua.}
\label{cuad:balance}
\end{table}

\textbf{Y cuatro: el validador local filtra más flojo que la propia aplicación.} Con el
filtro completo de la app sobreviven \textbf{36 de esos 47} y A1 se queda con un solo
texto, de modo que el orden creciente sobre cinco niveles está medido sobre un conjunto
que la aplicación habría filtrado más. Por qué ocurre, qué regla se desactiva y qué
textos caen se explica entero en el §\ref{sec:filtro}, al final del
§\ref{sec:modelacion}.

\subsection{El \emph{baseline} propio, y qué aporta el modelo por encima de él}
\label{sec:baseline}

\begin{table}[htbp]
\centering
\small
\begin{tabular}{@{}L{7.6cm} R{5.4cm}@{}}
\toprule
\textbf{Señal} & \textbf{Spearman con el nivel oficial (47 textos)} \\
\midrule
Modelo v4 & \textbf{0,860} \\
Contar las palabras del texto, sin modelo & 0,807 \\
Modelo, descontando la longitud (correlación parcial) & \textbf{0,668} \\
\bottomrule
\end{tabular}
\caption{El modelo frente a la señal trivial. Las tres cifras las calcula \cod{validar\_local.py} sobre los mismos 47 textos, con Spearman sobre rangos y la correlación parcial en la misma escala.}
\label{cuad:baseline}
\end{table}

Un contador de palabras llega casi tan lejos como el modelo en ese corpus, y omitirlo
sería deshonesto. Lo que sostiene al modelo es la tercera fila: \textbf{aporta señal por
encima de la longitud}, que es precisamente lo que un contador de palabras no puede
hacer —y conviene decirlo con precisión: no es que su correlación parcial sea peor, es que
\textbf{no está definida}, porque descontar la longitud de la propia longitud deja un
residuo de varianza cero—. Además, el modelo no recibe la longitud del documento como
variable —puntúa
bloques de tamaño fijo y promedia—, aunque su salida correlacione 0,732 con ella, porque
los corpus de examen gradúan también la extensión de la tarea.

\textbf{La escala está comprimida, y la aplicación lo asume.} Contra el ancla real, el
ajuste ronda $\text{predicho} \approx 0{,}72 \times \text{real} + 0{,}38$
(figura~\ref{fig:compresion}). Conviene decir de dónde sale esa recta, porque no está
ajustada sobre los textos uno a uno: son \textbf{mínimos cuadrados sobre las cuatro medias
por nivel que reportó Colab} (A2 a C1; el catálogo gratuito de British Council no tiene C2
y aquella tanda no traía A1), y dan $0{,}7196x + 0{,}3760$. Rehecha sobre las cinco medias
locales de hoy sale $0{,}67x + 0{,}41$: cambian los decimales, no la conclusión, que es que
la pendiente está muy por debajo de 1. De modo que un texto A2 tiende a puntuar cerca de
B1. Por
eso la interfaz muestra \textbf{bandas} (básico $< 0{,}6 \leq$ intermedio $< 1{,}4 \leq$
avanzado) y la letra MCER solo como orientación, en lugar de fingir una precisión de dos
decimales que el modelo no tiene.

Los cortes de esas bandas salen del MCER, y esa es una corrección posterior. Al principio
estaban en 0,8 y 1,6, que además \textbf{ni siquiera partían el 0--2 en tres partes
iguales}: dejaban dos tramos de 0,8 y uno de 0,4. Pero como el modelo comprime
sus predicciones, con esos cortes un texto C1 caía en «intermedio» y «avanzado» solo lo
alcanzaba un C2: tres de los cinco niveles medidos compartían banda. Ahora la banda se
deriva del ancla MCER, con el agrupamiento del propio marco —usuario básico A1--A2,
independiente B1--B2, competente C1--C2—, y las fronteras quedan en \textbf{0,6 y 1,4}.
Así la palabra y la letra no pueden contradecirse por construcción. Ninguna puntuación
cambió: solo cambió cómo se agrupan.

\textbf{Un intento de ampliar la validación que hubo que retirar, y por qué se deja
escrito.} El 12 de septiembre de 2026 se midió el v4 sobre los dos conjuntos ingleses de
referencia de UniversalCEFR \parencite{imperial2025universalcefr}, \cod{cambridge\_exams\_en} y
\cod{elg\_cefr\_en}: 739 textos, Spearman \textbf{0,946}. Era más alto que la validación
interna (0,908), y eso bastó para no creérselo. Leyendo el guion de entrenamiento del v4
salió la causa: \textbf{el v4 se entrenó con ELG CEFR English}, que es exactamente
\cod{elg\_cefr\_en} (654 aceptados + 57 cortos + 1 descartado = 712, el tamaño del
conjunto). Y los 325 textos distintos de Cambridge están \textbf{todos} dentro de él: 130
idénticos y 195 casi idénticos, con una mediana del 92\pct{} de sus 8-gramas compartidos y
el mismo nivel en los 195. Aquel 0,946 medía el modelo sobre su propio entrenamiento. No
hay en UniversalCEFR un subconjunto inglés limpio para el v4, y la validación de este
apartado sigue siendo la del cuadro~\ref{cuad:validacion}.

\textbf{Una validación externa de verdad: OneStopEnglish.} Con la composición del anclaje ya
conocida se pudo buscar un corpus que los pesos del v4 no hubieran visto, y comprobarlo en vez
de suponerlo: ninguno de los textos de OneStopEnglish \parencite{vajjala2018onestop} comparte
con ELG CEFR English ni una secuencia de ocho palabras. Son 189 artículos de prensa escritos
tres veces, en nivel Elementary, Intermediate y Advanced. Se mide una puntuación por artículo
y nivel ---560 unidades--- con \cod{validar\_onestop.py}, que repite la comprobación de fuga
cada vez que se ejecuta (cuadro~\ref{cuad:onestop}).

\begin{table}[htbp]
\centering
\small
\begin{tabular}{@{}lrr@{\hspace{2em}}lrr@{}}
\toprule
\textbf{Nivel} & \textbf{n} & \textbf{Media v4} & \textbf{Señal} & \textbf{Modelo} & \textbf{Palabras} \\
\midrule
Elementary & 185 & 0,934 & Spearman con el nivel & \textbf{0,800} & 0,692 \\
Intermediate & 189 & 1,272 & Parcial, sin la longitud & \textbf{0,688} & --- \\
Advanced & 186 & 1,515 & Artículos en orden (de 185) & \textbf{183} & 179 \\
\bottomrule
\end{tabular}
\caption{Validación externa con OneStopEnglish, un corpus que los pesos del v4 no vieron. La
media va en la escala pública 0--2; «Palabras» es contar palabras, sin modelo.}
\label{cuad:onestop}
\end{table}

Las medias suben en los tres niveles y el modelo conserva 0,688 descontando la longitud: aporta
señal propia en un corpus que no vio. \textbf{Pero la última fila hay que leerla con cuidado},
porque es tentadora: dentro de cada artículo el modelo ordena bien las tres versiones en 183 de
185, y \textbf{la longitud sola ya acierta 179}, porque las versiones avanzadas son más largas
(523, 667 y 809 palabras de media). Donde la longitud falla ---seis artículos--- el modelo
acierta en cinco. Es una ventaja real, y es pequeña. Dos matices más: la escala vuelve a salir
comprimida y desplazada hacia arriba ---un texto Elementary puntúa 0,93, en la banda
intermedia---, y OneStopEnglish \textbf{sí se usó para desarrollar los modelos anteriores} (el
v2 se afinó con él), así que es externo para los pesos del v4, no virgen para el proyecto.

\textbf{Lectura graduada por sus editores.} Tres bibliotecas infantiles abiertas ponen a cada
libro un nivel propio: StoryWeaver del 1 al 4, African Storybook y Let's Read del 1 al 5.
\textbf{No son niveles del MCER}, así que solo se puede validar el orden. Se midieron con
\cod{validar\_graduados.py}, con la misma comprobación de fuga contra ELG CEFR English (0 textos
contaminados en las tres) y dos decisiones que hay que declarar: los cuentos de menos de 250
palabras se unen en bloques de al menos 250 \emph{sin mezclar fuente ni nivel}, y la longitud
que se descuenta es la natural de los cuentos, no la del bloque. Los intervalos son bootstrap al
95\pct{} con 2.000 remuestreos (cuadro~\ref{cuad:graduados}).

\begin{table}[htbp]
\centering
\small
\resizebox{\linewidth}{!}{%
\begin{tabular}{@{}lrlrll@{}}
\toprule
\textbf{Fuente} & \textbf{Unid.} & \textbf{Spearman v4} & \textbf{Longitud} & \textbf{Parcial sin longitud} & \textbf{Medias crecientes} \\
\midrule
StoryWeaver (1--4) & 467 & \textbf{0,426} (0,348 a 0,502) & 0,542 & \textbf{0,206} (0,111 a 0,295) & 99,6\pct{} \\
African Storybook (1--5) & 174 & 0,416 (0,279 a 0,537) & 0,810 & 0,133 ($-$0,029 a 0,289) & 54\pct{} \\
Let's Read (1--5) & 1.000 & 0,178 (0,118 a 0,235) & 0,620 & $-$0,083 ($-$0,144 a $-$0,026) & 1,8\pct{} \\
\bottomrule
\end{tabular}}
\caption{Validación con lectura graduada por sus editores. «Longitud» es el Spearman de la
longitud natural con el nivel; «Medias crecientes», la proporción de remuestreos en que las
medias por nivel salen en orden. Entre paréntesis, el intervalo bootstrap al 95 por ciento. StoryWeaver es una muestra aleatoria de 649 libros intercalada
por niveles.}
\label{cuad:graduados}
\end{table}

Las tres dicen cosas distintas, y conviene no promediarlas. \textbf{StoryWeaver valida}: los
cuatro niveles salen separados y en orden, y el modelo aporta orden \emph{además} de la longitud
(la parcial excluye el cero). \textbf{En African Storybook la longitud predice el nivel mucho mejor
que el modelo}, y eso describe al corpus: sus niveles son casi una escala de longitud. El v4 no
ve la longitud ---puntúa ventanas fijas de 250 palabras--- y dentro de cada nivel su correlación
con ella es prácticamente nula; ordena por contenido, pero no se puede afirmar que aporte orden
más allá de la longitud, y el orden entre los niveles 1 y 2 no está establecido. \textbf{Let's Read
no valida, y no por el modelo}: más de la mitad de sus libros son traducciones al inglés que
heredan el nivel del original en laosiano, jemer o indonesio; en esas 459 traducciones la
correlación es 0,023, y en los 360 originales, 0,207. Su etiqueta describe el libro en su idioma.

Esta sección se corrigió tras una revisión metodológica independiente: la primera redacción
leía la parcial de African Storybook como «lo explica sobre todo la longitud», daba sus medias
por «estrictamente crecientes» sin intervalos y no advertía que el filtro de prosa rechazaba 42
de sus 65 cuentos de nivel 1, sesgando ese nivel hacia arriba. Los detalles están en
\cod{auditoria/\allowbreak validacion\_graduados\_2026-09-12.md}.

% ---------------------------------------------------- FIGURA: compresión
\begin{figure}[htbp]
\centering
\begin{tikzpicture}
\begin{axis}[
  ejell,
  width=0.70\linewidth,
  height=7.0cm,
  xlabel={Ancla del nivel oficial en la escala 0--2},
  ylabel={Media predicha por el modelo},
  xmin=-0.12, xmax=1.78,
  ymin=-0.12, ymax=1.78,
  xtick={0,0.4,0.8,1.2,1.6},
  xticklabels={A1\\{\scriptsize 0,0},A2\\{\scriptsize 0,4},B1\\{\scriptsize 0,8},B2\\{\scriptsize 1,2},C1\\{\scriptsize 1,6}},
  xticklabel style={align=center},
  ytick={0,0.4,0.8,1.2,1.6},
  grid=both,
  ymajorgrids=true, xmajorgrids=true,
  legend style={at={(0.03,0.97)}, anchor=north west},
]

% identidad
\addplot[llgris!70, dashed, line width=0.8pt, domain=-0.1:1.75, samples=2] {x};
\addlegendentry{calibración perfecta ($y=x$)}

% recta ajustada documentada
\addplot[llrojo, line width=1.0pt, domain=-0.1:1.75, samples=2] {0.72*x+0.38};
\addlegendentry{$y = 0{,}72\,x + 0{,}38$}

% puntos observados
\addplot[only marks, mark=*, mark size=2.6pt, mark options={fill=llcian, draw=llazul}]
  coordinates {(0.0,0.5040) (0.4,0.5561) (0.8,0.9276) (1.2,1.2414) (1.6,1.4952)};
\addlegendentry{media medida por nivel}

\node[font=\scriptsize, color=llazul, anchor=north west] at (axis cs:0.02,0.50) {A1};
\node[font=\scriptsize, color=llazul, anchor=north west] at (axis cs:0.42,0.545) {A2};
\node[font=\scriptsize, color=llazul, anchor=north west] at (axis cs:0.82,0.925) {B1};
\node[font=\scriptsize, color=llazul, anchor=north west] at (axis cs:1.22,1.245) {B2};
\node[font=\scriptsize, color=llazul, anchor=north west] at (axis cs:1.62,1.492) {C1};

\draw[{Stealth[length=2mm]}-{Stealth[length=2mm]}, llambar!85, line width=0.8pt]
  (axis cs:0.4,0.40) -- (axis cs:0.4,0.556);
\node[font=\scriptsize, color=llambar!70!black, anchor=west]
  at (axis cs:0.44,0.30) {los niveles bajos se sobreestiman};
\draw[{Stealth[length=2mm]}-{Stealth[length=2mm]}, llambar!85, line width=0.8pt]
  (axis cs:1.6,1.60) -- (axis cs:1.6,1.4952);
\node[font=\scriptsize, color=llambar!70!black, anchor=east]
  at (axis cs:1.55,1.66) {los altos se subestiman};

\end{axis}
\end{tikzpicture}
\caption[La escala comprimida]{Por qué la interfaz muestra bandas y no decimales. Los
puntos son las medias medidas por nivel (cuadro~\ref{cuad:validacion}) frente al ancla
teórica de ese nivel. El orden se conserva —que es el objetivo—, pero la pendiente es
claramente menor que $1$: los niveles bajos se sobreestiman y los altos se subestiman.
La recta dibujada es $0{,}72x + 0{,}38$, ajustada sobre las \textbf{cuatro medias que
reportó Colab} (A2 a C1); ajustada sobre los \textbf{cinco puntos locales} que se ven
aquí sale $0{,}67x + 0{,}41$. Se dibuja la primera porque es la que cita el texto, y las
dos dicen lo mismo: la pendiente está muy por debajo de 1. Calibrar
esa recta afín sin tocar los pesos del modelo es el accionable A7
(cuadro~\ref{cuad:accionables}).}
\label{fig:compresion}
\end{figure}

\subsection{Cuadro de decisiones metodológicas}

Incluye las hipótesis que se probaron y se rechazaron (cuadro~\ref{cuad:decisiones}).
Reportarlas es parte del resultado.

\begin{small}
\begin{longtable}{@{}L{3.7cm} L{4.0cm} L{6.0cm}@{}}
\toprule
\textbf{Decisión} & \textbf{Razón} & \textbf{Evidencia / consecuencia} \\
\midrule
\endfirsthead
\toprule
\textbf{Decisión} & \textbf{Razón} & \textbf{Evidencia / consecuencia} \\
\midrule
\endhead
\bottomrule
\endfoot
\bottomrule
\caption{Decisiones metodológicas, incluidas las hipótesis rechazadas.}
\label{cuad:decisiones}
\endlastfoot
Regresión, no clasificación & La dificultad es continua; clasificar tira información ordinal & El error ordinal medio baja de 1,39 a 0,67 en la progresión completa \\
\addlinespace
Partición agrupada por libro & Evitar fuga de datos entre extractos \textbf{del mismo libro}. No separa autores: dos obras del mismo autor pueden caer a ambos lados, y eso no se comprobó & Aplicada a \emph{todos} los modelos comparados, no solo al final \\
\addlinespace
\cod{distilroberta} y no \cod{roberta-base} completo & La aplicación debe correr en CPU en el equipo de un estudiante & 328 MB y 0,9 s de carga, frente a 1,5$\times$ más parámetros (125 M contra 82 M). \textbf{La diferencia de calidad entre las dos no se midió}: la decisión se tomó por tamaño y velocidad, no porque se comprobara que costaba poco \\
\addlinespace
Promediar bloques sin ponderar, con la \textbf{ventana} como unidad & Simplicidad y reproducibilidad; ponderar por longitud reintroduce la longitud como señal & Constante del entrenamiento. La app y el precálculo del catálogo usan ahora el mismo estimador: durante un tiempo la app promedió las medias de cada capítulo y el catálogo todas las ventanas, y con capítulos de tamaño desigual \textbf{no son el mismo número} (0,004 de diferencia máxima sobre cinco libros; la medición se hizo en su momento y \textbf{no queda guion que la reproduzca}, así que se da como orden de magnitud, no como cifra auditable) \\
\addlinespace
\textbf{Rechazada: XGBoost sobre variables clásicas} & Se esperaba que un ensamble superase al modelo lineal & No lo hace: 0,28 frente a 0,29, y error ordinal peor (1,42 \emph{vs.} 1,39). El límite estaba en las variables \\
\addlinespace
\textbf{Rechazada: modelo híbrido \emph{embedding} + variables} & Se esperaba complementariedad & Ganancia nula (0,37 y 0,37; 0,91 \emph{vs.} 0,93). El \emph{embedding} ya contiene esa información \\
\addlinespace
\textbf{Rechazada: escalar la lista léxica al SPGC completo por rendimiento} & Se esperaba mejorar \cod{rare\_word\_ratio} & La correlación pasa de $-0{,}330$ a $-0{,}326$. Se conserva igualmente por cobertura, no por métrica \\
\addlinespace
\textbf{Rechazada: cuantización dinámica int8} & $+18\pct$ de velocidad en CPU & MAE de 0,210 sobre una escala de 0 a 2 y sesgo sistemático de $-0{,}21$. Un desplazamiento así cambia la letra MCER en el \textbf{52\pct} de la escala y la banda en el \textbf{21\pct}: no degrada la cifra, le cambia la etiqueta al libro. Inaceptable \\
\addlinespace
\textbf{Rechazada: precisión media (fp16) en la GPU local} & Se esperaba acelerar & Es \textbf{4,8$\times$ más lenta} que fp32 en una GTX 1650 SUPER: el chip TU116 no tiene \emph{tensor cores}. En una A100 sí compensa, y allí la desviación observada fue de unos 0,00037 de media; \textbf{tampoco quedó guardada}: mismo caso que la fila de arriba \\
\addlinespace
\textbf{Rechazada: \cod{inference\_mode} en lugar de \cod{no\_grad}} & Se esperaba menos sobrecarga & 0,4\pct{} de diferencia, por debajo del ruido \\
\addlinespace
\textbf{Rechazado: esfuerzo de razonamiento «mínimo» en la explicación léxica} & Es el doble de rápido (7,4 s frente a 14,9 s) & Con él, «had peeped» se explicó como pasado simple cuando es pasado perfecto. La gramática es justo lo que se evalúa \\
\addlinespace
Filtrar prosa continua para la validación externa & Excluir textos porque el modelo falla en ellos invalida la validación & El criterio salió del diagnóstico, no es \emph{a priori}, y así se dice en §\ref{sec:calidad}. Lo que lo sostiene es que se definió por \textbf{dos vías independientes} —título y estructura— que coinciden en el 89,8\pct{} sin volver a mirar las puntuaciones \\
\end{longtable}
\end{small}

\subsection{Qué filtra el validador local y qué filtra la aplicación}
\label{sec:filtro}

Este apartado desarrolla la cuarta salvedad del §\ref{sec:validacion}: el guion que
reproduce el cuadro~\ref{cuad:validacion} no aplica el mismo filtro que la aplicación, y
conviene ponerlo aquí y no esconderlo en el código.

\textbf{Qué regla se desactiva, y por qué.} El guion de reproducción une las líneas del
PDF con un espacio, y eso \textbf{desactiva una de las cuatro reglas} de
\cod{scoring.continuous\_prose}, la del maquetado de lista: sus dos comprobaciones —más del
17,5\% de las palabras en líneas de maquetado, o más de un tercio de líneas muy
cortas— necesitan saltos de línea para ver nada. Es decir, este validador filtra \textbf{más flojo que la propia aplicación}, que
sí recibe saltos de párrafo. Medido cambiando una cosa y solo una: con saltos de línea
sobreviven \textbf{36 textos en vez de 47}, y A1 se queda con uno solo, con lo que la
comprobación de orden creciente ya no se puede hacer sobre cinco niveles.

\textbf{Por qué se conserva así.} Porque la unión con espacio es la que reproduce el
cuadro~\ref{cuad:validacion} de este documento, que es la reproducción local. \textbf{No
reproduce las cifras de Colab}, y nada aquí lo hace: aquella tanda son 38 textos y A2 se
desvía 0,09. Pero el guion imprime las dos cuentas en cada ejecución, así que la
diferencia queda a la vista de quien lo ejecute.

\textbf{Qué textos se caen, y por qué regla.} En el cuadro~\ref{cuad:balance}, la primera
resta es de extracción. La segunda no: son \textbf{diez textos —nueve de A1 y uno de B1—
que rechaza \cod{scoring.continuous\_prose}}, y conviene decir con qué regla, porque el
filtro tiene cuatro. Los diez caen por la misma: \textbf{se quedan por debajo de las 150
palabras} una vez recortados entre «Reading text» y «Tasks». Son pósters, menús, horarios
y tarjetas, es decir, textos que en el original ya son cortos; el detector de maquetado de
lista ni siquiera llega a dispararse, porque el de longitud actúa antes. Los dos entornos
filtraron con reglas distintas, no solo extrajeron distinto.

\textbf{Lo que esto le hace al resultado principal.} Los tres textos A1 que sostienen el
escalón más débil del cuadro~\ref{cuad:validacion} \textbf{no pasan el filtro de la
aplicación}: pasan el de este guion, que es el más flojo de los dos. Con el filtro
completo, dos de los tres caen por maquetado de lista. Dicho de otro modo: una parte de
esos 47 son textos que la app rechazaría por parecer listas, y el escalón A1 es justo
donde más pesa. Los nombres de los 59 archivos están en
\cod{entrega/manifiesto\_corpus\_externo.md}, y \cod{validar\_local.py} imprime los diez
rechazados con su motivo.

% ===========================================================================
\section{Resolución del problema}
\label{sec:resolucion}
% ===========================================================================

\subsection{De la métrica a la decisión del lector}

Una correlación de 0,908 no le sirve de nada a un adolescente. Lo que le sirve son tres
decisiones, y el producto está construido para producir exactamente esas tres
(cuadro~\ref{cuad:decisiones_lector}).

\begin{table}[htbp]
\centering
\small
\begin{tabularx}{\textwidth}{@{}L{3.6cm} Y L{3.2cm}@{}}
\toprule
\textbf{Pregunta del lector} & \textbf{Qué la responde} & \textbf{Coste para él} \\
\midrule
¿Cuál es mi nivel? & Seis pasos, uno por pantalla: tres preguntas sobre hábitos observables (con qué frecuencia lee, si ha terminado un libro entero, qué puede contar de una página que acaba de leer) y tres fragmentos reales en escalera; o bien subir un libro que ya lee cómodo & Menos de dos minutos \\
\addlinespace
¿Qué leo ahora? & Búsqueda por significado en español sobre el catálogo; el nivel del perfil ordena lo más parecido y lo de otros niveles sale aparte, sin filtrar & Segundos \\
\addlinespace
¿Por dónde se va a complicar? & Curva de dificultad capítulo a capítulo, medida sobre el libro entero sin muestrear (figura~\ref{fig:curva_wp}) & De 12,7 s (vistazo, medido sobre \emph{Crime and Punishment}) a minutos (libro completo) \\
\addlinespace
¿Qué significa esta palabra \emph{aquí}? & Mantener presionada la palabra (con ratón, doble clic): lema, IPA británica y americana cuando el diccionario la trae (el 58\pct{} de las consultas sobre un libro real), definición de Wiktionary, el sentido concreto en esa frase, gramática y audio & Unos 0,7~s en caliente y unos 5~s la primera consulta de la sesión; la explicación contextual tarda más \\
\addlinespace
¿Lo estoy pronunciando bien? & Lectura en voz alta marcada palabra por palabra, con el sonido concreto que falla & Una frase \\
\bottomrule
\end{tabularx}
\caption{Del modelo a las decisiones del usuario.}
\label{cuad:decisiones_lector}
\end{table}

\subsection{Un recorrido completo, con números reales}
\label{sec:recorrido}

Tomemos el caso concreto de una lectora cuyo perfil da 1,20 sobre 2 (banda intermedia,
ancla B2). La figura~\ref{fig:recorrido} resume el recorrido completo, empezando por el
cuestionario de perfil; los seis pasos que siguen arrancan ya con ese perfil hecho
y detallan el resto.

% ---------------------------------------------------- FIGURA: recorrido
\begin{figure}[htbp]
\centering
\begin{tikzpicture}[node distance=6mm and 4mm, font=\footnotesize]

  \node[cajall, text width=22mm] (perfil) {\textbf{1. Perfil}\\3 preguntas\\+ 3 fragmentos};
  \node[cajall, text width=22mm, right=of perfil] (buscar) {\textbf{2. Buscar}\\búsqueda semántica\\en español};
  \node[cajall, text width=22mm, right=of buscar] (abrir) {\textbf{3. Abrir}\\descarga de\\Project Gutenberg};
  \node[cajall, text width=22mm, right=of abrir] (ficha) {\textbf{4. Ficha}\\todo contado\\sobre el fichero};
  \node[cajall, text width=22mm, right=of ficha] (curva) {\textbf{5. Curva}\\dificultad\\por capítulo};

  \node[cajall, text width=22mm, below=13mm of buscar] (palabra) {\textbf{6. Palabra}\\lema, IPA,\\sentido en la frase};
  \node[cajall, text width=22mm, right=of palabra] (voz) {\textbf{7. Voz}\\corrección\\palabra a palabra};

  \node[etiqll, below=1.5mm of perfil] {$<$ 2 min};
  \node[etiqll, below=1.5mm of buscar] {segundos};
  \node[etiqll, below=1.5mm of abrir] {un clic};
  \node[etiqll, below=1.5mm of ficha] {instantáneo};
  \node[etiqll, below=1.5mm of curva] {12,7 s a minutos};

  \node[etiqll, below=1.5mm of palabra] {0,7 s en pantalla};
  \node[etiqll, below=1.5mm of voz] {una frase};

  \draw[flechall] (perfil) -- (buscar);
  \draw[flechall] (buscar) -- (abrir);
  \draw[flechall] (abrir) -- (ficha);
  \draw[flechall] (ficha) -- (curva);
  \draw[flechall] (abrir.south) to[out=-90,in=90] (palabra.north);
  \draw[flechall] (palabra) -- (voz);

  \begin{scope}[on background layer]
    \node[draw=llverde!45, fill=llverde!7, rounded corners=5pt, inner sep=7pt,
          fit=(palabra)(voz), label={[font=\scriptsize\itshape, color=llverde!55!black]below:mientras lee}] {};
    \node[draw=llcian!45, fill=llcian!6, rounded corners=5pt, inner sep=7pt,
          fit=(perfil)(curva), label={[font=\scriptsize\itshape, color=llcian!55!black]above:antes de leer}] {};
  \end{scope}

\end{tikzpicture}
\caption[El recorrido del lector]{El recorrido completo, con los tiempos medidos de cada
paso. No es un modelo con una interfaz encima: es una secuencia de decisiones que el
lector no podía tomar antes. Lo único lento es la explicación contextual, que sale a la
red; por eso la ficha se pinta en dos etapas.}
\label{fig:recorrido}
\end{figure}

\begin{enumerate}[leftmargin=*, itemsep=4pt]
  \item Busca. Escribe \emph{«Crimen y castigo de Dostoyevski»}. Con su perfil de 1,20
        la pantalla avisa de que no hay libros en su nivel ni un poco por encima, y
        \emph{Crime and Punishment} sale el primero en «De otros niveles», sin pulsar
        nada: a 390~px queda a 679~px del principio de la página (ensayo del 15 de
        septiembre de 2026 sobre los reinicios 28 y 30b). Escribir solo
        \emph{«dostoyevski»} con un perfil puesto es peor: «En tu nivel» se llena de
        cuentos sin relación y el libro cae a 3.789~px, que es exactamente la
        limitación~15. El buscador no busca por
        título sino por significado, sobre \emph{embeddings} multilingües. Podría haber
        escrito en español lo que le apetece —\emph{«una novela de detectives
        victoriana»}— y el buscador respondería igual. Medido el 11 de septiembre de
        2026 con su perfil de 1,20 y el ajuste al nivel puesto (entonces una casilla;
        desde el 15 de septiembre \emph{Encontrar} ordena siempre por nivel y aparta en
        otro bloque lo que queda lejos de él), esa frase en español devuelve \emph{The Secret Adversary} (Christie, B2), un cuento infantil
        A1 de StoryWeaver, \emph{Great Expectations}, \emph{North and South},
        \emph{Bleak House} y, en séptimo lugar, \emph{The Memoirs of Sherlock Holmes}.
        Que el segundo puesto sea un A1 para una lectora de 1,20 no es un fallo del
        buscador: es el libro cuyo texto más se parece a la palabra «detective», la
        ajuste \textbf{acerca y no filtra}, y su nivel va escrito al lado.

        \textbf{Y aquí conviene decir algo que el propio recorrido enseña.} Lo que
        devuelve la búsqueda \textbf{no está filtrado por el nivel del lector}: ordena
        por cercanía al tema, y la dificultad va al lado de cada resultado para que el
        lector decida. En este recorrido eso importa, porque el libro que abre mide 1,70
        y ella está en 1,20: le queda por encima, y el valor de la herramienta es
        precisamente que \textbf{lo sabe antes de la primera página}.
  \item Abre el libro. Se descarga de Project Gutenberg y se abre dentro de la
        aplicación, sin salir a buscarlo.
  \item Aparece la ficha, \textbf{contada sobre el texto y no sobre metadatos},
        aunque no toda sobre el libro entero: las palabras y el tiempo de lectura
        salen del fichero completo, y la longitud de frase, el diálogo y el
        vocabulario se calculan sobre una muestra con \textbf{tope de 60.000 palabras},
        repartido a partes iguales entre los capítulos —la variedad léxica, sobre las
        primeras 20.000 de esa muestra—, porque recorrer \emph{Guerra y
        paz} entero para pintar una ficha haría esperar al lector. Lo que trae:
        palabras reales, tiempo de lectura estimado a 180 o 110 palabras por minuto según
        el texto le quede por debajo o por encima de su nivel —no es su ritmo medido, que
        la aplicación no conoce—, longitud media de frase,
        proporción de diálogo, variedad de vocabulario y palabras más repetidas. La
        proporción de diálogo distingue de verdad entre obras —\emph{Alice} 73\pct,
        \emph{Pride and Prejudice} 28\pct, \emph{Crime and Punishment} 31\pct—, cosa que
        no hacía cuando se medía por párrafo en lugar de por oración, porque un epub no
        conserva los saltos de párrafo y las tres daban 100\pct.
  \item Pide la curva de dificultad. Sobre \emph{Crime and Punishment} (PG2554), el modelo
        recorre \textbf{840 ventanas de 203.699 palabras en 134 segundos} y devuelve una
        media de \textbf{1,70 sobre 2} —orientativo C1—, además de decir si el libro se va
        complicando o soltando y cuáles son su capítulo más llevadero (el \textbf{19},
        1,52) y su capítulo más exigente (el \textbf{36}, 1,87). La segunda vez
        tarda 0,02~s porque queda guardada. Las figuras~\ref{fig:curva_wp}
        y~\ref{fig:curvas_comp} muestran ese producto sobre tres libros medidos.
  \item Mantiene presionada una palabra que no conoce (con ratón, doble clic). El
        trabajo del servidor —lema, categoría gramatical, consulta al diccionario,
        desambiguación por similitud, IPA— son \textbf{40 milisegundos} \emph{en
        caliente}; lo que ve el lector, con la red de por medio, son unos 0,7~s
        (0,64--0,73 medidos) y unos 5~s la primera consulta de la sesión, que carga el
        diccionario, spaCy y el modelo de similitud. Medido el 15 de septiembre de 2026
        sobre el \emph{front} publicado, a 390~px. La explicación
        contextual, que sale a la red, es lo único lento siempre, y por eso la ficha se
        pinta en dos etapas: lo local aparece al instante y la explicación se añade
        cuando llega.

        Esa explicación se pide con \textbf{salidas estructuradas} (\emph{structured
        outputs}): un esquema Pydantic con un campo por cada cosa que se va a pintar
        ---definición en contexto, explicación gramatical, versión adaptada al nivel y
        ejemplo simplificado--- en lugar de texto libre que luego haya que trocear.
        El esquema fija la forma de la respuesta, \textbf{no garantiza su veracidad}.
        Si Kaikki trae una acepción utilizable, la definición se conserva y el servicio
        explica y traduce el contexto; ante una elección dudosa solo puede seleccionar
        otra glosa existente. Desde el 13 de septiembre, si no hay acepción utilizable
        y se autoriza la llamada, puede ofrecer una definición \textbf{marcada como
        generada por IA}, nunca atribuida al diccionario. Sin clave se conserva la ficha
        local y el aviso de ausencia. Las salidas estructuradas también se usan en
        pronunciación y en «¿de qué trata?».
  \item Lee una frase en voz alta y recibe la corrección palabra por palabra.
\end{enumerate}

Ese recorrido es la unidad de valor del producto: \textbf{no es un modelo con una
interfaz encima, es una secuencia de decisiones que el lector no podía tomar antes}.

% --------------------------------------------------- FIGURA: curva War and Peace
\begin{figure}[htbp]
\centering
\begin{tikzpicture}
\begin{axis}[
  ejell,
  width=\linewidth,
  height=7.2cm,
  xlabel={Capítulo (orden de lectura)},
  ylabel={Dificultad en la escala 0--2},
  xmin=0, xmax=366,
  ymin=1.25, ymax=2.06,
  xtick={1,50,100,150,200,250,300,362},
  ytick={1.3,1.4,1.5,1.6,1.7,1.8,1.9,2.0},
  ymajorgrids=true,
  legend style={at={(0.5,-0.20)}, anchor=north, legend columns=3},
]

% banda "avanzado" (>= 1,4), que es el corte vigente
\addplot[draw=none, forget plot, name path=techo] coordinates {(0,2.06) (366,2.06)};
\addplot[draw=none, forget plot, name path=corte16] coordinates {(0,1.4) (366,1.4)};
\addplot[llrojo!7, forget plot] fill between[of=techo and corte16];

% serie cruda, aligerada
\addplot[llgris!55, line width=0.4pt, each nth point=2,
         filter discard warning=false, unbounded coords=discard]
  table[x=capitulo, y=dificultad] {datos_figuras/curva_War_and_Peace_ma.dat};
\addlegendentry{capítulo a capítulo (362 medidos)}

% media móvil con TODOS los puntos
\addplot[llazul, line width=1.4pt]
  table[x=capitulo, y=media_movil] {datos_figuras/curva_War_and_Peace_ma.dat};
\addlegendentry{media móvil centrada, 15 capítulos}

\addplot[llverde, dashed, line width=1.0pt, domain=0:366, samples=2] {1.735};
\addlegendentry{media del libro: 1,735}

% extremos
\addplot[only marks, mark=*, mark size=2.4pt, forget plot,
         mark options={fill=llverde, draw=llverde!60!black}]
  coordinates {(34,1.333)};
\addplot[only marks, mark=*, mark size=2.4pt, forget plot,
         mark options={fill=llrojo, draw=llrojo!60!black}]
  coordinates {(202,1.990)};

\node[font=\scriptsize, color=llverde!55!black, anchor=west] at (axis cs:42,1.333)
  {capítulo más llevadero: 34 (1,333)};
\node[font=\scriptsize, color=llrojo!70!black, anchor=south] at (axis cs:202,2.000)
  {capítulo más exigente: 202 (1,990)};
\node[font=\scriptsize\itshape, color=llrojo!60!black, anchor=north east]
  at (axis cs:363,2.045) {banda «avanzado» ($\geq 1{,}4$)};

\end{axis}
\end{tikzpicture}
\caption[Curva de dificultad de \emph{War and Peace}]{\textbf{La figura que mejor explica
el producto.} Los 362 capítulos de \emph{War and Peace} medidos sin muestrear —2.459
ventanas de 250 palabras, análisis completo— sobre la escala pública 0--2. La serie fina
es la medición capítulo a capítulo (dibujada con un punto de cada dos por legibilidad); la
línea gruesa es la media móvil centrada de 15 capítulos, calculada sobre \emph{todos} los
puntos. Un libro no tiene una dificultad: tiene un recorrido. Aquí va de 1,333 a 1,990
—0,657 de la escala— con media 1,735. Lo que la curva enseña no es una pendiente sino
una meseta con dientes: los tres tercios promedian 1,723, 1,733 y 1,747, y el máximo de
la media móvil cae en el capítulo 289, ya en el último tercio. La variación que importa
al lector está entre capítulos vecinos, no entre el principio y el final. Fuente: \cod{datos\_figuras/\allowbreak curva\_War\_and\_\allowbreak Peace.dat}.}
\label{fig:curva_wp}
\end{figure}

% --------------------------------------------------- FIGURA: curvas comparadas
\begin{figure}[htbp]
\centering
\begin{tikzpicture}
\begin{groupplot}[
  group style={group size=2 by 1, horizontal sep=0.9cm},
  ejell,
  width=0.47\linewidth,
  height=5.4cm,
  ymin=1.15, ymax=2.0,
  ymajorgrids=true,
  xlabel={Capítulo},
  ylabel={Dificultad (0--2)},
]

\nextgroupplot[title={\emph{Crime and Punishment}}, xlabel={Capítulo (30 medidos)},
  xmin=0, xmax=31, xtick={1,10,20,30}]
\addplot[llazul, line width=0.9pt, mark=*, mark size=1.1pt]
  table[x=capitulo, y=dificultad] {datos_figuras/curva_Crime_and_Punishment.dat};
\addplot[llgris!80, dashed, line width=0.8pt, domain=0:31, samples=2] {1.6793};
\node[font=\scriptsize, color=llgris, anchor=south west] at (axis cs:0.6,1.16) {media del libro: 1,679};

\nextgroupplot[title={\emph{Alice in Wonderland}}, xlabel={Capítulo (12 medidos)},
  ylabel={}, xmin=0, xmax=13, xtick={1,4,8,12}]
\addplot[llverde!80!black, line width=0.9pt, mark=*, mark size=1.4pt]
  table[x=capitulo, y=dificultad] {datos_figuras/curva_Alice_s_Adventures_in_Wonder.dat};
\addplot[llgris!80, dashed, line width=0.8pt, domain=0:13, samples=2] {1.4968};
\node[font=\scriptsize, color=llgris, anchor=south west] at (axis cs:0.3,1.16) {media del libro: 1,497};

\end{groupplot}
\end{tikzpicture}
\caption[Curvas de dificultad en modo vistazo]{La dispersión interna es la norma, no la
excepción. Dos libros más, medidos en \emph{modo rápido} (el vistazo que la aplicación
ofrece en segundos): \emph{Crime and Punishment} recorre 1,270--1,891 y \emph{Alice}
1,296--1,878, pese a que sus medias las separa casi dos décimas. Elegir libro por su media
esconde justamente esto. Fuente: \cod{datos\_figuras/\allowbreak curva\_Crime\_and\_\allowbreak Punishment.dat} y
\cod{datos\_figuras/\allowbreak curva\_Alice\_s\_\allowbreak Adventures\_in\_\allowbreak Wonder.dat}.}
\label{fig:curvas_comp}
\end{figure}

\subsection{Escalado: el catálogo entero, precalculado}
\label{sec:escalado}

La curva de dificultad de un libro largo son minutos de cómputo. Hacer esperar eso a cada
visitante no es aceptable, de modo que el catálogo se precalcula. Las mediciones que
hicieron viable esa decisión están en el cuadro~\ref{cuad:velocidad} y en la
figura~\ref{fig:velocidad}.

\begin{table}[htbp]
\centering
\small
\begin{tabular}{@{}l l R{2.4cm} R{1.8cm}@{}}
\toprule
\textbf{Entorno} & \textbf{Configuración} & \textbf{Ventanas/s} & \textbf{Factor} \\
\midrule
CPU (Ryzen 5 2600) & 12 hilos, lote 64 & 8,19 & 1$\times$ \\
GPU local (GTX 1650 SUPER) & fp32, lote 128 & 112,7 & 13,8$\times$ \\
GPU A100 (Colab) & fp16, lote grande & 2.022 & 247$\times$ \\
\bottomrule
\end{tabular}
\caption{Rendimiento medido de extremo a extremo, incluida la tokenización y las
transferencias. Son cifras del precálculo, que agrupa en lotes fijos de 64 en CPU y 128
en GPU. \textbf{La aplicación en vivo no llega a esos lotes}: agrupa por capítulo, y no por
conveniencia sino para poder ir repintando la curva mientras mide; un capítulo son
unas veinte ventanas, así que el
rendimiento interactivo real está más cerca de las 6,96 ventanas/s del agrupado por
capítulo que de las 8,19 de esta fila. Fuente:
\cod{datos\_figuras/\allowbreak velocidad.dat}.}
\label{cuad:velocidad}
\end{table}

% ---------------------------------------------------- FIGURA: velocidad
\begin{figure}[htbp]
\centering
\begin{tikzpicture}
\begin{axis}[
  barrall,
  width=0.86\linewidth,
  height=5.0cm,
  xbar,
  bar width=13pt,
  xmode=log,
  log origin x=infty,
  xmin=1, xmax=12000,
  xlabel={Ventanas de 250 palabras por segundo \emph{(escala logarítmica)}},
  symbolic y coords={A100,GPU_local,CPU},
  ytick={A100,GPU_local,CPU},
  yticklabels={GPU A100 (Colab)\\{\scriptsize fp16, lote grande}, GPU local (GTX 1650 S)\\{\scriptsize fp32, lote 128}, CPU (Ryzen 5 2600)\\{\scriptsize 12 hilos, lote 64}},
  yticklabel style={align=right, font=\scriptsize},
  enlarge y limits=0.30,
  xmajorgrids=true,
]
\addplot[fill=llverde!35, draw=llverde!85] table[x=ventanas_por_segundo, y=entorno]
  {datos_figuras/velocidad.dat};

\node[font=\scriptsize\bfseries, color=llgris, anchor=west] at (axis cs:10.5,CPU) {8,19};
\node[font=\scriptsize\bfseries, color=llgris, anchor=west] at (axis cs:135,GPU_local) {112,7};
\node[font=\scriptsize\bfseries, color=llgris, anchor=west] at (axis cs:2430,A100) {2.022};

\node[font=\scriptsize\itshape, color=llgris, anchor=west] at (axis cs:5200,CPU) {$1\times$};
\node[font=\scriptsize\itshape, color=llgris, anchor=west] at (axis cs:5200,GPU_local) {$13{,}8\times$};
\node[font=\scriptsize\itshape, color=llgris, anchor=west] at (axis cs:5200,A100) {$247\times$};
\end{axis}
\end{tikzpicture}
\caption[Velocidad de inferencia por entorno]{Tres órdenes de magnitud entre el equipo del
autor y la GPU alquilada, medidos de extremo a extremo (tokenización y transferencias
incluidas). La escala es logarítmica porque el rango va de 8,19 a 2.022 ventanas por
segundo. Es lo que convierte «medir los 61.642 libros en inglés de Project Gutenberg» en
unas pocas horas de cómputo en lugar de un mes seguido de CPU: a 8,19 ventanas/s y
350 ventanas por libro son \textbf{731,7 h}, o sea \textbf{30,5 días}. Datos de
\cod{datos\_figuras/\allowbreak velocidad.dat} (cuadro~\ref{cuad:velocidad}).}
\label{fig:velocidad}
\end{figure}

Dos comprobaciones acompañan a ese cuadro, porque acelerar no puede costar exactitud:

\begin{itemize}[leftmargin=*, itemsep=3pt]
  \item \textbf{La GPU da los mismos números que la CPU.} Sobre 840 ventanas de
        \emph{Crime and Punishment}, 839 son idénticas a tres decimales y una difiere en
        0,001; el MAE es 0,000001 sobre la escala 0--2.
  \item \textbf{Reagrupar en lotes distintos no mueve ninguna cifra que se publique}, y
        conviene decirlo con precisión, porque hay dos casos y solo uno es exacto. Cuando
        se trocea \textbf{un texto de más de 250 palabras}, todas sus ventanas miden 250
        palabras exactas y, en prosa inglesa real, pasan siempre de 256 \emph{tokens} BPE,
        así que se truncan a 256 —comprobado: la ventana más corta de un texto de prueba
        da 258; un texto artificial de 250 monosílabos daría 252 y no llegaría, pero eso
        no es prosa—: no hay
        relleno, y la fila de cada ventana no depende de con quién vaya agrupada. Es el
        caso de la comprobación: las mismas \textbf{300 ventanas} puntuadas de una en una
        y en un solo lote de 300 salen \textbf{300/300 idénticas}, con MAE 0,000000. El
        precálculo del catálogo, en cambio, mete en el mismo lote las ventanas de
        \textbf{todos los capítulos} de un libro, y un capítulo de menos de 250 palabras
        aporta una ventana corta: ahí el relleno sí actúa hasta la ventana más larga del
        lote, y el resultado \textbf{no es idéntico bit a bit}. También está medido, sobre
        \emph{Alicia en el País de las Maravillas}, 112 ventanas: \textbf{103 coinciden
        exactamente y el resto difiere como mucho en $4{,}8\times 10^{-7}$}. Las cifras del
        catálogo se guardan con tres decimales, así que una diferencia de ese tamaño no
        puede mover ninguna.
\end{itemize}

Con 2.022 ventanas por segundo, medir los 61.642 libros en inglés de Project Gutenberg
salían \textbf{entre 3 y 5 horas de A100}, y en su momento se dijo de qué dependía esa
horquilla en vez de dar una cifra sola: del número medio de ventanas por libro. Con las
350 que asume el código del precálculo salen \textbf{2,96 h}; con las 606,4 que promedian
las cinco novelas largas medidas en local —una muestra sesgada hacia arriba, porque son
clásicos extensos— salen \textbf{5,14 h}.

\medskip
\noindent\textbf{Después se ejecutó, y la estimación resultó ser la mitad de la
respuesta.} El ritmo real medido en la A100 fue de \textbf{61 libros por minuto}, o sea
\textbf{unas 17 horas} para los 61.642 —cinco veces la cifra de tarjeta—. Y el desglose
dice exactamente dónde se va el tiempo: un libro medio son \textbf{0,17 s de GPU} contra
\textbf{\textasciitilde1 s de reloj}, así que \textbf{el 80 \% del tiempo la tarjeta está
parada} esperando a que el sistema de archivos de Google Drive entregue el \texttt{.txt}.

\medskip
\noindent La estimación no era falsa: era \textbf{incompleta}, y de una manera que
conviene nombrar porque es un error de método y no de aritmética. Medía la parte del
sistema que se sabía medir —el modelo— y daba por supuesto que era la que manda. No lo
era. Es la misma lección que la agrupación por capítulo en vez de por libro, y la misma
que la descarga: \textbf{el cuello de botella de este proyecto nunca ha estado en el
modelo}. En serie, a 0,3 libros/s, la descarga tardaría 57 horas. Y no
era la pausa de cortesía entre peticiones, era la latencia —cada petición al espejo se
pasa casi todo su tiempo esperando la ida y vuelta—, así que con seis hilos esas esperas
se solapan. Seis y no treinta a propósito: el espejo lo mantiene alguien de forma gratuita,
y cada hilo conserva su pausa.

\medskip
\noindent\textbf{Después se buscó una vía más rápida, y la decisión que importa no fue
técnica.} El archivo único de Gutenberg —11,3 GB, una petición en lugar de las 29.668 que
faltaban el 9 de septiembre— parecía el atajo evidente y no lo era: medido en dos máquinas distintas, el espejo
entrega \textbf{0,29 MB/s}, o sea \textbf{casi once horas}, y con seis conexiones
simultáneas sube a 0,46 y rechaza la séptima. No es el ancho de banda de nadie: es un
servidor que mantiene un voluntario. La lección es del mismo tipo que las dos anteriores:
\textbf{un atajo no es un atajo hasta que se cronometra}.

\medskip
\noindent La vía rápida existe y también se midió: pedir los libros uno a uno a
\texttt{www.gutenberg.org} con treinta y dos hilos da \textbf{3.362 libros por minuto},
unos nueve minutos para lo que faltaba. \textbf{Se descartó.} La política del sitio lo
dice sin ambigüedad —«está pensado para usuarios humanos; cualquier uso aparente de
herramientas automatizadas resultará en el bloqueo de su dirección IP»— y autoriza tres
caminos: un espejo privado, el punto de acceso \texttt{/robot/harvest} con dos segundos de
espera entre ficheros, o el catálogo. Que una puerta esté abierta y sea rápida no la
convierte en una entrada.

\medskip
\noindent Lo que sí resuelve el problema es no pedirle nada a Gutenberg.
\texttt{common-pile/project\_gutenberg} es un volcado ya publicado \parencite{kandpal2025commonpile} de esos mismos textos
de dominio público, servido por una red de distribución de contenidos. Medido sobre el
primer fragmento: \textbf{8,0 MB/s}, unos \textbf{veintidós minutos} para los 10,5 GB, y
el identificador de cada documento \textbf{es} el número de Gutenberg —239 de 239 cruzaron
con el catálogo oficial en la comprobación, sin una sola excepción—. Veintisiete veces más
rápido que el archivo único, y con la infraestructura de los voluntarios intacta.

\subsection{Dos decisiones de producto que merecen argumento}
\label{sec:producto}

\textbf{Se ofrecen dos velocidades de análisis y siempre se dice cuál se usó.} El vistazo
rápido toma 30 capítulos repartidos por la obra y 2 ventanas de cada uno; el análisis
completo no muestrea. El número 30 no es redondo por gusto: sale de la velocidad medida de
la máquina. Y el resumen declara cuál de los dos se ejecutó, porque un vistazo que
presumiera de haber leído el libro entero acabaría citado como si lo hubiera hecho. La
diferencia entre ambos modos se ve al comparar la figura~\ref{fig:curva_wp} (completo, 362
capítulos) con la figura~\ref{fig:curvas_comp} (vistazo).

\textbf{No se filtra el catálogo por contenido.} Se implementó un filtro basado en las
materias que asigna el propio Project Gutenberg —90 libros de 61.642— y después se retiró,
deliberadamente. El razonamiento: esto es una plataforma de literatura, y la literatura es
donde esas cosas aparecen. Filtrar por etiqueta habría eliminado también obras que son
literatura por derecho propio, como \emph{Fanny Hill}, perseguida judicialmente y
relevante precisamente por eso. En su lugar se avisa, y el aviso ya está en la
aplicación: una lista curada a mano (\cod{avisos\_contenido.csv}) pone en la ficha del
libro un recuadro «Aviso sobre el contenido» antes de abrirlo, sin bloquear nada. A
fecha de \cod{db366b8} lleva 2 títulos, las dos ediciones de \emph{Little Black Sambo}
de Helen Bannerman (PG17824 y PG1330), por estereotipos raciales. El buscador no
cambia de orden por el aviso; al recomendar por nivel, los 5 primeros puestos son
para títulos sin aviso y los que lo llevan siguen detrás, sin quitar ninguno. Aparte,
el resumen que escribe el modelo puede añadir su propio «Antes de empezar». \textbf{Se
informa al lector y se le deja elegir; no se elige por él.} Es una decisión discutible y
por eso se argumenta en lugar de ocultarse.

\subsection{Contra quién compite esto}
\label{sec:competencia}

Un proyecto que no sabe contra qué compite no puede decir qué aporta. La versión larga
está en el documento de estrategia; aquí van las tres alternativas reales y la cifra que
resume cada una.

\begin{table}[htbp]
\centering
\small
\begin{tabular}{@{}l R{3.6cm} R{3.2cm}@{}}
\toprule
\textbf{Alternativa} & \textbf{Cifra publicada} & \textbf{La nuestra} \\
\midrule
Fórmulas clásicas (FKGL, ARI) & Pearson \textbf{0,517} en CLEAR & \textbf{0,89} \\
Mejor LLM \emph{zero-shot} publicado & Pearson \textbf{0,725} en CLEAR & \textbf{0,89} \\
Lexile, tarifa de creador de contenido & \textbf{3.824--54.636} USD/año \parencite{metametrics2026lexile} & A100 de pago hasta el 10 sep y T4 gratuita después; sin coste en USD registrado \\
\bottomrule
\end{tabular}
\caption{Cifras publicadas frente a las de este proyecto. \textbf{No son comparaciones
cara a cara}: las dos primeras filas vienen de un estudio ajeno \parencite{grossman2026zeroshot} sobre el mismo corpus, con
otra partición y otro protocolo. Justifican la decisión de afinar un codificador; no
demuestran una victoria medida.}
\label{cuad:competencia}
\end{table}

\begin{figure}[htbp]
\centering
\begin{tikzpicture}
\begin{axis}[
  barrall,
  width=0.82\linewidth,
  height=4.6cm,
  xbar,
  bar width=11pt,
  xmin=0, xmax=1.02,
  xlabel={Correlación de Pearson con el juicio humano, corpus CLEAR},
  symbolic y coords={literalevel,llm,formulas},
  ytick={formulas,llm,literalevel},
  yticklabels={Fórmulas clásicas,LLM \emph{zero-shot},\textbf{Nuestro codificador}},
  yticklabel style={font=\scriptsize},
  enlarge y limits=0.28,
  nodes near coords={\pgfmathprintnumber[fixed,precision=3]{\pgfplotspointmeta}},
  nodes near coords align={horizontal},
]
\addplot[fill=llgris!25, draw=llgris!70] coordinates {(0.517,formulas)};
\addplot[fill=llgris!25, draw=llgris!70] coordinates {(0.725,llm)};
\addplot[fill=llcian!45, draw=llcian!90] coordinates {(0.89,literalevel)};
\end{axis}
\end{tikzpicture}
\caption[Comparación con las alternativas]{Las fórmulas clásicas cuentan sílabas y llegan
a 0,517. Un modelo de lenguaje grande con un buen \emph{prompt}, sin afinar, llega a
0,725. El codificador afinado de este proyecto, a 0,89. Las tres cifras son sobre el mismo
corpus, pero —y esto hay que repetirlo aquí, no solo en el pie del
cuadro~\ref{cuad:competencia}— \textbf{no son comparaciones cara a cara}: las dos
primeras vienen de un estudio ajeno, con otra partición y otro protocolo, y el 0,89 es
el del \cod{roberta-base} afinado sobre CLEAR del~\S\ref{sec:progresion}, que
\textbf{no es el modelo que se despliega}. Justifican la decisión de afinar un
codificador; no demuestran una victoria medida.}
\label{fig:competencia}
\end{figure}

Los 53.179 libros del catálogo medido caen en el tramo más alto de la tarifa de Lexile. Medirlos
aquí sí usó cómputo de pago, y no hay cifra en dólares que dar: no quedó registrado cuántas
unidades se gastaron. Hasta el 10 de septiembre de 2026 el cuaderno de Colab se abría con una
A100, que consume unidades de pago de la cuenta Colab Pro, y hasta esa tarde se midieron
27.776 libros (los 29.875 medidos el 10 menos los 2.099 de la primera tanda en T4); el 9 de
septiembre quedaban 1,9 unidades (\cod{colab\_precalculo.py:903}). Desde entonces se midió en una T4 gratuita, en
tandas reanudables ---la primera, 2.099 libros; la última, 9.795 en 55 minutos---, y los
textos libres en la CPU local. Las «2,8 h de A100» del registro del proyecto no son horas
gastadas: son la estimación de tarjeta para los 61.642 libros (332,3 ventanas por libro a
2.022 por segundo). Al ritmo medido en la A100, 61 libros por minuto, los 53.179 habrían sido
unas 14,5 horas de reloj (\S\ref{sec:escalado} da 17 para los 61.642). \textbf{No es una razón
monetaria estricta} —una licencia anual y unas horas de cómputo no se dividen— y tampoco
son servicios equivalentes: Lexile vende una escala aceptada por instituciones. Lo que sí
se sostiene es que el \textbf{coste de procesamiento está en otro orden de magnitud}, y eso
es lo que hace viable medir un catálogo entero en un trabajo de estudiante.

\textbf{Dónde se pierde, que también hay que decirlo.} Lexile tiene treinta años de
validación y una escala que \textbf{acepta un colegio}; la de aquí no la acepta nadie.
\emph{TeakReader} \parencite{teakreader} ya da nivel MCER por capítulo y es una aplicación publicada en dos
tiendas, mientras esto vive en un enlace. Y las lecturas graduadas de editorial tienen algo
que ningún modelo da: el nivel se lo puso una persona que leyó la obra.

\textbf{Lo que queda, y basta para justificar el trabajo}, son tres cosas: se recorre el
libro \textbf{entero, capítulo a capítulo}, en lugar de dar un número global; está escrito
para un \textbf{hispanohablante que elige qué leer}, no para un profesor que clasifica
textos; y es \textbf{auditable en lo que se puede auditar y de acceso gratuito para el
lector}, con el modelo publicado y el código a la vista.

Esa tercera hay que acotarla: el código, la inferencia desde los pesos públicos \textbf{y}
las métricas internas del v4 se pueden rehacer ---recalculadas el 12 de septiembre de 2026,
idénticas a cuatro decimales---. Y las
fórmulas clásicas también son gratuitas y reproducibles, así que la ventaja no está en la
auditabilidad sola, sino en tenerla junto con un modelo que lee el texto. Lo que ninguna
alternativa reúne son las \textbf{tres a la vez}.


% ===========================================================================
\section{Limitaciones}
\label{sec:limitaciones}

\textbf{Antes de la lista, de dónde sale.} Este trabajo se sometió a \textbf{quince
auditorías externas} hechas por modelos de dos familias distintas, con el encargo
explícito de ser duros y \textbf{a ciegas}: sin decirles qué se había corregido antes,
para que lo que volviera a salir fuera lo que seguía mal. Entre todas señalaron más de
un centenar de hallazgos: fugas de datos, cifras que no cuadraban entre
archivos, métricas atribuidas al corpus equivocado, pruebas que no podían fallar y
conclusiones más fuertes que su evidencia. Los informes están en \cod{auditoria/} y las
correcciones, en el historial del repositorio.

Varias de las limitaciones que siguen no las encontró el autor, y algunas contradicen lo
que este mismo documento afirmaba en una versión anterior. Parece razonable decirlo una
vez, al principio de la lista, en lugar de repetirlo en cada párrafo corregido.
% ===========================================================================

\begin{enumerate}[leftmargin=*, itemsep=4pt]
  \item \textbf{La escala 0--2 es propia y no es una certificación.} El MCER fue el ancla
        del entrenamiento. LiteraLevel no acredita ningún nivel, y decirlo en el documento
        no basta: la aplicación lo repite en pantalla.
  \item \textbf{La escala está comprimida.} El ajuste contra el ancla real ronda
        $\text{predicho} \approx 0{,}72 \times \text{real} + 0{,}38$
        (figura~\ref{fig:compresion}; mínimos cuadrados sobre las cuatro medias por nivel
        de Colab, y $0{,}67x + 0{,}41$ con las cinco locales), así que los niveles bajos se
        sobreestiman y los altos se subestiman. Es la razón de mostrar bandas y no decimales.
  \item \textbf{Las letras del catálogo son orientativas.} Ya no son heredadas: los 53.179
        libros están medidos enteros por el v4 (cuadro~\ref{cuad:catalogomedido}). Pero la
        escala del modelo sale comprimida y desplazada hacia arriba ---un texto Elementary
        de OneStopEnglish puntúa 0,93---, así que la letra de un libro fácil tiende a subir.
        La aplicación lo dice al pie de cada tabla, en la pantalla y no solo en el PDF.
        Hasta el 12 de septiembre de 2026 esta limitación era peor: con perfil 1,20,
        «Sorpréndeme» proponía \emph{The Secret Adversary} con etiqueta heredada
        \textbf{B2}, y el v4 leyendo el texto devolvía \textbf{C1}.
  \item \textbf{La literatura del catálogo es, en su inmensa mayoría, C1--C2}: 1 libro A1,
        40 A2 y 557 B1 de 53.179. Para un lector B1 la respuesta no son libros enteros sino
        capítulos (3.644 en 1.499 libros), y por nivel solo se ofrecen lecturas literarias:
        288 entre A1 y B1.
  \item \textbf{Los tamaños muestrales de la validación externa son pequeños}: 38 textos
        en Colab, 47 en local, con 3 textos en A1. Una diferencia de 0,05 entre dos
        niveles adyacentes con esas $n$ no es significativa; lo que se afirma es el orden
        global, no cada comparación por pares.
  \item \textbf{La validación no es enteramente independiente de la exploración.} El
        corpus de British Council se usó primero para diagnosticar (descubrir que el
        modelo no sabe leer menús) y después para validar. El filtro de prosa se definió
        por dos vías independientes precisamente para acotar ese problema, pero salió de
        ese diagnóstico y el conjunto ha sido mirado y ya no es virgen.
  \item \textbf{La longitud del texto es un confusor en el corpus de examen.} Contar
        palabras alcanza 0,807 donde el modelo alcanza 0,860 (cuadro~\ref{cuad:baseline}).
  \item \textbf{El modelo solo mide prosa continua en inglés.} Poesía, teatro, diálogo muy
        fragmentado y cualquier lengua distinta del inglés quedan fuera.
  \item \textbf{El corpus de anclaje MCER es pequeño y de una sola fuente}: 654 documentos
        aceptados de ELG CEFR English. Hasta el 12 de septiembre de 2026 ni su composición
        ni el guion de entrenamiento estaban en el repositorio, y una auditoría externa lo
        señaló con razón: sin esa lista nadie podía comprobar que las fuentes de validación
        fueran ajenas a las de entrenamiento. \textbf{Ya están}: guion, tarjeta de métricas y
        las 144 predicciones, en \cod{entrenamiento\_v4/}, que no viaja con los
        entregables y se comparte a petición. Y la comprobación resultó necesaria
        de verdad: la primera validación con UniversalCEFR se hizo sobre este mismo corpus y
        hubo que retirarla.
  \item \textbf{British Council no se puede redistribuir, pero sí reconstruir.} Tiene
        derechos y no acompaña a la entrega; se reconstruye desde su manifiesto de URLs y
        SHA-256, que sí viaja —va dentro del cuaderno— y se verifica archivo por archivo.
        Comprobado el 15 de septiembre de 2026 sin cuenta: las 59 URLs responden desde un
        navegador y los 59 archivos dan el mismo SHA-256 que el manifiesto
        (\cod{auditoria/\allowbreak ronda4\_enlaces\_4e\_\allowbreak 2026-09-15.md}); antes,
        desde una carpeta vacía el 12 de septiembre con \cod{tools/\allowbreak
        bajar\_corpus\_bc.py}, salían las mismas cifras del cuadro~\ref{cuad:validacion}.
        Ese guion, \cod{validar\_local.py} y \cod{entrenamiento\_v4/} \textbf{no viajan con
        los entregables y se comparten a petición}: lo que cualquiera puede comprobar por su
        cuenta es la cadena del cuaderno —el modelo v4 en Hugging Face, CLEAR, el catálogo de
        Gutenberg y estos hashes—. Las cifras internas del v4 (0,908 y 97,2\pct) se
        recalculan con el guion y las 144 parejas de \cod{entrenamiento\_v4/}.
  \item \textbf{El cuaderno corre entero en un Colab limpio, pero correr no es
        recalcular.} Los pesos del v4 están publicados en Hugging Face
        (\cod{MikeMaki1763/\allowbreak literalevel-v4}, repositorio público, descarga sin
        credenciales), así que el cuaderno se ejecuta de principio a fin y demuestra la
        inferencia sobre texto real. Ejecutado así el 16 de septiembre de 2026 en Colab, con
        CPU y cuenta gratuita, y con las salidas guardadas en el propio cuadernillo: 57 de 57
        celdas de código con salida y ningún error. El cuaderno no reentrena el v4, que
        pide una GPU, pero \textbf{sus métricas sí se recalculan fuera de él}: los datos de
        entrenamiento son
        públicos (ELG CEFR English, en Hugging Face), y el guion y las 144 parejas de nivel
        y predicción están en \cod{entrenamiento\_v4/}, que no viaja con los entregables y se
        comparte a petición. Recalculadas en CPU en menos de un minuto, dan las mismas cifras
        a cuatro decimales.
  \item \textbf{El despliegue está en el escalón más barato, y eso tiene dos
        consecuencias.} La aplicación está publicada en dos Spaces de Hugging Face con
        dirección fija que sobrevive a apagar el equipo del autor: el \emph{front} de
        Angular en un Space estático (\cod{MikeMaki1763/\allowbreak literalevel-front}),
        gratuito y sin suspensión, y la API con el motor de Gradio en
        \cod{MikeMaki1763/\allowbreak literalevel}, hardware \cod{cpu-basic}. Este segundo
        cuesta \textbf{7 USD al mes}: un Space de Gradio exige cuenta Pro, y
        \cod{cpu-basic} es el hardware que no añade coste sobre esa suscripción. Pero ese
        hardware \textbf{suspende el servicio tras 48 h sin visitas} (\cod{gcTimeout} $=$
        172.800 s según la API de Hugging Face), de modo que el \emph{front} abre al
        instante, pero su primera llamada a la API después de una pausa paga un arranque
        en frío mientras se cargan \emph{torch} y los 328 MB del modelo; el \emph{front} lo
        dice con palabras («tarda unos 30 segundos en despertar», \cod{nucleo/api.ts},
        que se queda corto) en vez de fallar. Está medido, y cambia de un arranque a otro:
        en el despliegue del 15 de septiembre (\cod{0319290}, con la imagen del contenedor
        ya construida), \textbf{46 s hasta que el runtime dijo \cod{RUNNING}} y el buscador
        listo —\cod{buscador\_listo} en \cod{/api/salud}, precargado al arrancar desde el
        reinicio 26— \textbf{14 s después}; en los otros arranques medidos ese día, entre
        12 y 55 s después de \cod{RUNNING} (\cod{DESPLIEGUE.md}). Una medida anterior,
        sobre el Space sin el buscador precargado, dio 31 s hasta \cod{RUNNING} y 32 s
        hasta que la página respondía. Todas son cotas inferiores: despertar de una
        suspensión de 48 h añade traer la imagen. El disco del
        contenedor \textbf{sí es efímero}, y por eso la estantería, el avance de
        lectura, el cuaderno y el perfil del lector se guardan en un \textbf{bucket de
        almacenamiento montado en \cod{/data}} (18 USD por TB y mes en repositorio
        privado, y esto ocupa megas). Verificado de punta a punta: se fijó un perfil,
        se reinició el Space y el perfil seguía ahí. Era la acción A4 de la estrategia
        y está cerrada.
  \item \textbf{La prueba de usuario tiene $n = 3$.} Esmeralda y Laura son dos lectoras,
        y el tercero es Ismael, que la probó en un Android los días 14 y 15 de septiembre:
        pidió que se entendiera mejor, dijo que al pulsar el micrófono no pasaba nada y que
        las pantallas se sentían saturadas. De ahí salieron las tres tandas de ayuda; la
        causa de lo del micrófono \textbf{está en revisión}, porque se supo después que no
        había iniciado sesión (apartado~\ref{sec:ismael}).
        Su valor es cualitativo y de diagnóstico, no estadístico: sirve para encontrar
        fallos de usabilidad graves —y encontraron varios—, no para estimar satisfacción
        media. Desde el 14 de septiembre hay cuentas repartidas a los primeros usuarios (\cod{auditoria/\allowbreak feedback\_usuarios\_\allowbreak 2026-09-14.md}), y la primera respuesta ya enseñó un fallo: una lectora B1 buscó libros de amor y encontró pocos. No faltaban en el catálogo; fallaba la búsqueda. En la franja que la aplicación da a un lector B1 (dificultad de 0,35 a 1,25) hay 114 textos libres con \emph{love} o \emph{romance} en el título, frente a 4 libros de Gutenberg (recuento sobre el catálogo que carga la aplicación, \cod{auditoria/\allowbreak ronda1\_consolidado\_\allowbreak 2026-09-14.md}); pero la búsqueda devolvía los más parecidos de todo el catálogo, y con «amor» o «Historias de amor» y nivel B1 los 24 primeros eran de Gutenberg, sin ninguno de esos cuentos. Desde el reinicio 25 (15 de septiembre), detrás de esos resultados van hasta 12 textos de la franja del lector que se parecen a lo buscado: con «amor» y nivel B1 se marcan 12, y 8 llevan el amor en el título. En \emph{Encontrar}, contados a mano, unos 7 de los resultados de «Historias de amor» que caen en el nivel del lector o un poco por encima son de amor, en sentido amplio. Recontado el 15 de septiembre con el catálogo que carga la aplicación y con el Space publicado (\cod{auditoria/\allowbreak chips\_reinicio25\_\allowbreak 2026-09-15.md}).
  \item \textbf{La explicación contextual depende de un servicio externo de pago.} Sin
        clave, la aplicación sigue funcionando en todo lo que no sale a la red desde el
        servidor: diccionario, fonética, buscador y puntuación. Los enlaces de audio de
        Wiktionary son aparte, porque los descarga el navegador del lector: necesitan
        conexión, aunque no clave. Lo que se apaga sin clave son tres cosas: la
        explicación en contexto, la corrección de pronunciación y el resumen de
        «¿de qué trata?» de la ficha. Con cuenta, cada corrección gasta un uso del cupo del
        perfil, salvo cuando la grabación no trae voz: entonces la aplicación dice «No te
        oí. Acércate al micrófono y vuelve a grabar. Este intento no cuenta.» y devuelve el
        uso, con un tope de 10 devoluciones por cuenta y día UTC (\cod{app/api.py:1945},
        \cod{app/cuentas.py:306}). Si lo que falla es la red o la clave, no hay devolución.
  \item \textbf{La búsqueda por franja tiene tres límites conocidos,} medidos el 15 de
        septiembre sobre el Space del reinicio 28 y dejados sin cambio de código porque
        los arreglos probados estropeaban más búsquedas de las que salvaban
        (\cod{DESPLIEGUE.md}).
        \begin{itemize}
          \item \textbf{Un autor buscado por una sola palabra, con perfil, trae cuentos
                sin relación.} Con un perfil de 1,20, «dostoyevski» marca 15 textos en la
                franja y «tolstoy», «cervantes» y «oscar wilde», 12 cada uno: con una sola
                palabra, la franja compara el título de los textos libres con «historias
                de X» y «libros sobre X», y el modelo de parecido acerca a esas frases
                títulos de misterio, bestias o fantasía aunque no tengan que ver con el
                autor. Quitar el título solo cuando casa con el marco no lo arregla, y
                decidir por el autor del primer resultado quitaba cuentos de su nivel a
                «león», «brown», «hope» o «animales de la granja» (reinicio 29, vuelto
                atrás). Con título y autor no pasa: «Crimen y castigo de Dostoyevski» da 0.
          \item \textbf{En las consultas de varias palabras, el título pesa demasiado:}
                «Historias de amor» marca cuentos titulados «My stories», y contados a mano
                solo unos 5 de sus resultados con A2 y 7 con B1 son de amor; con «Terror
                gótico» pasa lo mismo. Ninguna de las tres variantes probadas mejoraba lo
                desplegado en 14 consultas de calibración (85,4\,\% frente a
                80,7--84,4\,\%).
          \item \textbf{Con una sola palabra, un título que solo la contiene sube:}
                «rosa» trae «The Water Supply of the El Paso…» con A2 y B1, y «luna», con
                B1, «Luna, The Rajapalayam Puppy». El arreglo medido cambiaba 70 de las 200
                listas de las 100 palabras de tema y hacía que «dracul» perdiera
                \emph{Dracula}.
        \end{itemize}
\end{enumerate}

% ===========================================================================
\section{Trabajo futuro}
% ===========================================================================

\begin{enumerate}[leftmargin=*, itemsep=3pt]
  \item \textbf{Hecho el 12 de septiembre de 2026:} remedir el catálogo completo y sustituir
        las etiquetas heredadas por medidas del v4 sobre el libro entero (53.179 libros,
        585.111 capítulos).
  \item \textbf{Hecho el 12 de septiembre de 2026:} publicar el guion de entrenamiento y la
        composición del corpus de anclaje, que resultó ser público (ELG CEFR English). Las
        métricas del v4 ya se recalculan en lugar de citarse.
  \item Ampliar el corpus de anclaje MCER, que hoy es el eslabón más delgado de la cadena
        de evidencia.
  \item Calibrar explícitamente la escala comprimida (una transformación afín de dos
        parámetros ajustada sobre entrenamiento, sin tocar los pesos) y comprobar si el
        orden se conserva tras calibrar.
  \item Repetir la prueba de usuario con cinco o seis lectores de perfiles distintos, que
        según \textcite{nielsen2000five} es donde se detecta la mayor parte de los
        problemas de usabilidad.
  \item Medir el efecto real sobre el abandono: comparar la tasa de libros terminados
        entre lectores que eligen con recomendación y sin ella. Es la única métrica que
        verdaderamente valida la propuesta.
  \item \textbf{Convertir las cinco acciones en herramientas de un agente}
        (\emph{function calling}). Hoy cada una tiene su botón: medir un capítulo,
        buscar por significado, abrir un libro, explicar una palabra y corregir la
        pronunciación. Las cinco ya son funciones con argumentos y salida tipada
        ---que es justo lo que un esquema de herramienta necesita---, así que el
        trabajo no sería reescribirlas sino declararlas y dejar que el modelo elija:
        «quiero algo de misterio victoriano que pueda leer» debería encadenar la
        búsqueda, el filtro por nivel y la ficha sin que el lector sepa que hay tres
        pasos. \textbf{No se hizo por una razón de calendario, no de diseño}: la
        aplicación llevaba dos días congelada y auditada cuando esto se planteó, y
        meter un bucle de herramientas en el camino que ya funciona habría cambiado
        el riesgo de la entrega por una mejora que no estaba medida.
\end{enumerate}

% ===========================================================================
\section{Conclusiones}
% ===========================================================================

\textbf{Se cumplió el objetivo ordinal, que era el objetivo.} Hay dos resultados y vienen
de corpus distintos, así que conviene no mezclarlos. Sobre
\textbf{textos externos con nivel MCER oficial} —los 47 de British Council— el modelo
ordena los niveles de forma estrictamente creciente (0,5040 · 0,5561 · 0,9276 · 1,2414 ·
1,4952; figura~\ref{fig:validacion}) y alcanza Spearman \textbf{0,829} con 38 textos,
\textbf{0,860} con los 47. El \textbf{0,908 y el 97,2\pct} son otra cosa: salen del
conjunto de validación de \textbf{144 documentos del propio corpus de anclaje}, que es
interno. Atribuirlos a los externos sería un error. Para el problema planteado —elegir qué leer— acertar el orden
es lo que importa; acertar la letra exacta es un lujo.

\textbf{Y aquí hay que poner un límite que el documento no ponía.} Ese orden está medido
sobre \textbf{textos de examen}, de A1 a C1 y sin C2, no sobre novelas. La aplicación
puntúa novelas y dibuja sus curvas, pero \textbf{la dificultad de una novela no se ha
contrastado contra ninguna etiqueta externa ni contra ningún lector}. Que el modelo ordene
bien cinco niveles de material de examen es una señal razonable de que la escala tiene
sentido; no es una demostración de que ayude a elegir mejor un libro. Ese vínculo queda
pendiente.

\textbf{En las configuraciones probadas, el cuello de botella no estuvo en el algoritmo
sino en la representación del texto.} Aumentar la potencia del clasificador sobre las
mismas variables no mejoró nada —XGBoost con trescientos árboles queda por detrás de una
regresión logística, en las dos tandas—, mientras que cada vez que se enriqueció la
representación el resultado subió: de conteos a \emph{embedding} y de ahí a leer el texto
crudo con atención, hasta 0,67 niveles de error ordinal (figura~\ref{fig:progresion}).

\textbf{Dicho eso, esto no es un experimento controlado y no conviene venderlo como tal.}
Entre la fila del \emph{embedding} y la del transformador cambian a la vez la
representación, la arquitectura, el preentrenamiento y el objetivo —clasificación de seis
clases frente a regresión—, así que el salto no se puede atribuir a un solo factor. Lo que
la evidencia sostiene es la observación: \textbf{subir la potencia del clasificador no
movió el resultado y cambiar la representación sí}. Aislar cuál de los cuatro factores
pesa más exigiría comparaciones que este trabajo no hizo.

\textbf{Sobre el tamaño de ese salto conviene ser preciso}, porque el
cuadro~\ref{cuad:progresion} tiene dos tandas y es fácil restar filas que no se pueden
restar. Contra las tres variables clásicas de la tanda histórica, el paso al
\emph{embedding} va de 1,39 a 0,93. Pero el cuadernillo entregado ejecuta una regresión
logística sobre tres variables clásicas, sobre esa misma partición, y obtiene
\textbf{1,295}; con diecisiete variables en vez de tres baja a \textbf{1,062}.

\textbf{No son exactamente el mismo modelo, y decir que sí lo eran era impreciso.} Las
tres variables se llaman igual en las dos tandas, pero \cod{palabras\_raras} no contiene
lo mismo: en el cuadernillo se calcula contra las 3.000 palabras más frecuentes
\textbf{del entrenamiento de CLEAR}, y en la tanda histórica contra las listas externas de
NLTK y del SPGC (cuadro~\ref{cuad:saturacion}). Misma partición y misma semilla, sí; misma
definición de las variables, no. Esa es la razón más probable de los 0,10 de diferencia.

Medido así, contra unas variables clásicas bien surtidas, la ventaja del \emph{embedding}
es bastante menor de lo que sugiere la
comparación con tres. \textbf{La brecha que no se cierra con más conteos es la del
transformador.} La conclusión cualitativa
se sostiene; la cifra con la que se vendía, no.

\textbf{Escalar el corpus no siempre mejora nada, y conviene decirlo.} Multiplicar por
2.500 el corpus de referencia léxica —de 18 libros a los 45.345 del SPGC— llevó la
correlación de $-0{,}379$ a $-0{,}326$: \textbf{empeoró}. La muestra intermedia de 397
libros daba $-0{,}330$. La ley de Zipf explica por qué. El paso se
conservó por cobertura, no por métrica, y así se reporta.

\textbf{La velocidad se ganó sin mover ninguna cifra publicada, y con la salvedad dicha.}
Doce
hilos y lotes de 64 en CPU, fp32 con lote 128 en GPU: $+33\pct$ y $\times 13{,}8$
respectivamente (figura~\ref{fig:velocidad}), con 839 de 840 ventanas idénticas a tres
decimales y la restante a 0,001. No es exactitud bit a bit —eso sería falso— sino
diferencias por debajo del redondeo con el que se guarda todo. Y se
rechazaron dos palancas más rápidas —int8 y fp16 en la tarjeta local— porque una desviaba
0,210 sobre una escala de 0 a 2 y la otra era 4,8 veces más lenta. \textbf{Medir antes de
optimizar, y medir también lo que se descarta.}

\textbf{Lo que convierte esto en un producto no es el modelo, es lo que se hizo con lo que
dijo una lectora de quince años.} Un ciclo completo —observar, escuchar la queja,
diagnosticar el fallo propio y corregirlo— vale más que un punto de correlación. Ese ciclo
se documenta en el apartado~\ref{sec:esmeralda}, y su resultado son seis etiquetas
reescritas, cinco paletas verificadas contra el criterio de contraste WCAG AA y una
interfaz que se puede usar con el pulgar.

\textbf{Y queda pendiente lo que más pesa.} La aplicación funciona y ya tiene dirección
fija ---el \emph{front} de Angular en un Space estático y la API con el motor en otro---,
pero el motor corre sobre el hardware más barato: se suspende por inactividad. El catálogo ya está
medido entero ---53.179 libros--- y las métricas del v4, que el documento citaba sin poder
recalcular, ya se recalculan fuera del cuaderno con el guion y las 144 parejas de
\cod{entrenamiento\_v4/}, que se comparten a petición. Lo que queda está planificado con coste
y criterio de terminado en el cuadro~\ref{cuad:accionables}.

\textbf{Y para qué sirve esto, dicho como decisión y no como métrica.} Un lector que hoy
elige libro por la portada puede, con esta herramienta, hacer tres cosas que antes no
podía: saber \textbf{antes de empezar} si un libro concreto le queda grande, ver
\textbf{en qué capítulo} se va a atascar y decidir si lo salta o lo prepara, y llevarse
las palabras que consultó \textbf{con la frase en la que aparecieron}, que es la unidad
con la que se recuerda un idioma. Para quien enseña, el cambio es de otro tipo: un
catálogo de dominio público con dificultad medida convierte «lee algo de tu nivel» en una
lista concreta, y el coste marginal de añadir un libro es de segundos de cómputo, no de
una tarifa por título.

\textbf{Lo que este proyecto recomienda hacer a continuación, en orden de valor por hora
de trabajo}, está en el cuadro~\ref{cuad:accionables} con coste y criterio de terminado.
La remedición del catálogo, que era lo que separaba un prototipo de un producto, ya está
hecha, y el almacenamiento persistente también; lo siguiente es calibrar la escala, que es
la diferencia entre decirle a alguien que un libro es A2 y que lo sea.


% ===========================================================================
\section*{Declaración de alcance}
\addcontentsline{toc}{section}{Declaración de alcance}
% ===========================================================================

LiteraLevel es una herramienta de apoyo a la lectura. \textbf{No es un instrumento de
evaluación ni de certificación lingüística.} No debe usarse para asignar calificaciones,
para decidir la admisión de una persona a un programa, para acreditar un nivel del Marco
Común Europeo de Referencia ni para ninguna decisión con consecuencias para el evaluado.
La escala 0--2 es propia; las letras A1 a C2 son el ancla con la que se entrenó el modelo
y se muestran únicamente como orientación.

El sistema mide dificultad lingüística de prosa continua en inglés. No mide calidad
literaria, no mide comprensión del lector y no emite juicios sobre el contenido de las
obras. El catálogo no se filtra por contenido: se informa al lector mediante un aviso en
la ficha del libro y la decisión queda en sus manos.

Los datos personales del lector —los libros empezados, el vocabulario consultado y el
perfil de nivel— se guardan separados por navegador y no se comparten entre visitantes. Sí
se comparten las cachés de contenido ---definiciones, audio generado, curvas y resúmenes de
un mismo texto y el análisis de lectura natural de una frase---, que no son historial de
nadie. Ningún texto con derechos se distribuye ni se imprime fuera del entorno de cómputo:
de los corpus con licencia restrictiva solo salen estadísticas agregadas.

% ===========================================================================
\clearpage
\printbibliography[title={Referencias}]
\addcontentsline{toc}{section}{Referencias}

\end{document}
